% Options for packages loaded elsewhere
% Options for packages loaded elsewhere
\PassOptionsToPackage{unicode}{hyperref}
\PassOptionsToPackage{hyphens}{url}
%
\documentclass[
  korean,
  letterpaper,
  twoside,
  open=any,
  fontsize=12.5pt,
  headings=small]{scrbook}
\usepackage{xcolor}
\usepackage[paperwidth=148mm,paperheight=210mm,textwidth=87mm,inner=34mm,outer=27mm,top=24mm,textheight=148.5mm,bottom=37.5mm,headsep=7mm,footskip=14mm]{geometry}
\usepackage{amsmath,amssymb}
\setcounter{secnumdepth}{-\maxdimen} % remove section numbering
\usepackage{iftex}
\ifPDFTeX
  \usepackage[T1]{fontenc}
  \usepackage[utf8]{inputenc}
  \usepackage{textcomp} % provide euro and other symbols
\else % if luatex or xetex
  \usepackage{unicode-math} % this also loads fontspec
  \defaultfontfeatures{Scale=MatchLowercase}
  \defaultfontfeatures[\rmfamily]{Ligatures=TeX,Scale=1}
\fi
\usepackage{lmodern}
\ifPDFTeX\else
  % xetex/luatex font selection
\fi
% Use upquote if available, for straight quotes in verbatim environments
\IfFileExists{upquote.sty}{\usepackage{upquote}}{}
\IfFileExists{microtype.sty}{% use microtype if available
  \usepackage[]{microtype}
  \UseMicrotypeSet[protrusion]{basicmath} % disable protrusion for tt fonts
}{}
\makeatletter
\@ifundefined{KOMAClassName}{% if non-KOMA class
  \IfFileExists{parskip.sty}{%
    \usepackage{parskip}
  }{% else
    \setlength{\parindent}{0pt}
    \setlength{\parskip}{6pt plus 2pt minus 1pt}}
}{% if KOMA class
  \KOMAoptions{parskip=half}}
\makeatother
% Make \paragraph and \subparagraph free-standing
\makeatletter
\ifx\paragraph\undefined\else
  \let\oldparagraph\paragraph
  \renewcommand{\paragraph}{
    \@ifstar
      \xxxParagraphStar
      \xxxParagraphNoStar
  }
  \newcommand{\xxxParagraphStar}[1]{\oldparagraph*{#1}\mbox{}}
  \newcommand{\xxxParagraphNoStar}[1]{\oldparagraph{#1}\mbox{}}
\fi
\ifx\subparagraph\undefined\else
  \let\oldsubparagraph\subparagraph
  \renewcommand{\subparagraph}{
    \@ifstar
      \xxxSubParagraphStar
      \xxxSubParagraphNoStar
  }
  \newcommand{\xxxSubParagraphStar}[1]{\oldsubparagraph*{#1}\mbox{}}
  \newcommand{\xxxSubParagraphNoStar}[1]{\oldsubparagraph{#1}\mbox{}}
\fi
\makeatother


\usepackage{longtable,booktabs,array}
\newcounter{none} % for unnumbered tables
\usepackage{calc} % for calculating minipage widths
% Correct order of tables after \paragraph or \subparagraph
\usepackage{etoolbox}
\makeatletter
\patchcmd\longtable{\par}{\if@noskipsec\mbox{}\fi\par}{}{}
\makeatother
% Allow footnotes in longtable head/foot
\IfFileExists{footnotehyper.sty}{\usepackage{footnotehyper}}{\usepackage{footnote}}
\makesavenoteenv{longtable}
\usepackage{graphicx}
\makeatletter
\newsavebox\pandoc@box
\newcommand*\pandocbounded[1]{% scales image to fit in text height/width
  \sbox\pandoc@box{#1}%
  \Gscale@div\@tempa{\textheight}{\dimexpr\ht\pandoc@box+\dp\pandoc@box\relax}%
  \Gscale@div\@tempb{\linewidth}{\wd\pandoc@box}%
  \ifdim\@tempb\p@<\@tempa\p@\let\@tempa\@tempb\fi% select the smaller of both
  \ifdim\@tempa\p@<\p@\scalebox{\@tempa}{\usebox\pandoc@box}%
  \else\usebox{\pandoc@box}%
  \fi%
}
% Set default figure placement to htbp
\def\fps@figure{htbp}
\makeatother



\ifLuaTeX
\usepackage[bidi=basic,shorthands=off,provide=*]{babel}
\else
\usepackage[bidi=default,shorthands=off,provide=*]{babel}
\fi
\ifLuaTeX
  \usepackage{selnolig} % disable illegal ligatures
\fi


\setlength{\emergencystretch}{3em} % prevent overfull lines

\providecommand{\tightlist}{%
  \setlength{\itemsep}{0pt}\setlength{\parskip}{0pt}}



 


% ---- 한글 조판: kotex(xetexko) ----
\usepackage{kotex}
% 이야기 본문은 명조 계열로 통일한다. 메시지·앱 알림·메모는 아래 전용 매크로 글꼴을 유지한다.
% Bundle the OFL fonts used by the book so sandboxed XeTeX renders are stable.
\setmainfont[Path=tex/fonts/]{NanumMyeongjo.ttc}
\setsansfont[Path=tex/fonts/]{NanumGothic.ttc}
\setmonofont[Path=tex/fonts/]{D2Coding.ttc}
\setmainhangulfont[Path=tex/fonts/]{NanumMyeongjo.ttc}
\setsanshangulfont[Path=tex/fonts/]{NanumGothic.ttc}
\setmonohangulfont[Path=tex/fonts/]{D2Coding.ttc}
\newfontfamily\latinx{lmroman10-regular.otf}   % 나눔명조에 없는 라틴 확장 문자(é 등)
\newhangulfontfamily\titlehandwriting{BM Yeonsung} % 배달의민족 연성체: 표제지 전용
% 『시간 고양이』 1권 미리보기 실측 기준: 본문 폭 87mm, 쪽당 19줄, 줄 간격 약 8mm(22.7pt), 글자 약 12~12.5pt, 양끝 맞춤
\makeatletter\renewcommand\normalsize{\@setfontsize\normalsize{12.5pt}{22.7pt}}\makeatother\normalsize
\linespread{1}
\setlength{\parindent}{1em}
\setlength{\parskip}{0pt}
\setlength{\parfillskip}{0pt plus 0.8\linewidth}   % 문단 끝 줄이 너무 짧지 않게(최소 약 1/5)
\setlength{\emergencystretch}{1em}
% ---- 외톨이 줄·쪽 끝 제목 막기 ----
\clubpenalty=10000 \widowpenalty=10000 \displaywidowpenalty=10000
\brokenpenalty=10000
\raggedbottom
\usepackage{needspace}
\usepackage{etoolbox}
\pretocmd{\section}{\Needspace{5\baselineskip}}{}{}
\pretocmd{\subsection}{\Needspace{4\baselineskip}}{}{}
% ---- 머리말·쪽번호 ----
\usepackage[automark]{scrlayer-scrpage}
\automark[chapter]{chapter}
\clearpairofpagestyles
\ofoot*{\pagemark}
% 어린이 독서 흐름을 위해 쪽 위 반복 장 제목(머리글)은 표시하지 않는다.
\ohead{}
\renewcommand*{\chaptermarkformat}{}
\setkomafont{pageheadfoot}{\sffamily\small\color{gray}}
\setkomafont{pagenumber}{\sffamily\small\color{gray}}
\setkomafont{chapter}{\sffamily\bfseries}
\setkomafont{section}{\sffamily\bfseries\large}
\setkomafont{subsection}{\sffamily\bfseries}
\setkomafont{subsubsection}{\sffamily\bfseries}
% ---- 차례: 장 번호 고정 폭 ----
\newcommand{\tocnum}[1]{\makebox[3.3em][l]{#1}}
% ---- 로그 블록(ASCII 화면) ----
\usepackage{fancyvrb}
\usepackage{xcolor}
\usepackage{tikz}
\usetikzlibrary{arrows.meta,positioning,shapes.geometric}
\usepackage{eso-pic}
\definecolor{memoBg}{gray}{0.96}
\definecolor{memoLine}{gray}{0.48}
\definecolor{memoText}{gray}{0.16}
\newcommand{\memobox}[2]{%
  \par\medskip\noindent\fcolorbox{memoLine}{memoBg}{%
    \begin{minipage}{\dimexpr\linewidth-2\fboxsep-2\fboxrule}
    {\sffamily\bfseries\fontsize{11pt}{14pt}\selectfont\color{memoText}#1}\par
    \vspace{1.5mm}
    {\normalfont\fontsize{11.5pt}{18pt}\selectfont\color{black}#2}
    \end{minipage}}\par\vspace{4mm}}
\newcommand{\logfont}{\linespread{1}\fontsize{10.8}{15.2}\selectfont}
% ---- 메신저·앱 알림: 바깥 테두리 없이, 읽기 쉬운 회색 말풍선 ----
\newenvironment{chatthread}{%
  \par\Needspace{5\baselineskip}\medskip\noindent\begin{center}%
  \begin{minipage}{.90\linewidth}\vspace{1.5mm}%
}{%
  \vspace{1.5mm}\end{minipage}\end{center}\par\medskip%
}
\newcommand{\chatroomtitle}[1]{%
  {\noindent\hspace*{\dimexpr.20\linewidth-4.5mm\relax}%
   \sffamily\footnotesize\color{gray}#1\par}\vspace{1mm}%
}
\newcommand{\chatbubble}[2]{%
  \par\noindent\makebox[\linewidth][r]{%
  \begin{tikzpicture}
    \node[anchor=east,rounded corners=3mm,fill=gray!12,inner sep=2.5mm,
      text width=.80\linewidth,align=left] (bubble)
      {{\sffamily\footnotesize\color{gray}#1\par}\vspace{.5mm}%
       \sffamily\fontsize{11.2pt}{15.5pt}\selectfont\color{black}#2};
    \fill[gray!12] ([xshift=-8mm]bubble.south east) --
      ([xshift=-2mm]bubble.south east) --
      ([xshift=-1mm,yshift=-5mm]bubble.south east) -- cycle;
  \end{tikzpicture}}\par\vspace{.5mm}%
}
\newcommand{\appnotice}[2]{%
  \appnoticefrom{#1}{}{#2}%
}
\newcommand{\appnoticefrom}[3]{%
  \par\addvspace{6mm}\Needspace{5\baselineskip}\noindent\makebox[\linewidth][c]{%
  \begin{minipage}{.90\linewidth}
    {\noindent\sffamily\footnotesize\color{gray}#1\par}\vspace{1.5mm}%
    \noindent\colorbox{gray!12}{\parbox{.78\linewidth}{%
      \if\relax\detokenize{#2}\relax\else
        {\sffamily\footnotesize\color{gray}보낸 사람: #2\par}\vspace{1mm}%
      \fi
      \sffamily\fontsize{10.8pt}{15pt}\selectfont\color{black}#3}}%
    \vspace{2mm}%
  \end{minipage}}\par\addvspace{6mm}%
}
\newcommand{\stickynote}[1]{%
  \par\addvspace{7mm}%
  \begingroup\setlength{\fboxsep}{3mm}%
  \noindent\makebox[\linewidth][c]{%
  \fcolorbox{gray!45}{gray!5}{\begin{minipage}{.78\linewidth}
    {\sffamily\fontsize{10.5pt}{16.5pt}\selectfont\color{black}#1\par}
  \end{minipage}}}%
  \endgroup\par\addvspace{8mm}%
}
% ---- 장 첫 쪽 날짜 줄 ----
\newcommand{\dateline}[1]{\par\noindent{\sffamily\small\color{gray}#1}\par\medskip}
% ---- 표제지: 가로 제목 + 투명 컷아웃 그림 ----
\makeatletter
\renewcommand{\maketitle}{%
  \begin{titlepage}
  \thispagestyle{empty}
  \AddToShipoutPictureBG*{%
    \AtPageLowerLeft{\hspace*{10mm}\raisebox{-4mm}{\includegraphics[width=114.4mm]{images/title-adventure.png}}}%
  }%
  \vspace*{17mm}
  \begin{center}
  \makebox[0pt][c]{{\titlehandwriting\fontsize{30}{36}\selectfont 곰새코지 돌고래 목소리 사건}}\par
  \vspace{4mm}
  {{\titlehandwriting\fontsize{9.5}{13}\selectfont 이희근\par}}
  \end{center}
  \end{titlepage}
}
\makeatother
% ---- 장 도비라 ----
\usepackage{graphicx}
\newcommand{\bodystyle}{\setlength{\parindent}{1em}\setlength{\parskip}{0pt}\normalsize}
\newcommand{\fullpageimage}[1]{%
  \clearpage\thispagestyle{empty}%
  \AddToShipoutPictureBG*{%
    \AtPageLowerLeft{\includegraphics[width=\paperwidth,height=\paperheight]{#1}}%
  }%
  \null\clearpage}
\newcommand{\chapterartpage}[2]{%
  \IfFileExists{#1}{\fullpageimage{#1}}%
    {\IfFileExists{#2}{\fullpageimage{#2}}%
      {\clearpage\thispagestyle{empty}\vspace*{\fill}\begin{center}\fbox{\parbox[c][0.8\textheight][c]{0.8\textwidth}{\centering\color{gray}\small 장 전면 그림이 들어갈 자리}}\end{center}\vspace*{\fill}\clearpage}}}
\newcommand{\dobira}[4]{%
  \clearpage\thispagestyle{plain.scrheadings}\bodystyle%
  \addcontentsline{toc}{chapter}{\tocnum{#1}#2}%
  \markboth{#2}{#2}%
  \vspace*{42mm}%
  {\centering\sffamily{\small #1}\par\vspace{2mm}{\fontsize{19}{27}\selectfont\bfseries #2}\par}%
  \vspace{3mm}%
  \chapterartpage{#4}{#3}}
% ---- 부록(뒷부분): 들여쓰기 없이 문단 사이 띄우기 ----
\newcommand{\backmatterstyle}{\setlength{\parindent}{0pt}\setlength{\parskip}{0.45em plus 0.1em}\backmattersize}
% ---- 판권: 맨 마지막 짝수 쪽 아래 ----
\newcommand{\colophonstart}{\clearpage\thispagestyle{empty}\vspace*{\fill}\begingroup\small\setlength{\parindent}{0pt}\setlength{\parskip}{0.5em}\linespread{1.3}\selectfont}
\newcommand{\colophonend}{\par\endgroup}
% ---- 표: A5 폭 안에 들어오게 ----
\setlength{\tabcolsep}{3pt}
\AtBeginEnvironment{longtable}{\small\linespread{1.2}\selectfont}

% ---- 권말: 본문보다 작은 글자(11pt/18pt) ----
\newcommand{\backmattersize}{\fontsize{10.5pt}{16.5pt}\selectfont}

% ---- 차례 한 쪽에 ----
\RedeclareSectionCommand[tocbeforeskip=1pt plus 1pt]{chapter}
\setkomafont{chapterentry}{\normalfont\sffamily\small}
\setkomafont{chapterentrypagenumber}{\normalfont\sffamily\small}
\AfterTOCHead{\small\linespread{1}\selectfont}
% ---- 앞부분 빈 쪽 줄이기(차례 다음 바로 시작) ----
\AtBeginDocument{\let\cleardoublepage\clearpage}
\RedeclareSectionCommand[beforeskip=6mm,afterskip=5mm]{chapter}
\makeatletter\renewcommand\mainmatter{\clearpage\@mainmattertrue\pagenumbering{arabic}}\makeatother
% ======== 읽기 쉬운 앞부분: 인물 쪽 ========
\definecolor{peopleSea}{HTML}{EAF7F7}
\definecolor{peopleBlue}{HTML}{D9ECFF}
\definecolor{peopleLine}{HTML}{9CB8C9}
\definecolor{peopleText}{HTML}{24404F}
\definecolor{peopleGray}{HTML}{6F7D86}
\newenvironment{peoplepage}{%
  \par\setlength{\parindent}{0pt}\setlength{\parskip}{0pt}%
  \sffamily\fontsize{9.1pt}{12.1pt}\selectfont%
}{\par}
\newcommand{\peopleintro}{%
  \noindent\fcolorbox{peopleLine}{peopleSea}{%
    \begin{minipage}{\dimexpr\linewidth-2\fboxsep-2\fboxrule}
    \sffamily\bfseries\color{peopleText}누가 누구인지 먼저 훑어보기\par
    \vspace{0.2mm}
    {\normalfont\footnotesize\color{peopleGray}아이들 · 가족 · 연구소 · 마을 어른들 · 돌고래가 얽히며 사건이 움직입니다.}
    \end{minipage}}\par\vspace{1.8mm}}
\newcommand{\persongroup}[1]{%
  \par\addvspace{2.1mm}%
  \noindent\colorbox{peopleBlue}{%
    \makebox[\dimexpr\linewidth-2\fboxsep][l]{\sffamily\bfseries\footnotesize\color{peopleText}\quad #1}}%
  \par\vspace{0.8mm}}
\newcommand{\person}[4]{%
  \par\noindent\fcolorbox{peopleLine}{white}{%
    \begin{minipage}{\dimexpr\linewidth-2\fboxsep-2\fboxrule}
    {\sffamily\bfseries\fontsize{10.7pt}{12.5pt}\selectfont\color{peopleText}#1}%
    \if\relax\detokenize{#2}\relax\else\hspace{1.5mm}{\sffamily\scriptsize\color{peopleGray}#2}\fi
    \if\relax\detokenize{#4}\relax\else\hfill{\sffamily\scriptsize\color{peopleGray}소품: #4}\fi\par
    \vspace{0.2mm}
    {\normalfont\fontsize{9.1pt}{11.8pt}\selectfont #3}%
    \end{minipage}}%
  \par\vspace{0.7mm}}
\newenvironment{profilepage}{%
  \par\setlength{\parindent}{0pt}\setlength{\parskip}{0pt}%
  \sffamily\fontsize{8pt}{10pt}\selectfont%
}{\par}
\newcommand{\profileintro}{%
  {\sffamily\footnotesize\color{peopleGray}사진을 먼저 보고, 이름과 역할만 가볍게 훑어보세요.}\par
  \vspace{4mm}}
\newcommand{\profile}[4]{%
  \begin{minipage}[t]{0.305\linewidth}
  \centering
  \fcolorbox{peopleLine}{white}{\includegraphics[width=20mm,height=20mm]{#1}}\par
  \vspace{0.6mm}
  {\sffamily\bfseries\fontsize{10pt}{11.5pt}\selectfont\color{peopleText}#2}\par
  \if\relax\detokenize{#3}\relax\else{\sffamily\scriptsize\color{peopleGray}#3\par}\fi
  \vspace{0.3mm}
  {\normalfont\fontsize{7pt}{8.4pt}\selectfont #4\par}
  \end{minipage}}
\newcommand{\profilecompact}[4]{%
  \begin{minipage}[t]{0.31\linewidth}
  \centering
  \fcolorbox{peopleLine}{white}{\includegraphics[width=25mm,height=25mm]{#1}}\par
  \vspace{0.6mm}
  {\sffamily\bfseries\fontsize{10pt}{11.8pt}\selectfont\color{peopleText}#2}\par
  \if\relax\detokenize{#3}\relax\else{\sffamily\scriptsize\color{peopleGray}#3\par}\fi
  \vspace{0.4mm}
  {\normalfont\fontsize{7.8pt}{9.5pt}\selectfont #4\par}
  \end{minipage}}
\newcommand{\profilehero}[4]{%
  \begin{minipage}[t]{0.31\linewidth}
  \centering
  \fcolorbox{peopleLine}{white}{\includegraphics[width=25mm,height=25mm]{#1}}\par
  \vspace{0.8mm}
  {\sffamily\bfseries\fontsize{10.8pt}{12.5pt}\selectfont\color{peopleText}#2}\par
  \if\relax\detokenize{#3}\relax\else{\sffamily\fontsize{8.8pt}{10.5pt}\selectfont\color{peopleGray}#3\par}\fi
  \vspace{0.7mm}
  {\normalfont\fontsize{7.8pt}{9.8pt}\selectfont #4\par}
  \end{minipage}}
% ======== 이야기 지도형 차례 ========
\usepackage{tocbasic}
\DeclareTOCStyleEntry[
  beforeskip=4pt plus 1.5pt,
  entryformat=\normalfont\sffamily\fontsize{11pt}{19pt}\selectfont,
  pagenumberformat=\normalfont\sffamily\fontsize{10.5pt}{19pt}\selectfont\color{gray},
  linefill=\hfill,
  indent=7mm, numwidth=0pt
]{tocline}{chapter}
\definecolor{tocsea}{HTML}{16879A}
\renewcommand{\tocnum}[1]{\makebox[3.8em][l]{{\color{gray}\footnotesize #1}}}
\newcommand{\tocgroup}[1]{%
  \addvspace{6mm}%
  \noindent\textcolor{tocsea}{\rule{1.2mm}{4.2mm}}\hspace{2.5mm}%
  {\sffamily\bfseries\fontsize{12}{16}\selectfont\color{tocsea}#1}\par
  \addvspace{1.5mm}}
\AfterTOCHead[toc]{\enlargethispage{10mm}}
\makeatletter
\@ifpackageloaded{bookmark}{}{\usepackage{bookmark}}
\makeatother
\makeatletter
\@ifpackageloaded{caption}{}{\usepackage{caption}}
\AtBeginDocument{%
\ifdefined\contentsname
  \renewcommand*\contentsname{목차}
\else
  \newcommand\contentsname{목차}
\fi
\ifdefined\listfigurename
  \renewcommand*\listfigurename{그림 목록}
\else
  \newcommand\listfigurename{그림 목록}
\fi
\ifdefined\listtablename
  \renewcommand*\listtablename{표 목록}
\else
  \newcommand\listtablename{표 목록}
\fi
\ifdefined\figurename
  \renewcommand*\figurename{그림}
\else
  \newcommand\figurename{그림}
\fi
\ifdefined\tablename
  \renewcommand*\tablename{표}
\else
  \newcommand\tablename{표}
\fi
}
\@ifpackageloaded{float}{}{\usepackage{float}}
\floatstyle{ruled}
\@ifundefined{c@chapter}{\newfloat{codelisting}{h}{lop}}{\newfloat{codelisting}{h}{lop}[chapter]}
\floatname{codelisting}{목록}
\newcommand*\listoflistings{\listof{codelisting}{코드 목록}}
\makeatother
\makeatletter
\makeatother
\makeatletter
\@ifpackageloaded{caption}{}{\usepackage{caption}}
\@ifpackageloaded{subcaption}{}{\usepackage{subcaption}}
\makeatother
\usepackage{bookmark}
\IfFileExists{xurl.sty}{\usepackage{xurl}}{} % add URL line breaks if available
\urlstyle{same}
% fallback for those not using the hyperref driver hyperxmp:
\makeatletter
\@ifundefined{xmpquote}{\newcommand{\xmpquote}[1]{#1}}{}
\makeatother
\hypersetup{
  pdftitle={곰새코지 돌고래 목소리 사건},
  pdfauthor={이희근},
  pdflang={ko},
  hidelinks,
  pdfcreator={LaTeX via pandoc}}


\title{곰새코지 돌고래 목소리
사건}
\author{이희근}
\date{}
\begin{document}
\frontmatter
\maketitle

\renewcommand*\contentsname{차례}
{
\setcounter{tocdepth}{0}
\tableofcontents
}

\mainmatter
\addtocontents{toc}{\protect\tocgroup{읽기 전에}}

\bookmarksetup{startatroot}

\chapter*{곰새코지
사람들}\label{uxacf0uxc0c8uxcf54uxc9c0-uxc0acuxb78cuxb4e4}
\addcontentsline{toc}{chapter}{곰새코지
사람들}

\markboth{곰새코지
사람들}{곰새코지 사람들}

\backmatterstyle

\begin{profilepage}
\profileintro
\noindent\makebox[\linewidth][c]{\profilehero{images/portraits/danu.jpg}{고단우}{열두 살 · 6학년}{시민 라벨러 리더보드 1위. 코딩을 좋아한다.}\hspace{2mm}\profilehero{images/portraits/chaea.jpg}{양채아}{열두 살 · 6학년}{‘채아의 돌핀로그’를 하는 단우의 친구.}\hspace{2mm}\profilehero{images/portraits/haan.jpg}{다원}{열한 살 · 5학년}{틀린 데를 잘 찾아내는 미숙의 딸.}}\par\vspace{8mm}
\noindent\makebox[\linewidth][c]{\profilehero{images/portraits/grandma.jpg}{할머니}{해녀}{마지막 물질을 앞둔 해녀. 돌고래를 잘 알아본다.}\hspace{2mm}\profilehero{images/portraits/mom.jpg}{엄마}{간호사}{밤 근무가 많아 단우와 자꾸 엇갈린다.}\hspace{2mm}\profilehero{images/portraits/tak.jpg}{탁 박사}{과학자}{물마루 AI 연구소에서 돌고래 소리를 연구한다.}}
\clearpage
\thispagestyle{plain.scrheadings}
% 두 쪽 펼침에서 첫 인물 페이지와 같은 물리적 조판면을 쓴다.
% 짝수 쪽의 바깥 여백 차이(7mm)를 보정하고, 첫 사진 행 높이도 맞춘다.
\vspace*{-5.5mm}
\begin{center}{\sffamily\color{gray}곰새코지 사람들}\end{center}
\begin{center}{\sffamily\fontsize{8.5pt}{11pt}\selectfont\color{peopleGray}제주 바닷가에 있는 작은 마을, 곰새코지의 사람들입니다.}\end{center}
\vspace{3.5mm}
\noindent\hbox{\kern7mm\makebox[\linewidth][c]{\profilecompact{images/portraits/bong.jpg}{봉 소장}{연구소장}{AI 바당이의 첫 문장을 발표한 사람.}\hspace{2mm}\profilecompact{images/portraits/osamchon.jpg}{오삼촌}{배 조종}{연구소 배를 모는 아저씨. 성이 오씨다.}\hspace{2mm}\profilecompact{images/portraits/tu.jpg}{미숙}{해녀}{다원의 엄마. 곰새코지의 젊은 해녀.}}}\par
\noindent\hbox{\kern7mm\makebox[\linewidth][c]{\profilecompact{images/portraits/chaea-dad.jpg}{채아 아빠}{관광선 선장}{관광선 ‘곰새코지 3호’를 몬다.}\hspace{2mm}\profilecompact{images/portraits/bokja.jpg}{김복자}{해녀회장}{배가 작업장 가까이 오는 걸 싫어한다.}\hspace{2mm}\profilecompact{images/portraits/deoksu.jpg}{강덕수}{어촌계장}{마을 어민 모임의 대표. 항구 확장을 원한다.}}}\par
\noindent\hbox{\kern7mm\makebox[\linewidth][c]{\profilecompact{images/portraits/narae.jpg}{윤나래}{돌고래쉼표}{돌고래 보호 단체 활동가. 개인 방송을 한다.}\hspace{7mm}\profilecompact{images/portraits/bulssugi.jpg}{불쑥이}{돌고래}{카메라 앞에 불쑥 떠오른다.}}}
\end{profilepage}

\fullpageimage{images/pagefill/map.png}

\addtocontents{toc}{\protect\tocgroup{이야기}}\dobira{1장}{새벽 세 시, 내 아이디가 찍혔다}{images/pagefill/ch01.png}{images/pagefill/ch01-page.png}

\section{새벽 세 시의
라벨}\label{uxc0c8uxbcbd-uxc138-uxc2dcuxc758-uxb77cuxbca8}

\dateline{6월 24일 금요일 · 방송 특집까지 D-22}

화면 아래쪽에 자막이 떴다. 흰
글자 세 개와 마침표 하나.

\textbf{돌려줘.}

거실이 조용해졌다. 생중계 속
연구소장은 돌고래 넥타이를
만지며 숨을 골랐고, 할머니는
귤을 까던 손을 멈췄다. 단우는
자막 아래, 깨알 같은 글씨를
읽었다. 라벨러 danu\_0101.

단우의 아이디였다. 그런데
단우는 그 라벨을 붙인 기억이
없었다.

처음 1초 동안은 손끝이
찌릿했다. 전국 생중계에서 AI가
처음 내놓은 문장 밑에, 열두 살
고단우의 아이디가 박혀 있었다.

돌고래 소리 조각을 듣고
이름표를 붙이는 시민 라벨러는
전국에 수두룩했다. 그 많은
이름표 가운데 단우의 것이
뽑혔다. 한 달째 지켜 온 라벨러
순위표 맨 꼭대기, 리더보드
1위라는 말도 진짜처럼 들렸다.

두 번째 1초에는 목덜미가
서늘해졌다. 소파 쿠션을
슬그머니 무릎 위로 끌어당겼다.

화면 귀퉁이에 「바당이, 첫
문장을 말하다」가 떠 있었다.
바당이. 물마루 AI 연구소가
만든 돌고래 소리 분석 AI.
라벨러들이 붙인 이름표를
재료로 돌아가는 AI였다.

할머니는 그 화면을 노려보고
있었다.

``돌고래가 무사 사람 말을
해?''

``번역한 거예요. 바당이가
돌고래 소리를 사람 말로 바꿔
준 거라고요.''

``번역은 무신 번역.''

할머니는 귤을 입에 넣지도,
내려놓지도 않았다.

그때 덩치 큰 돌고래 하나가
카메라 바로 앞에 불쑥
떠올랐다. 렌즈에 물이 튀어
소장 얼굴이 뿌옇게 번졌다.

``저건 불쑥이여.'' 할머니가 귤
든 손으로 화면을 가리켰다.
``지느러미 끝 뜯긴 거 봐라.
어디든 불쑥 끼어들주.''

단우는 이름이 딱 맞다고
생각했다. 카메라 앞이면 꼭
불쑥 나타나는 녀석이었다.

불쑥이 옆으로 MR-012라는 작은
글씨가 떴다. 연구소는
등지느러미 흠집을 지문처럼
찍어 돌고래마다 번호를 붙인다.
할머니는 그 일을 눈으로 하고
있었다.

``이 첫 문장이 무엇을
뜻하는지, 이제 우리 모두 함께
생각해 볼 때입니다.''

엄마가 집에 있었으면
``밥은?''부터 물었을 것이다.
엄마는 오늘도 제주시 병원에서
야간 근무였다.

``저 방에 갈게요.''

할머니는 대답 대신 귤껍질을
반으로 접었다. 복도에서
돌아보니, 불쑥이가 이번엔 소장
등 뒤에서 또 불쑥 떠오르고
있었다.

``또 끼어드네.''

\Needspace{4\baselineskip}\begin{center}◇\end{center}\nopagebreak

단우는 방에서 라벨러 앱 `귀
기울임'에 자기 아이디를
검색했다. 맨 위에 첫 문장의
근거가 된 녹음 조각이 떠
있었다.

\begin{Verbatim}[samepage=true,frame=single,framesep=2.5mm,rulecolor=\color{gray},formatcom=\logfont]
#A-3318
03:02:08 | 라벨러 | danu_0101
03:02:08 | 라벨   | 어미-새끼 분리 위협
03:02:08 | 관광선 | 있음(80m)
\end{Verbatim}

새벽 3시 2분 8초. 어미-새끼
분리 위협. 엄마 돌고래와
새끼가 떨어지게 생겼다는
뜻이었다.

단우가 고른 말이 아니었다. 그
시간에 단우는 자고 있었다. 밤
9시면 눈꺼풀이 저절로 내려오는
체질이었다. 목록을 내리자
02:51, 03:02, 03:17, 04:05.
전부 이불 속에 있던
시간이었다.

그때 노트북 구석의 작은 창이
눈에 들어왔다. 자막을 본
순간부터 떠올리지 않으려
애쓰던 창이었다.

\begin{Verbatim}[samepage=true,frame=single,framesep=2.5mm,rulecolor=\color{gray},formatcom=\logfont]
labelbot.py | 상태 | 실행 중
labelbot.py | 연동 | 관광선 위치 앱
\end{Verbatim}

새벽 3시 2분 8초에 라벨을 붙인
건 단우가 아니었다. 단우 대신
라벨을 붙이게 만들어 둔 작은
프로그램, labelbot.py였다.

한 달 전, 유튜브 코딩 강의를
두 배속으로 보며 혼자 짠
프로그램이었다. 자는 동안
라벨을 대신 붙여 주는
프로그램. `귀 기울임'은 사람이
직접 듣고 붙여야 하는
앱이었다. 가입하던 날에는
`모든 라벨은 제가 듣고
붙이겠습니다' 칸에 체크도
했다. 그러고도 아무에게도
말하지 않았다. 리더보드 1위
옆에 금색 왕관이 붙은 건
그다음 주부터였다.

손가락이 책상을 두드리기
시작했다.
\mbox{\texttt{0, 1, 0, 1}}.

메모 앱을 열자 라벨러를 시작한
1월의 메모가 맨 위에 있었다.

\memobox{단우의 메모 \#1}{라벨 = 소리 조각에 붙이는 이름표.\\사진에 다는 해시태그 같은 것.}

연구소에 말해야 하나.
할머니한테? 아니, 끄면 아무도
모르지 않을까.

단우는 메모 앱을 닫고, 작은
창의 빨간 정지 버튼을 눌렀다.

\texttt{실행\ 중}이
\texttt{중지됨}으로 바뀌었다.

그걸로 다 끝난 것처럼, 단우는
노트북을 덮었다.

\Needspace{4\baselineskip}\begin{center}◇\end{center}\nopagebreak

10시가 넘자 책상 위 휴대폰이
빙글빙글 돌기 시작했다.

\begin{chatthread}
\chatroomtitle{곰새코지 6학년 수다방}
\chatbubble{채아 (6월 24일 밤 10:10)}{야, 고단우! 너 뉴스 나왔어!\\바당이 첫 문장에 네 라벨이 쓰였대!}
\chatbubble{채아 (6월 24일 밤 10:11)}{연구소 페이지에 찍힌 시간까지 떴어!\\근데 새벽 3시에 깨 있었냐?\\너 9시면 자잖아}
\end{chatthread}

일곱 명뿐인 6학년 단톡방이
밀려 올라갔다. 앱도 단우의
라벨이 바당이 첫 문장의 근거가
됐다고 축하했다.

단우는 단톡방에
\texttt{그거\ 데이터\ 오류일\ 수도}를
썼다가 지웠다.
\texttt{어제\ 좀\ 늦게\ 잤어}까지
쳤다가 그것도 지우고, 결국
휴대폰을 뒤집어 놓았다.

뒤집힌 휴대폰이 또 울렸다.
리더보드의 단우 이름 바로
아래에 댓글이 달려 있었다.

\begin{Verbatim}[samepage=true,frame=single,framesep=2.5mm,rulecolor=\color{gray},formatcom=\logfont]
새벽세시감귤: 축하드려요 1위님 ^^
\end{Verbatim}

리더보드 2위였다. 그
\texttt{\^{}\^{}}가 웃는
얼굴인지 비웃는 얼굴인지 알 수
없었다. 닉네임부터가 새벽 세
시였다. 이 사람은 대체 언제
자는 걸까.

방문 문턱에 할머니가 걸터앉아
주황색 테왁 끈을 고치고
있었다. 단우가 어릴 때
비치볼인 줄 알고 끌어안았던
둥근 부표였다. 해녀인 할머니는
물질 나가는 날이면 그걸 바다에
띄워 놓았다. 바닷속에서 성게를
따다가, 숨이 차면 올라와 그
위에 몸을 기댔다.

``낼 물때 좋다.'' 할머니가
끈을 이로 툭 끊었다. ``올해가
마지막이라, 하나라도 더
나가야주.''

단우 머릿속에 욕실 타일 틈의
거뭇한 때가 먼저 떠올랐다.
하지만 할머니의 물때는
바닷물이 빠지고 드는 때였다.
할머니는 달력을 안 보고도 그걸
알았다.

할머니가 고개도 들지 않고
말했다.

``테레비에 느 이름 나왔더라.''

``이름 아니고 아이디예요.''

``그게 그거주.''

할머니는 매듭을 당겨 보더니,
혼잣말처럼 덧붙였다.

``말젯이가 요새 안 뵌다.''

말젯이? 처음 듣는 이름이었다.
할머니가 `삼춘'이라 부르는
동네 해녀인지, 바다에 사는
무엇인지 알 수 없었다. 단우는
묻지 않고 이어폰을 꼈다.
음악은 틀지 않았다.

\Needspace{4\baselineskip}\begin{center}◇\end{center}\nopagebreak

11시 40분. 단톡방 숫자가
`99+'에서 멈추자, 단우는 불을
끄고 누웠다.

9시면 잔다던 애는 잠이 오지
않았다. 노트북을 다시 열어
봐도 작은 창에는
\texttt{중지됨}뿐이었다. 눈을
감으면 흰 글자 세 개가
떠올랐다. 마침표까지 꼬박꼬박.

천장에 파란 불빛이 번졌다.
베개 옆 휴대폰의 잠금 화면에
앱 알림이 떠 있었다.

\appnoticefrom{귀 기울임 · 연구소 공식 메시지}{탁은재 박사(물마루 AI 연구소)}{단우 학생, 라벨 \#A-3318 말이에요. 새벽 3시 2분 8초에 붙었던데—}

\chapterartpage{images/pagefill/ch02-page.png}{images/pagefill/ch02.png}

\section{다섯 개의
돌려줘}\label{uxb2e4uxc12f-uxac1cuxc758-uxb3ccuxb824uxc918}

\dateline{6월 25일 토요일 · 방송 특집까지 D-21}

단우는 그 알림을 밤새 열지
못했다.

아침 9시, 결국 엄지로 눌렀다.

\begin{quote}
단우 학생, 라벨 \#A-3318
말이에요. 새벽 3시 2분 8초에
붙었던데, 어떤 방식으로
붙였는지 설명해 줄 수 있어요?
월요일에 연구소로 와요.
보호자분께도 같은 내용을
보냈어요.
\end{quote}

`어떤 방식으로'에서 눈이
멈췄다. 마지막 줄에선 엄지도
멈췄다.

보호자. 엄마였다.

안방 문은 닫혀 있었다. 식탁 위
엄마 휴대폰에 손을 뻗었다가
거뒀다.

연구소는 월요일이었다. 사흘
남았다.

\Needspace{4\baselineskip}\begin{center}◇\end{center}\nopagebreak

10시쯤 단톡방이 울렸다.

\begin{chatthread}
\chatroomtitle{곰새코지 6학년 수다방}
\chatbubble{채아 (6월 25일 오전 10:02)}{포구에 방송차 두 대 왔어!}
\chatbubble{채아 (6월 25일 오전 10:03)}{소장님 넥타이 뭐야, 진짜야? ㅋㅋ}
\end{chatthread}

컵라면을 사러 가자 구멍가게
아저씨가 고개를 뺐다.

``단우 아니냐? 테레비
나왔다며.''

``아이디만요.''

어깨가 펴졌다. 아저씨는 봉지에
컵라면 하나를 더 넣어 주었다.

``유명인 서비스다.''

카메라들 한가운데에 그
연구소장, 봉 소장이 있었다.
넥타이엔 점박이 돌고래가
빼곡했다. 어젯밤엔 파란 돌고래
한 마리였는데.

채아가 옆에 와 있었다.

``넥타이 몇 개인지 내기할래?''

``확률적으로는 최소 두 개.''

``그건 나도 보여. 난 세 개에
건다. 지는 사람이 컵라면.''

단우는 봉지를 등 뒤로 돌렸다.
``이미 샀는데.''

``돌고래가 우리에게 처음 건넨
말입니다.''

기자 하나가 마이크를 내밀었다.
``확신도는 몇
퍼센트였습니까?''

``숫자보다 마음이 먼저
닿았잖아요.''

기자들은 받아 적었다. 그래서
몇 퍼센트인데요.

카메라들이 탈의장 쪽으로 몸을
틀었다. 해녀회장님은
마이크에도 멈추지 않았다.

``배가 작업장 코앞까지
들어오잖아. 바다를 돌려달라는
말이여.''

목소리 큰 어촌계장 아저씨가
끼어들었다.

``관광객 끊기민 마을 끝이여.
항구 넓히는 공사 들어오민 다시
온다.''

회장은 계장을 한 번 보고
들어갔다.

방파제 끝에선 보호 단체
`돌고래쉼표'가 `돌고래와 50m,
지켜 주세요' 현수막을 펼쳤다.
조금 떨어져 활동가 한 사람이
휴대폰을 들고 있었다. 채아가
그 방송을 보여 줬다.

\begin{Verbatim}[samepage=true,frame=single,framesep=2.5mm,rulecolor=\color{gray},formatcom=\logfont]
윤나래 개인 방송 | LIVE
시청자 | 3,204명
\end{Verbatim}

``자유를 돌려달라는 거예요.
새끼가 안 보인다면서요? 누가
데려간 건 아닌지도 봐야죠.''

새끼가 안 보인다. 어디서 들은
말 같았다.

`돌려줘'를 검색하자 구독자
80만 과학 채널 오션톡과 채아의
돌핀로그 영상이 나란히 떴다.

\begin{Verbatim}[samepage=true,frame=single,framesep=2.5mm,rulecolor=\color{gray},formatcom=\logfont]
검색 결과 | 새 영상
오션톡
인류에게 보내는 바다의 경고?
조회수 41만
채아의 돌핀로그
우리 배한테 화났대요?
조회수 1,002,418
\end{Verbatim}

``백만이야, 백만!''

``구독자 3만인데 백만이
봤다고. 이게 계산이 돼?''

``서른세 배쯤.''

``그런 계산 말고!''

채아가 휴대폰을 들었다.

``안녕, 돌핀로그 채아예요---
지금 곰새코지 포구는요---''

화면 가득 계장 아저씨 뒤통수가
들어왔다. 채아는 녹화를 껐다.

회장님은 바다, 계장님은
관광객, 그 활동가는 자유.
오션톡은 경고, 채아는 화난
돌고래. 같은 세 글자인데
얼굴이 다섯이었다.

``웃기지? 다들 자기 하고 싶은
말만 해.''

``확증편향이야. 보고 싶은 것만
골라 보는 거. 응원하는 팀이
지면 심판이 다 상대 팀 편드는
것 같잖아.''

``또 유튜브 두 배속으로
배웠지?''

``세 배속.''

``근데 넌? 네 라벨이잖아. 무슨
뜻인데?''

``그건 바당이가 번역한 거고.''

``그러니까 넌 무슨 뜻으로
붙였냐고.''

단우는 대답 대신 화면을 껐다.
꺼진 화면에 자기 얼굴이
비쳤다.

``됐다, 확증편향 박사님.''

채아는 코웃음 치고 댓글을
내렸다.

\begin{Verbatim}[samepage=true,frame=single,framesep=2.5mm,rulecolor=\color{gray},formatcom=\logfont]
바다덕후: 돌고래도 할 말은 하네 ㅋㅋ
제주토박이: 선장 아빠 처벌해라
\end{Verbatim}

웃음소리가 그쳤다.

계단 아래에 채아 아빠, 양
선장이 앉아 있었다. 휴대폰에
같은 글자가 쌓였다. 예약 취소.
예약 취소. 예약 취소.

선장이 고개를 들어 채아를
보았다.

``왔냐.''

포구를 떠날 때까지 채아는 웃지
않았다.

\Needspace{4\baselineskip}\begin{center}◇\end{center}\nopagebreak

토요일 저녁, 엄마가 깨는
소리에 단우는 방문을 닫고
이어폰을 꼈다.

발소리가 방문 앞에서 멈췄다.

``단우야, 자?''

단우는 눈을 감고 숨을 골랐다.
현관문이 닫힌 뒤 식탁에 쪽지가
있었다.

\stickynote{월요일 연구소.\\[1mm]연구소 오삼촌이 학교 앞으로 데리러 간대.\\[1mm]엄마는 근무.\\[1mm]끝나면 전화해.}

어떻게 붙였느냐는 말은 한 줄도
없었다. 단우는 쪽지를 반으로,
또 반으로 접었다.

일요일 저녁, 할머니가 딴
성게로 끓인 미역국이 올라왔다.
단우는 한 숟가락 뜨고 코를
찡그렸다.

``이 국 버그 있어요.''

``버그는 느 머릿속에 있주.''

할머니는 단우 그릇에 한 국자
더 떴다.

뉴스에 채아 영상이 나왔다.
관광선 옆에서 솟구친 돌고래가
눈 깜짝할 새 뱃머리 앞에서 또
솟구쳤다.

할머니가 숟가락을 멈췄다.

``저 아이 영상은 꼭 중간을
먹어.''

``영상이 뭘 먹어요?''

뉴스는 금요일 생중계도 틀었다.

\textbf{돌려줘.}

\dateline{근거 녹음 조각 \#A-3318 | 라벨러 danu\_0101}

그날은 아이디밖에 보이지
않았다. 오늘은 구석 작은
글씨까지 보였다.

\dateline{확신도 31\%}

시험 점수라면 엄마한테 못 보여
줄 숫자였다. 단우는 검색창에
`확신도'까지 치고는 지웠다.

\Needspace{4\baselineskip}\begin{center}◇\end{center}\nopagebreak

월요일 3교시, 선생님이 공문을
들고 들어왔다.

``푸른코드 청소년 AI 대회.
마감은 7월 16일 토요일 밤 11시
59분.''

선생님이 칠판 귀퉁이에 날짜를
적었다.

``1등 상은 서울 AI 캠프
참가권이랑 장학금.''

단우는 등받이에서 몸을 뗐다.
서울. 캠프. 화면 속에만 있는
도시였다.

``혼자도 되고 둘이도 돼. 대신
자기 손으로 만든 거여야 한다.
하고 싶은 사람은 칠판 옆
종이에 적어.''

쉬는 시간, 단우가 제일 먼저
일어났다. `내'를 쓸 때 펜 끝이
멈췄다.

\dateline{고단우 | 내 라벨이 만든 첫 문장}

교실에 ``오오'' 소리가 퍼졌다.
단우가 채아에게 펜을 흔들었다.

``너도 써. 백만 유튜버.''

``됐어. 난 AI 안 만들어.''

앞문을 지나던 5학년 다원이
멈췄다. 남의 글을 보면 틀린
데부터 찾는 애였다.

```만든'이 아니라 `만들어진'이
맞지 않아?''

``내 라벨이니까 `만든'이
맞아.''

``그럼 `내가 만든'이라고
쓰든가.''

단우는 못 들은 척 제목에
밑줄을 그었다. 다원은 공중에
`만들어진'을 써 보이고
사라졌다.

채아가 턱을 괴고 보았다.

``너 돌고래 실제로 본 적은
있냐?''

``데이터로는 1만 2천 번.''

``그거 다 네 귀로 들었어?''

책상 밑에서 단우 손가락이
움직였다.
\mbox{\texttt{0, 1, 0, 1}}.

점심시간, 교실 휴대폰들이
한꺼번에 울렸다.

\begin{Verbatim}[samepage=true,frame=single,framesep=2.5mm,rulecolor=\color{gray},formatcom=\logfont]
[속보] 물마루 AI 연구소 홍보실 발표
7월 16일(토) 정오 특집 생중계
'바당이, 바다에서 말하다'
시민 과학 소년 고단우 군 출연 예정
\end{Verbatim}

여섯 명의 눈이 단우 쪽으로
왔다.

``와, 고단우 생방송!''

고단우 군. 아이디가 아니라
이름이었다.

출연한다고 한 적은 없었다.
`어떤 방식으로'를 물어 온 그
연구소였다. 오늘 방과 후에
가기로 한.

단우 눈이 칠판 귀퉁이 날짜로
갔다. 7월 16일 토요일. 정오에
생중계, 밤 11시 59분에 마감.
같은 날이었다.

``너 나와?''

``\ldots 나도 지금 알았어.''

\dobira{2장}{잘린 장면은 누구 편일까?}{images/pagefill/ch03.png}{images/pagefill/ch03-page.png}

\section{AI가 말을 만든
곳}\label{aiuxac00-uxb9d0uxc744-uxb9ccuxb4e0-uxacf3}

\dateline{6월 27일 월요일 · 방송 특집까지 D-19}

교문 앞에 `물마루 AI
연구소'라고 적힌 하얀 트럭이
서 있었다.

운전석 창문이 내려갔다. 엄마
쪽지 속 그 오삼촌이었다.
친척은 아니고, 성이 오씨라
다들 오삼촌이라 불렀다.

``고단우? 엄마한테 문자
받았다. 타라.''

조수석 문을 열자 귤 냄새가
풍겼다. 오삼촌은 출발도 전에
하품부터 했다.

``소장님이 꼭 그대로 전하래.
미리 발표해서 미안해요. 해 줄
거죠?''

물어보는 건 발표하기 전에 하는
거 아닌가. 단우는 안전벨트만
당겼다.

``탁 박사님은\ldots{} 무서운
분이에요?''

``무섭긴. 커피를 물처럼 마시는
사람이지.'' 오삼촌이 또
하품했다. ``왜, 혼날 일
했냐?''

``그냥 물어본 거예요.''

빨간불에 오삼촌이 휴대폰을
봤다. 화면이 어두워 조수석에선
글자 하나 안 보였다.

``단우야, 물고기 중에 제일
공부 잘하는 물고기가 뭔지
아냐?''

``고등어요.''

``\ldots 아는구나.''

``2학년 때 들었어요.''

``요즘 애들은 빨라.''

\Needspace{4\baselineskip}\begin{center}◇\end{center}\nopagebreak

연구소는 옛 생선 창고였다.
`얼음 반입구' 글씨 아래
모니터들이 줄지어 있었다.

입구 벽 지도에 빨간 점 세 개가
찍혀 있었다.

``우리 귀 세 개. 물속
마이크야. 방파제 마이크, 큰여
마이크, 등대 마이크.''

``큰여가 뭔데요?''

``물속 바위. 해녀들 일하는
데.''

세 점을 이으면 삼각형이었다.
그 바깥 남쪽에
`그물먹는여'라는 바위 무리가
있었다.

``저기는 왜 마이크가 없어요?''

``바위투성이라 배가 못 들어가.
그래도 조용한 날엔 등대
마이크로 그쪽 소리도 넘어와.''

오삼촌이 지도를 손등으로 쳤다.

``엔진 끄면 다 들려. 사람 말도
그렇고.''

단우는 예의상 ``하하'' 웃어
주었다.

\Needspace{4\baselineskip}\begin{center}◇\end{center}\nopagebreak

복도 끝 문틀에 머그잔 든
사람이 기대 있었다.

``단우 학생? 탁은재예요.''

금요일 밤 잠금 화면의 그
이름이었다.

``이제부터 말 편하게 할게.''
탁 박사가 머그잔을 들어
보였다. ``연료 다섯 잔째.''

``커피를 연료라고 해요?''

``사람도 뭘 넣어야 돌아가지.''

화이트보드 위에 액자가 걸려
있었다.

\begin{Verbatim}[samepage=true,frame=single,framesep=2.5mm,rulecolor=\color{gray},formatcom=\logfont]
물마루 AI 연구소 원칙
1. 녹음한 소리를 바다에 틀지 않는다
2. 모르는 건 모른다고 쓴다
3. 이름표에는 붙인 사람 이름을 남긴다
\end{Verbatim}

세 번째 줄에서 단우는 눈을
화이트보드로 내렸다.

탁 박사는 앉으란 말도 없이
네모 넷을 그렸다.

``바당이는 공장이야. 바다
소리가 들어가서 사람 문장이
나오는. 칸은 네 개.''

첫째 칸에 `듣기'를 쓰자
스피커에서 높은 휘파람 같은
소리가 났다. 딸깍딸깍딸깍,
볼펜 꼭지 누르는 소리.

``앞에 건 휘슬. 돌고래끼리
주고받는 소리야.''

``뒤에 건요?''

``클릭. 어둠 속 손전등 같은
거야. 소리를 쏘고, 되돌아오는
걸 듣고 먹이를 찾지. 첫째
칸에서는 휘슬만 오려 내.''

둘째 칸에는 `나누기'.

``오려 낸 휘슬을 모양이 같은
것끼리 나눠 담아. W-01 상자,
W-02 상자, 이런 식으로.''

``모양이 다 달라요?''

``여기 사는 남방큰돌고래도
저마다 자기만 내는 휘슬이
있다는 연구가 있어.
이름처럼.''

``돌고래가 출석도 불러요?''

``그건 아직 모르는 거고.''

셋째 칸에는 `짝 맞추기'.

``여기가 너희 라벨러들 칸이야.
이 휘슬이 날 때 사람들이 무슨
상황이라고 붙였는지 세서 짝을
맞춰.''

옆 모니터, 방파제 끝 관측
카메라 화면에 지느러미 끝이
뜯긴 회색 등이 솟았다.

``012번이네. 또.''

``불쑥이에요. 할머니가 그렇게
불러요.''

오삼촌이 문간에 고개를
내밀었다.

``돌고래가 제일 잘하는
발차기가 뭔지 알아?
돌려차기.''

아무도 웃지 않았다. 불쑥이마저
물속으로 들어가 버렸다.

탁 박사가 말했다. ``덕분에
화면은 비었네요.''

``네 번째 칸은요?''

마커가 넷째 칸에 `문장
만들기'를 썼다.

``앞 칸에서 나온 짝을 사람
말로 바꾸는 칸.''

``그럼 `돌려줘'는 몇 번째
칸에서 나온 거예요?''

마커가 멈췄다. 탁 박사는 그
칸에 물음표를 그렸다.

``여기. 이 칸은 빈칸 채우는
받아쓰기야.''

``받아쓰기는 선생님이 불러
주는 걸 쓰는 거잖아요.''

``그러니까. 아무도 안 불러
줬는데 쓰는 받아쓰기.''

``그럼\ldots{} 지어낸
거예요?''

``그럴듯하게 채운 거지.
맞는지는 아무도 몰라.''

\begin{Verbatim}[samepage=true,frame=single,framesep=2.5mm,rulecolor=\color{gray},formatcom=\logfont]
[듣기] -> [나누기] -> [짝 맞추기]
-> [문장 만들기 ?]
\end{Verbatim}

금요일 밤 본 흰 글자 세 개가
저 물음표 하나에서 나왔다.

``네 번째 칸은 내가
반대했어.'' 탁 박사가 마커
뚜껑을 닫았다. ``소장님이
붙이신 칸이야.''

\Needspace{4\baselineskip}\begin{center}◇\end{center}\nopagebreak

``뉴스 구석에 있던 31\%는 무슨
뜻이에요?''

``그 얘긴 네 라벨 다시 보고
나서 하자.'' 탁 박사가
노트북을 돌렸다. ``W-07이라는
휘슬 기록이야.''

\begin{Verbatim}[samepage=true,frame=single,framesep=2.5mm,rulecolor=\color{gray},formatcom=\logfont]
W-07 | 라벨 기록 일부
03:02:04 | 배 있음 | 분리 위협
03:02:08 | 배 있음 | 분리 위협
05:40:31 | 배 없음 | 먹이 활동
09:12:15 | 배 없음 | 휴식
\end{Verbatim}

``바당이는 이런 줄을 수만 개
세어서, W-07은 `분리 위협'이랑
짝이라고 배웠어.''

단우의 아이디로 붙은 `분리
위협'이, 넷째 칸을 지나
`돌려줘'가 됐다.

날짜별 라벨 수 막대는 한 달
전쯤부터 키가 컸다. 금색
왕관이 붙은 무렵이었다.

``밤마다 이삼백 개씩 붙였더라.
대부분 새벽에.'' 탁 박사가
단우를 똑바로 봤다. ``어떻게
했어?''

``빨리 들었어요.''

생각보다 대답이 먼저 나왔다.

탁 박사가 고개를 끄덕였다.
믿는 건지 다 안다는 건지 알 수
없었다.

``우리 시민 과학 소년
왔어요?''

봉 소장이었다. 오늘은 줄무늬
넥타이에 돌고래들이 헤엄쳤다.
세 번째 넥타이. 컵라면은 채아
쪽으로 기울었다.

``미리 발표해서 미안해요.''
소장이 단우 손을 덥석 잡았다.
``해 줄 거죠?''

트럭에서 들은 말과 한 글자도
다르지 않았다.

``생방송에서 저는 뭘
하는데요?''

``어렵지 않아요. 라벨을 어떻게
붙였는지 이야기해 주면 돼요.''

탁 박사가 머그잔을 소리 나게
내려놨다.

``아, 그리고 그날 푸른코드
대회 심사위원장님도 스튜디오에
계세요. 좋은 기회잖아요.''

``대회는 자기 손으로 만든 걸
내는 건데요.''

``그러니까요! 자기 손으로 붙인
라벨로 온 국민 앞에 서는 열두
살. 최고의 자기소개죠.''

칠판 귀퉁이 날짜가 떠올랐다.
7월 16일 토요일. 스튜디오 조명
아래 심사위원장이 박수 치며
이름을 부른다. 고단우 군. 화면
속에만 있던 서울이 한 칸
다가왔다.

``그래서, 해 줄 거죠?''

단우는 이번엔 고개를 끄덕였다.

``좋아요! 국민이 감동할
거예요.''

소장은 성큼성큼 나갔다. 탁
박사는 화이트보드 물음표를
지우지 않았다.

\Needspace{4\baselineskip}\begin{center}◇\end{center}\nopagebreak

철문 밖은 주황색이었다.
오삼촌은 운전석에서 하품
중이었다.

``단우 학생.''

탁 박사는 빈 머그잔을 감싸
쥐었다.

``이 문장, 과학자들 모이는
학회엔 못 내. 그런데 소장님은
온 나라 생중계로 낸대. 네
라벨, 한 번만 다시 봐 줄래?''

단우는 입을 열었다가 다시
닫았다. 그리고 트럭에 올랐다.

끝나면 전화해. 엄마 쪽지가
떠올라 번호를 찾는데 알림이
먼저 왔다. `내 라벨 전체
보기'가 열렸고, 라벨은
\mbox{12,408개였다.}

엄마 번호를 누르려던 엄지가
멈췄다. 오늘 교실에서 `1만 2천
번'이라고 자랑한 숫자였다.
단우는 그 숫자를 둘로 쪼갤 수
있었다. 1월부터 하나씩 귀로
듣고 붙인 \mbox{3,100개.}
그리고 단우가 자는 동안 붙은
\mbox{9,308개.}

\chapterartpage{images/pagefill/ch04-page.png}{images/pagefill/ch04.png}

\section{채아가 자른
장면}\label{uxcc44uxc544uxac00-uxc790uxb978-uxc7a5uxba74}

\dateline{6월 27일 월요일 · 방송 특집까지 D-19 · 채아}

채아는 월요일 해 질 녘, 연구소
언덕길에서 휴대폰으로 바다를
당겨 보고 있었다. 아빠 배는
사흘째 묶여 있었다.

``돌핀로그 오늘의 바다. 돌고래
없음. 아빠 배도 없음.''

언덕 위 연구소 철문이 끼익
열렸다. 남자아이 하나와 머그잔
든 어른이 나왔다.

``고단우?''

어른이 뭐라 하자 단우는 입을
열었다가 다시 닫았다. 소리는
언덕길까지 닿지 않았다.

단우가 탄 트럭이 채아 옆을
지나갔다.

``야, 고단우!''

단우는 무릎 위 휴대폰만 보고
있었다. 엄지가 화면 위에 뜬
채였다.

``쟤 연구소 편이구나.''

채아는 녹화를 멈췄다. 정작
바다는 한 컷도 못 건졌는데.

집 식탁에 빨간 도장 찍힌
종이가 펼쳐져 있었다. 채아가
아는 낱말은 `배' 하나뿐이었다.

아빠는 식탁 끝에서 휴대폰을
보고 있었다.

\begin{Verbatim}[samepage=true,frame=single,framesep=2.5mm,rulecolor=\color{gray},formatcom=\logfont]
곰새코지 3호 | 예약 관리
예약 취소 | 23건
\end{Verbatim}

``아빠, 나 오늘 연구소 앞에서
고단우 봤어.''

취소 줄은 내려도 내려도
이어졌다.

``걔 그 생방송에 나간대.
연구소 사람들이랑.''

아빠는 그 종이를 반으로 접어
휴대폰 밑에 깔았다.

``밥 먹자.''

``내일은 배 나가?''

``모르겠다.''

``모레는?''

``먹어라.''

``나 오늘 구독자 또 늘었어.''

``\ldots 그래.''

저녁 내내 아빠는 한숨만
쉬었다. 국그릇이 빌 무렵, 한숨
조회수는 두 자리가 됐다.
영상이었으면 망한 숫자인데,
한숨으론 신기록.

방에서 채널 관리 화면을
열었다. 인기 영상 맨 위엔
\texttt{우리\ 배한테\ 화났대요?},
그 아래엔
\texttt{말젯이\ 첫\ 점프!(6월\ 4일)}가
있었다.

맨 위 영상은 요즘 열지 않았다.
채아가 누른 건 둘째 줄.

6월 내내 반 단톡방에 올리던
영상이었다. 반 애들은 그때마다
`귀엽다'를 달아 주었다.
고단우만 한 번도 대꾸가
없었다.

어미 옆에서 새끼가 통 튀어
올랐다. 말젯이. 단우 할머니가
붙인 별명이었다. 언제부턴가
새끼는 화면에 들어오지 않았다.

\Needspace{4\baselineskip}\begin{center}◇\end{center}\nopagebreak

목요일 방과 후, 포구는
토요일보다 시끄러웠다.
돌고래쉼표 사람들이 한 줄로
피켓을 들었다.

\begin{Verbatim}[samepage=true,frame=single,framesep=2.5mm,rulecolor=\color{gray},formatcom=\logfont]
접근 거리를 지켜 주세요
\end{Verbatim}

탈의장 앞엔 해녀들이, 맞은편엔
선장들이 모였다.

``돌고래가 돌려달라고
했다잖아!''

``그게 우리 배 보고 한 말인지
어떻게 알아!''

아빠는 선장들 맨 뒤에서 팔짱을
꼈다.

``양 선장, 한마디 해 봐!''

아빠는 팔짱을 풀지 않았다.

채아는 휴대폰을 높이 들었다.

``안녕, 돌핀로그 채아예요---
지금 곰새코지 포구는요---''

화면 귀퉁이로 김복자 회장님이
들어왔다. 토요일에 카메라들을
본 척도 않던, 바로 그
해녀회장님. 회장님은 성게가
불룩한 그물 자루, 망사리를
내려놓았다. 그러고는 채아
휴대폰을, 그다음엔 채아를
똑바로 보았다.

채아 입이 `요' 모양으로
굳었다.

회장님이 돌아서자 채아는
휴대폰을 가슴께로 내리고
카메라를 바깥으로 돌렸다. 녹화
버튼은 끄지 않았다.

조금 떨어져 토요일의 그
활동가, 윤나래가 방송
중이었다.

``예전에도 이 바다에서
돌고래를 잡아 수족관에 판 일이
있었어요. 새끼가 그렇게 된 건
아닌지 누가 알아요?''

회원 하나가 윤나래 팔을
붙잡았다.

``나래 씨, 그건 우리 단체 생각
아니에요.''

``제 개인 방송이에요.''

``그래도 우리 피켓
앞이잖아요.''

가슴께 휴대폰은 그 장면을
처음부터 끝까지 담았다.

\Needspace{4\baselineskip}\begin{center}◇\end{center}\nopagebreak

금요일 방과 후, 구멍가게 앞
평상에 고단우가 혼자 앉아
있었다. 컵라면 뚜껑 위
휴대폰을 켰다 껐다 했다.

지나치는데 가게 텔레비전
소리가 들렸다.

``\ldots 특집 생중계에는 시민
과학 소년 고단우 군이
함께합니다\ldots.''

``축하해, 유명인.''

고개를 든 단우 눈 밑이
거뭇했다.

``\ldots 왔냐.''

``왔다. 연구소 편 구경하러.''

``월요일에 봤어. 옆으로
지나가면서 쳐다보지도
않더라.''

``그때 뭐 좀 보느라고.''

``그래, 바쁘시겠지. 생방송도
하시고.''

``나가고 싶어서 나가는 거
아니야.''

``그럼 왜 나가는데?''

단우는 젓가락 포장지만 감았다
풀었다.

``우리 아빠 배, 이번 주에 한
번도 못 떴어. 댓글에 `선장
아빠 처벌해라'가 하루에 몇
개씩 달려.''

``그 댓글, 네 영상에 달린
거잖아.''

채아 귀가 뜨거워졌다.

``네 아이디 밑에 뜬 문장
때문에 달린 거잖아.''

``그건 바당이가 번역한---''

``또 바당이. 너는 맨날 바당이
뒤에 숨더라.''

``너네 연구소 때문에 우리 집
망해.''

``데이터가 그렇게 나온 거야.''

``데이터?'' 채아가 코웃음을
쳤다. ``그거 네가 붙인 거잖아.
새벽 세 시에.''

단우 손가락이 멈췄다. 휴대폰이
평상에 툭 떨어졌다. 채아는
처음으로 단우를 말로 이겼다고
생각했다.

``확률적으로는---''

``구독 취소.''

채아는 돌아섰다. 유리문에 비친
단우는 휴대폰 옆에 그대로 앉아
있었다. 컵라면에서 김만
올라왔다.

\Needspace{4\baselineskip}\begin{center}◇\end{center}\nopagebreak

그날 밤, 채아는 목요일 포구
영상을 하나씩 잘라 냈다.

\begin{Verbatim}[samepage=true,frame=single,framesep=2.5mm,rulecolor=\color{gray},formatcom=\logfont]
포구 시위 | 편집
해녀 선장 고함     | 자름
회장님이 본 장면   | 자름
윤나래 개인 방송   | 남김
"단체 생각 아니에요" | 자름
\end{Verbatim}

회장님 얼굴은 어차피 못
올렸다. 그 눈빛이 찍지 말라고
했으니까. 나머지는 아빠 때문에
잘랐다.

남은 영상은 조용한 피켓 줄로
시작해 윤나래 목소리로 끝났다.

``\ldots 새끼가 그렇게 된 건
아닌지 누가 알아요?''

그 말대로면 새끼가 안 보이는
건 아빠 배 탓이 아니었다.

채아는 이런 일을 잘했다. 단체
사진에서 웃는 얼굴만 남기고
잘라 내면, 아무도 안 싸운
날처럼 보였다.

제목 칸엔
\texttt{포구\ 시위\ 현장\ (전체)}라고
쳤다.

아빠가 문틀에 기대 서 있었다.

``안 자냐.''

``편집 중. 아빠한테 좋게
잘랐어.''

``\ldots 뭘 잘랐는데.''

``시끄러운 거. 싸우는 거 다.''

아빠 눈길이 제목 칸
`(전체)'에서 멈췄다. 입이 뭔가
말하려다 말았다.

``\ldots 일찍 자라.''

방문이 닫혔다.

첫 화면엔 아빠 배가 돌고래들과
멀찍이 달리는 장면이 필요했다.
채아는 6월 원본 폴더를 열었다.
파일 이름은 찍은 날짜와
시각이었다.

빨리감기로 넘기다 한 파일에서
눈이 멈췄다.

\begin{Verbatim}[samepage=true,frame=single,framesep=2.5mm,rulecolor=\color{gray},formatcom=\logfont]
0612_1640.mp4
\end{Verbatim}

물보라 사이로 어미와 새끼가
나란히 떠올랐다. 새끼
가슴지느러미 쪽에 초록색
무언가가 걸쳐 있었다.

채아는 빨리감기에서 손을 떼지
않았다.

화면 아래로 불쑥, 아빠 배
뱃머리. 새끼와 뱃머리가 한
화면에 꽉 찼다.

아빠 배가 가깝게 나온 컷은 늘
버리는 쪽이었다. `선장 아빠
처벌해라'가 떠오르니까.

휴지통 아이콘은 6월 내내 버린
원본들로 볼록했다. 파일을 집어
든 커서가 휴지통 위에서 잠깐
멈췄다.

채아는 그 파일을 늘 그랬듯
휴지통으로 끌어다 놓았다.
휴지통 비우기는 누르지 않았다.

\dobira{3장}{잠든 사이 라벨을 붙인 건 누구?}{images/pagefill/ch05.png}{images/pagefill/ch05-page.png}

\section{라벨봇의
규칙}\label{uxb77cuxbca8uxbd07uxc758-uxaddcuxce59}

\dateline{7월 2일 토요일 · 방송 특집까지 D-14}

닷새를 버틴 끝에, 토요일 아침
단우는 `내 라벨 전체 보기'
알림을 눌렀다. `분리 위협'만
골라 시각 순으로 세웠다.

02:51, 03:02, 03:17\ldots{}

9시면 자는 단우가 붙였을 리
없는 시각들. 열 줄에 여덟
줄꼴이었다.

단우는 labelbot.py를 열었다.
켜려는 게 아니라, 속을 보려고.

\begin{Verbatim}[samepage=true,frame=single,framesep=2.5mm,rulecolor=\color{gray},formatcom=\logfont]
만약 관광선_거리 < 100m 이고
   휘슬_횟수(1분) >= 5 이면
      라벨 = "어미-새끼 분리 위협"
아니면
      라벨 = "이동"
\end{Verbatim}

``배가 100미터 안에 있고,
1분에 휘슬이 다섯 번 이상이면
`분리 위협'. 아니면 `이동'.''

손으로 붙이며 알아챈 거였다.
`분리 위협'엔 휘슬이 잦았고,
근처엔 늘 관광선이 있었다.
그땐 똑똑한 지름길 같았다.

``관광선, 휘슬. 관광선,
휘슬\ldots.''

어디에도 `새끼'라는 말은
없었다.

단우는 아무렇지 않은 척
물었다.

\begin{chatthread}
\chatroomtitle{귀 기울임 · 연구소 메시지}
\chatbubble{탁 박사}{응. 왜?}
\end{chatthread}

단우는 `왜?'를 오래 보았다.

``배가 오면 휘슬이
많아지고\ldots{} 휘슬이 많으면
`분리 위협'\ldots. 그럼 배만
오면 무조건 `분리
위협'이잖아.''

봇이 잰 건 `새끼가 떨어졌다'가
아니라 `배가 왔다'였다.

``\ldots 배였네.''

봇은 공책에 글씨가
빽빽한지로만 숙제 검사하는
선생님이었다. 낙서로 채워도
빽빽하긴 마찬가지인데.

단우의 아이디로 붙은 `분리
위협'이 넷째 칸을 지나
`돌려줘'가 됐다. 처음부터 틀린
라벨이었다.

배가 와서 찍힌 라벨이 이제 배
탓이 됐다. `너네 연구소 때문에
우리 집 망해.' 채아 말이
따라왔다.

``그래도\ldots{} 휘슬이 많았던
건 진짜잖아.''

단우는 `왜?'에 답하지 않고
앱을 닫았다.

\Needspace{4\baselineskip}\begin{center}◇\end{center}\nopagebreak

토요일 오후, 단우는 컵라면을
사러 구멍가게에 갔다.

``유명인, 오늘은 혼자냐?
채아랑 싸웠냐? 목소리가 가게
안까지 들리던데.''

``\ldots 안 싸웠어요.''

아저씨가 봉지에 컵라면 하나를
더 넣으려다 멈췄다.

``유명인 서비스는 웃는
사람한테만이다.''

단우는 입꼬리를 억지로 올렸다.
``됐죠?''

아저씨가 혀를 찼다.
``\ldots 그것도 웃은 걸로 쳐
주마.''

단우는 어제 채아가 돌아선
평상에 앉았다. 배가 보이자
등을 돌렸다.

`휘슬 상자별 보기'를 눌렀다.
막대가 길수록 라벨이 많다는
뜻.

\begin{Verbatim}[samepage=true,frame=single,framesep=2.5mm,rulecolor=\color{gray},formatcom=\logfont]
내 라벨 | 휘슬 상자별
상자      | 배 있음  | 배 없음
W-07      | ======== | =====
다른 상자 | ===      | ===
미분류 #0 | .        | .
\end{Verbatim}

W-07 줄이 길었다. 배가 없을
때도. 단우 손가락이 그 줄에서
멈췄다가 내려갔다.

`미분류 \#0' 물음표를 눌렀다.

\begin{Verbatim}[samepage=true,frame=single,framesep=2.5mm,rulecolor=\color{gray},formatcom=\logfont]
미분류 #0
어느 상자에도 안 맞는 소리
대부분 잡음(딱총새우 등)
\end{Verbatim}

들어 보기를 누르자 따닥따닥,
지지직.

``라디오 채널 안 맞냐?
지지직거리게.''

``새우 소리래요. 딱총새우.''

``새우도 소리를 내?''

``그냥 잡음이래요.''

딱총새우 잡음이겠지. 새우
소리까지 들어 줄 시간은
없었다. 단우는 그 줄을 접어
두었다.

\begin{Verbatim}[samepage=true,frame=single,framesep=2.5mm,rulecolor=\color{gray},formatcom=\logfont]
귀 기울임 | 지난주 리더보드
1  새벽세시감귤
2  danu_0101
\end{Verbatim}

금색 왕관은 이제 새벽세시감귤
옆이었다. 생중계 날 밤
\texttt{\^{}\^{}} 댓글을 달던
그 2위.

``사람이 이걸 손으로 한다고?''

``뭘 손으로 해? 그거 몇
등이냐?''

``\ldots2등이요. 1등
뺏겼어요.''

봇을 멈춘 그날 밤부터라는 말은
안 나왔다.

``누구한테?''

``새벽세시감귤이요.''

``귤한테 졌냐? 그래서 표정이
그 모양이구나.''

가게 텔레비전에 예고가 나왔다.

``\ldots 특집 생중계에는 시민
과학 소년 고단우 군이
함께합니다\ldots.''

``십육일엔 저거 크게 틀어 놓을
거다. 유명인 나오는
날이니까.''

``\ldots 안 그래도 돼요.''

닉네임부터 새벽 세 시였다.
저쪽도 봇이 틀림없다. 단우는
괜히 `감귤' 글자를 노려보았다.

1위 화면은 늘 육지 공사장
아빠에게 보냈다. 답장은 늘
엄지 하나. 오늘은 보낼 화면이
없었다.

뚜껑을 여니 면이 퉁퉁 불었다.

``버그다.''

\Needspace{4\baselineskip}\begin{center}◇\end{center}\nopagebreak

저녁에 온 탁 박사 메시지를
단우는 미리보기만 보고 껐다.
9시가 넘어도 잠이 오지 않았다.

``단우야, 안 자냐?''

``자요, 지금 자요.''

``자는 애가 대답은 잘한다.''

단우는 불을 끄고 메시지를
열었다.

\begin{chatthread}
\chatroomtitle{귀 기울임 · 연구소 메시지}
\chatbubble{탁 박사}{아까 휘슬은 왜 물어봤어?\\라벨도 다시 봤어?}
\end{chatthread}

생중계에서 누가 이렇게
묻는다면.

``고단우 군, 라벨은 어떻게
붙였어요?''

정직한 대답은 바탕화면에
있었다. labelbot.py. 이름부터
너무 정직했다.

``숙제.py는\ldots{} 너무
수상하고.''

``연습.py는 너무 착해
보이고.''

labelbot.py라는 이름이 보이면,
누가 봐도 라벨봇인 걸 알 것
같았다.

결국 이렇게 쳤다.

\mbox{\texttt{test\_old.py}}

단우는 이름을 숨긴 파일을
`옛날 것' 폴더 맨 아래로
끌어다 놓았다.

``어차피 다시 안 켤 거니까.''

\texttt{사실\ 제가}까지 쳤다가
한 자씩 지웠다.

결국 단우는 ``제 라벨은
문제없어요.''라고 보냈다.

`입력 중\ldots'이 떴다가
사라지고, 다시 떴다가
사라졌다.

잠시 뒤 탁 박사에게서
``그래.''라는 답이 왔다. 고작
한마디였다. 빈 머그잔을 감싸
쥐던 탁 박사가 떠올랐다.

손가락이 책상을 두드렸다.
\mbox{\texttt{0, 1, 0, 1}}.
리듬 게임이면 퍼펙트였을 박자.
단우는 두 손을 엉덩이 밑에
깔았다.

\Needspace{4\baselineskip}\begin{center}◇\end{center}\nopagebreak

일요일 저녁, 할머니가 빨랫줄에
고무옷을 널며 말했다.

``낼 새벽 물에 들어간다. 일찍
자라.''

``저 9시면 자요.''

``자기는 무슨. 어젯밤에
문틈으로 파란 불빛이 다
새더라.''

``\ldots 게임 좀 했어요.''

``불 끄고 하는 게임도 있냐.''

방에서 단우는 W-07 녹음을
틀었다. 문제없다고 보냈으니
들어는 봐야 했다.

높은 휘파람 같은 소리.
올라갔다 내려가고, 다시
올라갔다가 길게. 잠깐 쉬고
똑같은 모양이 또 나왔다.

빨래 바구니를 든 할머니가
방문을 열었다.

``그 소리 안다.''

단우 손이 정지 버튼 위에서
멈췄다.

``\ldots 이 소리를요?''

``반달이 소리여.''

``반달?''

``지느러미 뒤쪽이 반달처럼
떨어져 나간 놈. 돌고래주.''

할머니는 소리가 끝날 때까지
꼼짝 않았다. ``한 번 더 틀어
봐라.''

``그럼 이 돌고래가 뭐라고 하는
건데요?''

``돌고래가 무사 사람 말을
해?''

``반달 소리인 건 어떻게
알아요? 데이터라도 있어요?''

``물때 봐야 알주.''

단우가 노트북을 돌렸다.

``할머니, 이게 생중계에서
`돌려줘'가 된 소리예요.''

할머니는 화면을 오래
내려다보다가, 문고리를 잡은 채
돌아보았다.

``알고 싶으민 낼 새벽 다섯 시.
불턱.''

\chapterartpage{images/pagefill/ch06-page.png}{images/pagefill/ch06.png}

\section{불턱에서 들은
이름}\label{uxbd88uxd131uxc5d0uxc11c-uxb4e4uxc740-uxc774uxb984}

\dateline{7월 4일 월요일 · 방송 특집까지 D-12}

새벽 4시 50분, 단우는 슬리퍼를
끌며 주황색 테왁을 낀 할머니를
따라 포구로 내려갔다.

바닷가 바위 위에 허리 높이
돌담이 나왔다. 바닥엔 불 피운
자국.

``여기가 불턱이다. 물에 들기
전에 모여서 어디로 들어갈지
정하고, 나와서는 불 쬐주. 딴
동넨 다 탈의장 쓰주만, 우린
아직 여기 모인다.''

돌담 안엔 해녀들이 둘러앉았다.
볼이 동그란 해녀 하나가 손뼉을
쳤다.

``덕순 삼춘네 손지가 바당엘 다
왔네! 해가 서쪽서 뜨켜.''

`덕순 삼춘'은 할머니, `손지'는
단우였다. 해녀들이 와르르
웃었다. 단우는 저도 모르게
하늘을 보았다. 해는 아직
뜨지도 않았는데.

``어제 큰여 바깥으로 곰새기
여남은 마리 지나갔주.''

``그럼 오늘은 큰여 안쪽으로
들자.''

정한 사람은 맨 안쪽 해녀회장,
김복자 회장이었다.

곰 새끼가 열 마리나 바다에?

젊은 해녀가 단우 쪽으로 몸을
숙였다.

``곰새기, 돌고래야. 삼춘들은
그렇게 불러. 나도 처음엔
몰랐어.''

곰새기. 곰새코지. 날마다
부르던 마을 이름에, 돌고래가
들어 있었다.

옆에서 하품을 참는 아이가
있었다. 단우네 학교 5학년
다원이었다.

``다원이 또 따라왔네.''

``따라온 거 아니에요. 엄마
데려다준 거예요.'' 다원이 젊은
해녀 등을 쳤다. ``맨날 내가
엄마 옆에서 설명하는데, 여기선
엄마가 다 알아들어.''

곰새기를 알려 준 해녀가 다원
엄마, 미숙이었다.

``느 저 바위에 앉아 있어라.
물엔 한 발짝도 들지 말고.''

``안 들어가요. 슬리퍼예요.''

``슬리퍼 아니어도 들지 마라.''

\Needspace{4\baselineskip}\begin{center}◇\end{center}\nopagebreak

6시가 넘자 해녀들이 헤엄쳐
나갔다. 검은 머리가 사라졌다
올라올 때마다 소리가 났다.

호오이---

휘파람 같기도, 긴 한숨 같기도
한 소리였다.

할머니는 일찍 나와 단우 옆에
앉았다.

``방금 그 소리 뭐예요?''

``숨비소리다. 물속에서 참은
숨을 한꺼번에 내쉬는 소리.''

호오이--- 할머니는 바다를
보지도 않았다.

``복자 올라왔다.''

멀리 회장 얼굴이 보였다. 잠시
뒤, 조금 짧은 소리.

``정임이.''

볼이 동그란 그 얼굴이 맞았다.

``안 보고 어떻게 알아요?''

``물때 봐야 알주.''

``또 물때요?''

올라오던 해녀 하나가 웃었다.

``삼춘은 말하기 싫으면 다
물때주.''

발소리만 듣고 엄마인 줄 알듯,
할머니는 숨비소리로 사람을
알았다. 사람마다 다른, 소리로
된 지문.

호오이---

이번엔 이름이 나오지 않았다.
할머니가 처음 고개를 돌려 미숙
얼굴을 확인했다.

녹음 버튼에 엄지를 올리는데,
소장 말이 떠올랐다.

`미리 발표해서 미안해요. 해 줄
거죠?'

물어보는 건 원래 발표하기 전에
하는 거 아닌가.

``할머니, 녹음해도 돼요?''

``뭐 하젠?''

``비교해 보려고요. 사람마다
숨비소리가 어떻게 다른지요.''

할머니는 휴대폰을 한 번, 단우
얼굴을 한 번 보았다.

``해라.''

할머니 소리는 됐다. 다른
소리는 할머니 것이 아니었다.
단우는 녹음 버튼에서 엄지를
뗐다.

\Needspace{4\baselineskip}\begin{center}◇\end{center}\nopagebreak

7시쯤 해녀들이 불턱으로
돌아왔다. 불을 피우자 젖은
고무옷마다 김이 모락모락.

눈이 쏠리자 단우 목소리가
기어들어 갔다.

``저\ldots{} 숨비소리 녹음해도
돼요? 비교해 보려고요.''

정임이 눈을 동그랗게 떴다.

``내 소리? 해라, 해. 테레비
나오는 거냐?''

``아니요. 저 혼자 보는
거예요.''

``에이, 아깝다. 화장하고
올걸.''

웃음이 잦아들자 회장이 입을
열었다.

``난 싫다.''

돌담 안이 조용해졌다.

``\ldots 네.''

회장은 말없이 엉덩이를 옆으로
옮겼다. 불 앞에 아이 하나 앉을
빈자리.

``나는 괜찮아.'' 장갑을 벗던
미숙이 단우를 올려다보았다.
``근데 나중에 나도 들려줘. 내
소리, 나는 한 번도 못
들었어.''

``네. 꼭요.''

``한 번 더 든다.'' 할머니가
테왁 끈을 고쳤다.

미숙만 쉬겠다며 바위에 남았다.
단우는 옆에서 녹음을 켰다.

``나는 돌고래 소리는 잘 몰라.
사람 소리는 조금 알아.''

호오이--- 코앞에서 회장 테왁이
솟았다. 그 소리가 휴대폰에
그대로 들어갔다.

단우는 녹음을 멈추고, 방금
들어간 끝부분을 골랐다.

\begin{Verbatim}[samepage=true,frame=single,framesep=2.5mm,rulecolor=\color{gray},formatcom=\logfont]
녹음 | 끝부분 선택됨
이 부분을 지울까요?
취소 | 지우기
\end{Verbatim}

회장은 싫다고 했다. 망설일
것도 없이, `지우기'.

다시 켜자 멀리서 웅웅 고깃배
소리. 반쯤 묻힌 정임의
호오이---.

배 소리가 멀어지자 먼 데서
소리가 왔다.

``삼춘 올라왔어.''

맨 바깥에 주황색 테왁 하나.
할머니였다.

``어떻게 알았어요?''

``삼춘 소리는 올라가서 오래
있다가 내려와.'' 미숙이 공중에
선을 그었다. 위로, 옆으로
길게, 아래로. ``삼춘 소리
들리면 나도 마음 놓여.''

``나 들어갈게. 내 소리도 잘
받아 둬.''

바위 앞에서 호오이--- 미숙의
소리가 휴대폰 속으로 또렷하게.

``단우, 학교 늦겠다!''

물에서 나온 미숙이 소리쳤다.
단우는 벌떡 일어나 뛰었다.

\Needspace{4\baselineskip}\begin{center}◇\end{center}\nopagebreak

그날 밤, 단우는 녹음을 소리
그림 프로그램에 넣었다.
할머니가 들어와 화면을 보았다.

``그건 또 뭐냐?''

``할머니 숨비소리요.'' 단우가
선을 짚었다. ``소리가 높아지면
선이 올라가고, 낮아지면
내려가요. 소리로 그린 산 능선
같은 거예요.''

\begin{Verbatim}[samepage=true,frame=single,framesep=2.5mm,rulecolor=\color{gray},formatcom=\logfont]
할머니 숨비소리
높 | ..========..
   | .=........=.
낮 | =..........=
\end{Verbatim}

할머니가 그 선을 짚었다.
아침에 미숙이 공중에 그린 선
그대로였다. 위로, 옆으로 길게,
아래로.

``올락내리락이네.''

``그래서 이렇게 적기로 했어요.
올라가면 U, 내려가면 D,
평평하면 F. 영어 첫 글자요.''

할머니가 혀를 찼다. ``숨에다
영어를 붙이네.''

단우는 확인해 둔 순서대로
이름을 붙이고 U, D, F를
적었다. 허락한 사람 소리만. 몇
번 들어도 모양은 그대로였다.

``복자는 없구나.''

``싫다고 하셨잖아요.''

정임 그림 하나는 아래쪽이
뿌옇게 번졌다.

``여긴 배가 지나가서 다
덮였어요.''

할머니 손가락이 멈췄다.

``내가 스무 살 적에, 같이 물에
드는 짝, 물벗이 하나 있었다.
큰여에서 숨 쉬러 올라오는데,
배가 바로 위에 와 있었다. 배
밑에서 도는 스크루에 다리를
크게 다쳤다. 그 뒤로 그 사람은
다시 물에 못 들었주.''

할머니가 손을 거두었다. 단우
눈이 마디 굵은 그 손에
머물렀다.

``그럼\ldots{} 배 안 오는 데는
없어요?''

``그물먹는여. 바위가 많아서
배가 못 들어가. 그 안쪽엔
바위틈 굴도 있다.''

``굴 안에도 가 봤어요?''

``물때 봐야 알주.''

할머니는 그 말만 하고 방을
나갔다. 이번엔 단우도 따지지
않았다.

혼자 남자 손가락이 W-07 파일로
갔다. `반달이 소리'도 그림으로
바꿔 표 맨 아래에 보탰다.

\begin{Verbatim}[samepage=true,frame=single,framesep=2.5mm,rulecolor=\color{gray},formatcom=\logfont]
누구    | 모양
정임    | F-U-D-D
정임    | F-U-D-D
할머니  | U-F-F-D
할머니  | U-F-F-D
미숙      | F-U-D-F
W-07    | U-D-U-F
\end{Verbatim}

같은 사람은, 같은 모양. 적어도
이 표에서는.

그럼 W-07은?

\dobira{4장}{W-07은 누구의 휘슬일까?}{images/pagefill/ch07.png}{images/pagefill/ch07-page.png}

\section{W-07의
정체}\label{w-07uxc758-uxc815uxccb4}

\dateline{7월 5일 화요일 · 방송 특집까지 D-11}

새벽 2시가 넘도록 단우는 눕지
못했다. 표 맨 아래 W-07 한
줄이 눈에 걸렸다. 같은 사람은,
같은 모양. 그럼 W-07은?

W-07은 돌고래 번호가 아니라,
바당이가 같은 모양 휘슬끼리
담아 둔 상자 번호였다. 단우는
그 상자만 골라 꺼내는 짧은
필터를 짰다. 필터는 많은 줄
가운데 보고 싶은 줄만 골라
내는 체 같은 프로그램이었다.
파일 이름은 listen.py.

실행하자 휘슬 줄이 주르륵
쏟아졌다. 줄마다 라벨과 배
기록이 붙어 있었다.

\begin{Verbatim}[samepage=true,frame=single,framesep=2.5mm,rulecolor=\color{gray},formatcom=\logfont]
W-07 | 모아 보기
상황      | 배   | 모양
먹이 활동 | 없음 | U-D-U-F
휴식      | 없음 | U-D-U-F
이동      | 있음 | U-D-U-F
분리 위협 | 있음 | U-D-U-F
\end{Verbatim}

``먹이 활동, \mbox{U-D-U-F.}
휴식, \mbox{U-D-U-F.} 이동,
\mbox{U-D-U-F.}''

``분리 위협\ldots{}
\mbox{U-D-U-F.}''

`분리 위협'은 배가 있던 줄에만
붙어 있었다.

``상황이 다 다른데, 왜
똑같지?''

단우는 숨비소리 표를 옆 창에
띄웠다. 성게를 땄든 빈손이든,
할머니 숨비소리는 늘
\mbox{U-F-F-D.} 먹이를 쫓든
쉬든, W-07은 늘
\mbox{U-D-U-F였다.}

연구소에서 들은 말이 떠올랐다.

``여기 사는 남방큰돌고래도
저마다 자기만 내는 휘슬이
있다는 연구가 있어.
이름처럼.''

그때 단우는 웃자고 되물었다.
``돌고래가 출석도 불러요?''

빈 방에 단우 목소리가 울렸다.

```돌려줘'가 아니었어. 이건
문장이 아니라\ldots{}
이름이잖아.''

단우는 W-07 줄을 차례로
틀었다. 먹을 때, 쉴 때, 배가
지나갈 때. 네 번 다, 같은
휘파람이었다. 그리고 처음부터
한 번 더.

창밖이 희뿌옇게 밝을 무렵,
할머니가 방문을 열었다.

``밤새 안 잔 거냐?''

``할머니, 정임 삼춘 숨비소리
들으면 `정임이' 하잖아요. 이
돌고래 소리도요, 뭘 하든---''

할머니는 화면을 잠깐
내려다보았다.

``그러니까 반달이 소리라
했주.''

단우는 입을 벌린 채 할머니
등만 바라보았다.

7시가 넘어 다원에게 메시지가
왔다. 단우는 listen.py 창만
찍어 보냈다. 숨비소리 표는
빼고.

\begin{chatthread}
\chatbubble{다원}{엄마가 녹음 잘 됐냐고 물어봐 달래}
\chatbubble{단우}{잘 됐어. 꼭 들려드릴게\\엄마 숨비소리 박사님께 보여 드려도 돼?}
\chatbubble{다원}{된대. 근데 listen.py는 뭐야?\\그냥 듣기.py라고 하든가}
\chatbubble{단우}{원래 영어로 짓는 거야}
\chatbubble{다원}{누가 그래?}
\end{chatthread}

단우는 휴대폰을 엎어 놓았다.
파일 이름은 그대로.

\Needspace{4\baselineskip}\begin{center}◇\end{center}\nopagebreak

화요일 저녁, 할머니는 망사리
터진 데를 깁고 있었다.

``할머니, 잠깐만 봐 주세요.''

``바늘 꿰는 중이다. 화면이나
대라.''

``그 휘슬이 울릴 때 연구소
배가 찍은 지느러미 사진이에요.
누군지 아세요?''

와이파이가 한 칸뿐이라 사진이
위에서부터 느릿느릿 내려왔다.
회색 지느러미 꼭대기가 막 보일
참이었다.

``반달.''

할머니는 힐끗 보고는 바늘귀로
눈을 내렸다.

``아직 반도 안 떴는데요.''

``반이면 됐다.''

``꼭대기만 보고 어떻게
알아요?''

``느도 문틈에 내 이마만 보이면
할머닌 줄 알 거 아니냐.''

사진이 다 뜨자 지느러미 뒤쪽이
보였다. 반달처럼 둥글게 떨어져
나간 자리.

다음 사진도, 그다음 사진도
할머니는 다 뜨기 전에 불렀다.
반달. 또 반달. 한 장에서는
혀부터 찼다.

``요건 불쑥이고. 또
끼어들었네.''

반달 옆에 끝이 뜯긴 지느러미가
불쑥. 오래된 사진일수록 반달
곁에 작은 지느러미가 자주
보였다.

할머니가 바늘 쥔 손으로 그
작은 지느러미를 짚었다.

``말젯이.''

``말젯이요? 말젯이가\ldots{}
돌고래였어요?''

``그럼 뭔 줄 알았냐.''

``\ldots 동네 해녀 삼춘인 줄
알았어요.''

생중계 날 밤 할머니가 한 말이
떠올랐다. ``말젯이가 요새 안
뵌다.''

W-07이 크게 울린 시각의 사진
속에는 거의 언제나 반달이
있었다.

``할머니, 내일 박사님한테
할머니 숨비소리 표 보여 드려도
돼요?''

할머니가 실 끝을 이로 똑
끊었다.

``해라.''

\Needspace{4\baselineskip}\begin{center}◇\end{center}\nopagebreak

수요일 저녁, 단우는 오삼촌의
트럭을 타고 연구소에 갔다.
화이트보드엔 지난번 네 칸이
그대로였다. 넷째 칸 물음표도.

탁 박사가 머그잔을 들어
보였다. ``연료 일곱 잔째.''

``지난번엔 다섯
잔이었잖아요.''

``오늘은 일이 좀 많았어.''

``박사님, 왼쪽은 숨비소리예요.
할머니랑 다원이 엄마요. 보여
드려도 된다고 허락받았어요.''

단우는 표의 정임 줄 두 개를
손바닥으로 가렸다.

``정임 삼춘 건 저 혼자 보기로
했거든요.''

``알았어.''

``사람마다 모양이 다르고, 같은
사람은 늘 똑같아요. 오른쪽
W-07은 뭘 하든
\mbox{U-D-U-F고요.}''

``그래서?''

``그래서 이건 문장이 아니라
이름 같은 거예요.''

탁 박사는 두 창을 훑더니
하나만 물었다.

``누구 소린지는?''

``연구소 개체 사진이요. W-07이
크게 잡힌 시각엔 거의 다
반달이 있었어요. 할머니는
사진이 반도 안 떴는데
맞혔고요.''

탁 박사의 입꼬리가 아주 조금
올라갔다.

``이런 걸 `시그니처
휘슬'이라고 불러. 이름 같은
휘슬. 남방큰돌고래 연구에도
나오고, 제주 무리에서도 보고가
있어.''

``그럼 W-07은 반달
이름이에요?''

``그건 아직 모르는 거고.''

``사진마다 반달이
있었다니까요.''

``반달 자기 이름일 수도 있고,
반달이 다른 돌고래를 부르는 걸
수도 있어. 그건 몰라.''

단우는 노트 맨 뒷장에 `모름'
칸을 만들었다.

\begin{quote}
\textbf{단우의 메모 \#2}
시그니처 휘슬 = 자기만 내는
휘슬. 이름 같은 것. W-07이
누구 이름인지는: 모름.
\end{quote}

고개를 드니 넷째 칸 물음표가
보였다.

``그럼 넷째 칸은\ldots{}
이름을 문장으로 바꿔 버린
거네요.''

탁 박사는 대답 없이 W-07
라벨별 막대를 띄웠다. `분리
위협' 막대만 혼자 길쭉했다.

``근데 이상하다.'' 탁 박사가
그 막대를 손끝으로 짚었다.
``아무 때나 나는 소리가, 왜
`분리 위협'에만 이렇게
몰렸지?''

단우의 귀에서 웅 소리가 났다.
그 막대의 열에 여덟은 잠든
사이 봇이 붙인 것이었다.
단우는 노트를 덮었다. 방금 쓴
`모름' 두 글자 위로.

``\ldots 잘 모르겠어요.''

탁 박사가 단우를 보았다.
``어떻게 했어?'' 하고 묻던
그날과 같은 눈이었다.

``단우야, 가자. 해 떨어진다.''

문간에서 오삼촌이 하품하며
트럭 열쇠를 흔들었다. 탁
박사는 더 묻지 않았다. 끝까지.

\Needspace{4\baselineskip}\begin{center}◇\end{center}\nopagebreak

목요일 밤, 단우는 W-07이 하루
평균 몇 번 울렸는지 막대로
세웠다.

\begin{Verbatim}[samepage=true,frame=single,framesep=2.5mm,rulecolor=\color{gray},formatcom=\logfont]
W-07 | 하루 평균 횟수
12일까지 | ===
13일부터 | =========
\end{Verbatim}

막대는 6월 13일에서 훌쩍
뛰었다. 그날부터 W-07은 하루
평균 세 배로 울렸다.

``13일에\ldots{} 무슨 일이
있었지?''

탁 박사의 메시지가 떠올랐다.

``서로 떨어졌을 때 늘어.''

단우는 말젯이가 반달 곁에 찍힌
날의 휘슬과, 13일 뒤의 휘슬을
겹쳐 보았다.

``모양은 똑같은데\ldots{} 13일
뒤 건 산이 낮아졌네.
오르내림도 작고.''

말젯이 소리일까, 반달
소리일까. 단우는 `모름' 칸에
보탰다. 반달? 말젯이?

마이크별로 나눠 보다가 멈췄다.
세 마이크에 다 잡히던 그
휘슬이, 어떤 날은 등대 마이크
하나에만. 단우는 노트 귀퉁이에
적고 넘어갔다.

\dateline{등대 쪽만? 나중에}

그래프를 보내자 답이 금방
왔다.

\begin{chatthread}
\chatroomtitle{귀 기울임 · 연구소 메시지}
\chatbubble{탁 박사}{비슷한 돌고래를 연구한 자료가 있어.\\새끼 곁에선 어미 휘슬이\\더 높이, 더 크게 오르내린대.\\그래도 이것만으로는 단정 못 해.}
\end{chatthread}

아기 앞에서는 어른들 목소리가
저절로 높아지잖아. 그런 걸까?

``그럼 낮아진 건\ldots{} 곁에
새끼가 없어서?''

단우는 `모름' 칸의 `반달?'
밑에 줄을 그었다. 물음표는
지우지 않았다.

거실에서 할머니가 귤을 까고
있었다.

``할머니, 말젯이 마지막으로 본
날이 언제예요?''

귤 까던 손이 멈췄다.

``13일. 월요일 물질 날.''

``확실해요?''

``물질 날은 안 잊는다.''

그래프가 뛰어오른 바로
그날이었다.

``말젯이는 무슨 뜻이에요?''

``셋째. 반달이 셋째 새끼라.''

할머니가 귤 한 쪽을 단우 손에
쥐여 주었다. 단우는 그 귤을
입에 넣지 못했다.

6월 13일부터 스물네 날. W-07이
누구 이름인지 단우는 몰랐다.
다만 숫자가 바뀐 날짜는
알았다.

그리고 탁 박사가 손끝으로 짚던
그 길쭉한 막대. 바다가 아니라
단우의 노트북에서 자란 막대.

`옛날 것' 폴더 맨 아래에는
이름을 바꿔 숨긴 봇,
test\_old.py가 있었다. 단우는
그 이름을 오래 보았다. 연구소
메시지 창을 열어 탁 박사에게
``내일 아침에 드릴 말씀이
있어요.''라고 보냈다.

\chapterartpage{images/pagefill/ch08-page.png}{images/pagefill/ch08.png}

\section{맞은 것과 안
것}\label{uxb9deuxc740-uxac83uxacfc-uxc548-uxac83}

\dateline{7월 8일 금요일 · 방송 특집까지 D-8}

아침 7시 반, 연구소 주차장. 탁
박사는 머그잔을 든 채 차에
기대어 있었다.

``일찍 왔네. 할 말이 뭔데?''

밤새 연습한 첫마디는 떠오르지
않았다. 단우는 노트북을 보닛에
올리고 손부터 움직였다.

`옛날 것' 폴더. 그 맨 아래,
test\_old.py.

``이거, 원래 이름은
labelbot.py예요.''

단우는 화면을 탁 박사 쪽으로
돌렸다.

``제가 짠 봇이에요. 배가
가깝고 휘슬이 많으면 `분리
위협'을 붙여요.''

``새벽 3시 2분 8초도 이거예요.
제가 자는 동안 붙은
\mbox{9,308개도요.}''

``\ldots\ldots.''

``빨리 들었다고 한 거,
거짓말이었어요.''

말하는 내내 단우의 손가락은
가만히 있었다. 0도, 1도
두드리지 않았다.

탁 박사는 머그잔을 입가로
가져가더니, 마시지 않고 차
지붕에 내려놓았다.

``\ldots 그래서 `분리
위협'에만 몰렸구나.''

화난 목소리는 아니었다.

``말해 줘서 고마워.''

따지는 말도, 혼내는 말도 오지
않았다. 단우 어깨에서 힘이 쑥
빠졌다.

``오늘 소장님 보고는 데이터만
얘기해. 네가 찾은 거. 이름
얘기.''

``봇 얘기는요?''

``그건 내가 생각할게. 네가
혼자 소장님 앞에서 할 얘기는
아니야.''

탁 박사는 노트북을 덮어
단우에게 안기고 철문 쪽으로
갔다. 머그잔은 차 지붕 위에
그대로 남았다.

\Needspace{4\baselineskip}\begin{center}◇\end{center}\nopagebreak

소장실 옷걸이엔 돌고래
넥타이가 줄줄이 걸려 있었다.
소장이 오늘 맨 건 은빛 반짝이
넥타이. 돌고래들이 번쩍번쩍.

넥타이 내기는 채아 쪽으로
기울었다. 말해 줄 수는
없었지만.

``우리 시민 과학 소년!
아침부터 웬일이에요?''

``단우가 찾은 게 있어요.'' 탁
박사가 말했다. ``직접 들어
보세요.''

단우는 W-07 표만 띄웠다.
숨비소리 표는
`박사님한테'까지만 허락받았다.

``W-07은 먹을 때도, 쉴 때도,
배가 있어도 없어도 모양이
똑같아요.''

``그래서요?''

``그래서 이건 문장이 아니라,
이름 같은 거예요.''

``시그니처 휘슬이라고 해요.''
탁 박사가 덧붙였다.
```돌려줘'라는 뜻으로 볼
근거는 없어요.''

소장이 넥타이를 한 번
쓸어내렸다. ``그래도 확신도
31\%면, 열 번에 세 번은
맞는다는 거잖아요.''

``31\%는 맞을 확률이
아니에요.'' 탁 박사가 말했다.
``닮은꼴 점수예요.''

연예인 닮은꼴 앱이 `31\%
닮았다'고 해도, 열 번에 세 번
그 사람인 건 아니다. 단우는
얼른 적었다.

\begin{quote}
\textbf{단우의 메모 \#3}
확신도 = 닮은꼴 점수. 맞을
확률이 아님.
\end{quote}

``흥미롭네요. 그런데 하나만
봐요.'' 소장이 몸을 기울였다.
``분리 위협 라벨이 실제로
맞았잖아요. 새끼가
없어졌다면서요.''

방 안이 조용해졌다.

``그럼 `돌려줘'도 틀린 게
아니죠. 국민이 감동했잖아요.''

``새끼가 없어진 거랑 그
라벨은---''

``단우 학생.'' 소장이 탁
박사의 말을 끊었다. ``시험 볼
때 찍은 문제도, 맞으면
동그라미 받죠?''

단우는 하마터면 고개를 끄덕일
뻔했다. 지난 수학 시험 마지막
문제를 찍었는데 동그라미였다.

``그거예요. 어떻게 맞혔든,
맞은 건 맞은 거예요.''

뭔가 이상했다. 어디가
이상한지는 짚을 수가 없었다.

``방송은 예정대로 해요. 이번
방송이 잘돼야 우리 연구소가
계속 바다를 들어요.'' 소장이
넥타이 매듭을 고쳐 맸다. 눈은
단우 쪽으로 오지 않았다.
``라벨이 문제면, 그건 시민
데이터 쪽 문제고요.''

단우는 노트북 모서리를 꽉
쥐었다. 시민 데이터. 그 시민은
단우였다. 진짜 붙인 쪽은 이
노트북 안에 있었다.

탁 박사가 빈손으로 책상을
짚었다.

``미성년자 아이디를 생중계에
낸 것부터 잘못이었어요.''

탁 박사의 목소리가 이렇게 커진
건 처음이었다.

``출연도 그래요. 이름부터
발표하시고, 열두 살한테 고개
한 번 끄덕이게 한 게
다였잖아요. 무슨 일인지 다
듣지도 못하고 한 끄덕임은
동의가 아니에요. 보호자께는
여쭤보셨어요?''

소장이 넥타이 매듭에서 손을
내렸다.

``\ldots 어머님께는 오늘
안내드릴 거예요. 시간이 없어서
순서가 조금 바뀐 거고요.''

탁 박사가 단우를 보더니
목소리를 낮췄다.

``단우야, 먼저 나가 있을래?''

따라 나온 탁 박사가 머그잔을
들어 한 모금 마셨다.

``연료가 다 식었네.''

``처음 온 날, 철문 앞에서 내가
그랬지. `네 라벨, 한 번만 다시
봐 줄래?' 그 문장, 방송에 내면
안 된다고. 그 말을 소장님한테
먼저 했어야 했어. 내가 못 한
말을 너한테 시켰네.''

탁 박사가 턱으로 언덕 아래를
가리켰다.

``학교 가. 늦겠다.''

\Needspace{4\baselineskip}\begin{center}◇\end{center}\nopagebreak

오후에 돌아오니, 늘 야간
근무를 나가던 엄마가 소파에서
텔레비전을 보고 있었다.

``왔어?''

``응.''

텔레비전에서 예고가
흘러나왔다.

\begin{Verbatim}[samepage=true,frame=single,framesep=2.5mm,rulecolor=\color{gray},formatcom=\logfont]
[예고] 7월 16일(토) 정오
바당이, 바다에서 말하다
시민 과학 소년 고단우 군(13)
\end{Verbatim}

``나 열두 살인데.''

``한 살 올려 줬네, 방송국이.''
엄마가 리모컨을 내려놓았다.

그때 엄마 휴대폰이 짧게
떨었다.

\begin{Verbatim}[samepage=true,frame=single,framesep=2.5mm,rulecolor=\color{gray},formatcom=\logfont]
물마루 AI 연구소 | 보호자 안내
특집 생중계 출연 동의서를
보내 드립니다.
\end{Verbatim}

``출연한다는 건 뉴스로 먼저
알았는데.'' 엄마가 휴대폰을
엎어 놓았다. ``동의서는 이제
오네.''

``단우야. 너\ldots{} 저거 진짜
네가 한 거야?''

칠판 옆에 적은 제목이
떠올랐다. `내 라벨이 만든 첫
문장.'

``\ldots 반쯤.''

``반쯤?''

단우는 더 말하지 못했다.
엄마도 다그치지 않았다.
거실에는 냉장고 소리뿐.

엄마가 입을 열었다가 닫고,
다시 열었다.

``밥은?''

``\ldots 아직.''

엄마는 밥을 퍼 주고 맞은편에
앉았다. 단우가 숟가락을 들
때까지, 아무것도 묻지 않았다.

\Needspace{4\baselineskip}\begin{center}◇\end{center}\nopagebreak

그날 밤, 단우는 `옛날 것'
폴더를 열었다. test\_old.py
옆에는 봇이 라벨마다 한 줄씩
적어 둔 기록 파일이 있었다.
붙인 건 한밤중이어도, 줄마다
맨 앞엔 녹음된 때가 적혀
있었다.

`시험 볼 때 찍은 문제도,
맞으면 동그라미 받죠?'

첫날 기록부터 내려 읽었다.
말젯이가 사라지기 전에도,
`분리 위협'이 빠진 날은
없었다.

날마다 같은 번호만 찍는 애도,
가끔은 동그라미를 받는다. 봇도
날마다 찍고 있었던 걸까?

불턱에서 단우는 얼굴을 보고
나서야 숨비소리에 이름을
적었다. 봇은 바다를 한 번도
보지 않았다.

봇이 `분리 위협'을 가장 많이
붙인 그 순간, 새끼가 바로 옆에
있었다면? 그게 가장 센 대답일
것 같았다.

\begin{Verbatim}[samepage=true,frame=single,framesep=2.5mm,rulecolor=\color{gray},formatcom=\logfont]
봇 기록 | 녹음 6월 4일 오전
10분 동안  | 배 있음
분리 위협  | =========
\end{Verbatim}

``아홉 번.'' 단우가
중얼거렸다. ``십 분 사이에.''

`개체 사진' 6월 4일 칸은 비어
있었다. 조사선이 안 나간
날이었다. 방파제 카메라엔
출렁이는 물뿐.

그날 바다를 찍은 사람이 있어야
했다.

단우는 반 단톡방을 위로, 또
위로 올렸다. 6월 내내 채아가
올린 영상들 밑에, 고단우라는
이름은 없었다.

한 번도 제대로 읽지 않은 제목
하나에서 엄지가 멈췄다.

\begin{chatthread}
\chatroomtitle{곰새코지 6학년 수다방}
\chatbubble{채아}{말젯이 첫 점프!(6월 4일)\\이거 진짜 귀엽다!}
\end{chatthread}

``말젯이\ldots{} 6월 4일.''

튀어 오르는 장면만 이어 붙인
영상. 찍은 시각은 어디에도
없었다.

할머니 말이 떠올랐다. ``저
아이 영상은 꼭 중간을 먹어.''

그때는 영상이 뭘 먹느냐고
되물었다. 이제는 알 것 같았다.

잘리기 전 원본, 찍은 시각이
그대로 남은 파일을 가진 사람은
하나뿐이었다. 양채아. 지난주
금요일, 구멍가게 앞에서 `구독
취소' 하고 돌아선 애.

\dobira{5장}{원본 영상이 찾아낸 증거}{images/pagefill/ch09.png}{images/pagefill/ch09-page.png}

\section{휴지통 속
원본}\label{uxd734uxc9c0uxd1b5-uxc18d-uxc6d0uxbcf8}

\dateline{7월 9일 토요일 · 방송 특집까지 D-7 · 채아}

채아는 토요일 아침, 이불
속에서 그 메시지를 또 열었다.

\begin{chatthread}
\chatbubble{고단우}{6월 4일 원본 있어?\\촬영 시각 그대로.}
\end{chatthread}

``구독 취소.'' 하고 돌아선
뒤로 처음 온 메시지였다.
안녕도 없고, 미안도 없었다.

``뭐야. 숙제 검사야?''

채아는 답장 칸을 눌렀다가 도로
닫았다. 읽음 표시는 떴을 테니,
그걸로 충분했다.

과학 채널 오션톡의 알림도
울렸다.

\begin{Verbatim}[samepage=true,frame=single,framesep=2.5mm,rulecolor=\color{gray},formatcom=\logfont]
오션톡 | 새 영상
관광선이 새끼를 쫓아냈나?
불법 포획설까지
\end{Verbatim}

썸네일 구석에 활동가 윤나래의
글이 보였다. `예전에도 이
바다에서 돌고래를 잡아\ldots'
지난주 포구에서 윤나래가
휴대폰에 대고 하던 그
말이었다.

채아는 얼른 자기 채널로
넘어갔다.

\begin{Verbatim}[samepage=true,frame=single,framesep=2.5mm,rulecolor=\color{gray},formatcom=\logfont]
포구 시위 현장 (전체) | 댓글
ㅇㅇ: 전체라면서 왜 다 잘렸음?
고래고래: 선장 딸 어느 학교임?
\end{Verbatim}

``\ldots 시끄러운 거만 자른
거거든.'' 채아가 화면에 대고
쏘아붙였다.

둘째 줄 앞에서는 그 말도
나오지 않았다. 휴대폰을 베개
밑으로 밀어 넣자, 거기서 또
부르르 떨었다.

``아, 또 뭔데.''

이번엔 마을 단톡방이었다.

\begin{chatthread}
\chatroomtitle{곰새코지 마을방}
\chatbubble{구멍가게}{방파제 마이크가 끊겼다더라.}
\chatbubble{바당지기}{선장들이 잘랐다는 말도 있던데.}
\end{chatthread}

``선장들이\ldots.'' 채아는
소리 내어 읽다 말았다. ``왜?
배 소리 듣는 마이크라서?''

연구소가 방파제 밑 바닷속에
달아 둔 마이크였다.
선장`들'이라면 아빠도 그 안에
있었다.

창밖 포구에서는 아빠가 묶어 둔
배의 밧줄만 풀었다 감았다
했다.

``아빠 배는 잘못 없다는
증거.'' 채아가 창틀을
두드렸다. ``그런 장면 하나만
있으면 돼.''

\Needspace{4\baselineskip}\begin{center}◇\end{center}\nopagebreak

오후, 채아는 6월 원본 폴더에서
4일 오전 영상부터 열었다.

앞부분 40초는 지느러미 끝이
뜯긴 불쑥이 차지였다.

``불쑥이, 좀 비켜. 너 주인공
아니거든?''

불쑥이가 끝까지 안 비켜서,
채아는 빨리감기를 눌렀다.

등지느러미 뒤쪽이 반달처럼
떨어져 나간 어미, 반달이
솟았다. 옆구리에는 단우
할머니가 말젯이라 부르는
새끼가 붙어 있었다.

``그래, 이거지.''

화면 귀퉁이에는 아빠 배 뱃머리
끝이 걸려 있었다. 손님들에게
아빠가 늘 하던 말은
`돌고래한테는 오십 미터.'였다.
거리는 지킨 거였다. 눈대중으로
70미터쯤은 됐다.

``거봐. 칠십은 되잖아.''
채아가 등받이에 털썩 기댔다.
``이런 거 몇 개만 더 있으면
돼.''

아빠 배가 크게 나온 컷은
댓글이 무서워서 늘 휴지통으로
보냈다. 비우기는 한 번도
누르지 않았으니, 6월 컷은 전부
거기 있었다.

채아는 노트북 카메라에 대고
브이를 만들었다.

``안녕, 돌핀로그 채아예요---
오늘은 휴지통 탐험! 구독,
좋아요, 그리고---''

녹화 버튼도 안 누른 채였다.
채아는 브이를 슬그머니 접었다.

``\ldots 누구한테 하는 거야,
이걸.''

휴지통 목록을 내리던 손이
`0612\_1640.mp4'에서 섰다.

``이건 아빠한테 불리해서 버린
건데\ldots.''

이번에는 빨리감기 대신 확대
버튼을 눌렀다.

물보라 사이로 반달과 새끼가
나란히 떠올랐다. 새끼의
가슴지느러미 한쪽 뿌리에
초록색 그물 조각이 너덜너덜
걸쳐 있었다.

그리고 화면 아래로 불쑥, 아빠
배의 뱃머리.

이번엔 눈으로 잴 것도 없었다.
50미터 안, 20미터 남짓까지
들어와 있었다. 규정
위반이었다.

마우스를 쥔 손바닥이 축축했다.

``\ldots 근데 저 그물은 뭐야.
우리 배엔 그물 없어.
관광선이잖아.''

새끼는 그물을 매단 채 반달
옆에서 숨을 쉬고 다시
들어갔다.

``적어도 이날은 누가 잡아간 게
아니잖아. 엄마 옆에서 제
힘으로 헤엄치고 있잖아.''

아침에 본 썸네일의 `불법
포획설까지'가 생각났다. 아빠
배 탓을 덜려고, 윤나래의 그
말을 영상 끝에 남겨 둔 건
채아였다.

채아는 파일의 `정보'를 열었다.

\begin{Verbatim}[samepage=true,frame=single,framesep=2.5mm,rulecolor=\color{gray},formatcom=\logfont]
파일   0612_1640.mp4
촬영   2033-06-12 16:40:12
위치   곰새코지 앞바다
\end{Verbatim}

``찍은 날, 찍은 시각, 찍은
자리\ldots.'' 채아가 한 줄씩
손가락으로 짚었다.

옛날 사진 뒷면에 적어 두던
날짜와 장소 같은 거였다.

``촬영 시각 그대로.'' 채아가
단우의 메시지를 소리 내어
읽었다. ``아, 이거였구나. 잘라
붙이면 이게 다 떨어져
나가니까.''

\begin{Verbatim}[samepage=true,frame=single,framesep=2.5mm,rulecolor=\color{gray},formatcom=\logfont]
휴지통 | 0612_1640.mp4
되살리기 | 영구 삭제
\end{Verbatim}

커서가 `영구 삭제' 위로
미끄러졌다. 이걸 누르면 아빠
배가 너무 가까운 장면은
사라진다.

마우스를 쥔 검지가 한 번
들썩였다.

화면 속 새끼는 아직 떠 있었다.
몸에 초록색을 매단 채.

``이거 지우면\ldots{} 저
그물도 같이 없어지는 거네.''

채아는 커서를 옆 칸으로
옮겼다. `되살리기'.

그러고는 6월 4일 파일과 되살린
0612를, 원본 그대로 USB에 옮겨
담았다.

\Needspace{4\baselineskip}\begin{center}◇\end{center}\nopagebreak

일요일 해 질 녘, 채아는
노트북을 끌어안고 포구로
내려갔다. 아빠는 오늘도 밧줄을
감고 있었다.

``아빠.''

``\ldots 왜.''

배에 올라 화면을 아빠 쪽으로
들자, 버릇처럼 입꼬리가
올라갔다.

``안녕, 돌핀로그 채아예요---
오늘 보여 드릴 영상은요, 우리
아빠 배가---''

다음 말이 목에 걸렸다. 채아는
재생 버튼만 눌렀다.

``6월 12일. 아빠 배 나와. 나
이거 고단우한테 줄 거야.
연구소에 가져가라고.''

아빠는 갑판에 앉아 빨리감기
없이 보았다. 뱃머리가 불쑥
들어오는 대목은 되돌려 한 번
더 보았다. 처음부터 끝까지,
보통 빠르기로.

화면이 멈췄다. 아빠는 오래
말이 없었다.

``\ldots 아빠가 너무
들어갔다.''

``걔 통해서 연구소에 줘라.
구조팀에도.''

아빠가 화면 속 초록색을
손끝으로 짚었다.

``그 애기 먼저 찾아야지.''

``근데 아빠, 이거 나가면 댓글
또 달려. 이번엔 진짜 아빠
배니까.''

아빠가 휴대폰으로 연 건 채아의
채널이었다. 화면 위쪽에 `구독
중'이 떠 있었다.

``\ldots 언제부터?''

아빠가 `(전체)' 영상의 댓글
창을 열었다. 엄지는 `선장 딸
어느 학교임?' 앞에서 더
내려가지 않았다.

``그런 거 맨날 달려. 신경 안
써.''

``이거, 댓글 닫는 거 할 줄
아냐.''

``\ldots 응.''

``닫아라. 이 영상도 내리고.''

``아빠한테 좋게 자른 건데.''

``안다. 전체도 아니잖아.''

아빠는 썸네일 줄을 천천히
내렸다. 영상마다 채아가 브이를
하고 웃고 있었다.

``얼굴은 당분간 넣지 마라.''

``그럼 오프닝은?''

아빠가 일어섰다. 밧줄 끝을
기둥에 단단히 걸고 조타실 문을
잠그더니, 주말 내내 서성이던
배에서 내려왔다.

``가자.''

아빠가 채아 품에서 노트북을
받아 겨드랑이에 꼈다.

\Needspace{4\baselineskip}\begin{center}◇\end{center}\nopagebreak

그날 밤, 채아는 우산을 쓰고
단우네 현관 앞에 섰다. 비
내리는 골목 어귀에서 아빠
트럭이 와이퍼만 움직이며
기다렸다.

초인종을 누르자마자 문이
열렸다. 고단우. 눈 밑이
지난주보다 더 거뭇했다.

채아는 손바닥을 내밀었다.
손바닥 위에 USB 하나.

``6월 4일. 그리고 하나 더.''

``하나 더?'' 단우가 채아를
보았다. ``\ldots 확률적으로
네가 올 줄 알았어.''

``확률 같은 소리.''

``읽고 답도 안 했으면서.''

``답장 대신 온 거거든?''

단우가 안으로 들어가더니,
컵라면 하나를 들고 나왔다.

``넥타이. 세 개도 넘더라. 셀
수가 없어.''

소장 넥타이 내기의
컵라면이었다.

``거봐. 내가 이겼지.''

``세 개도 틀렸거든.''

``두 개보단 가깝거든.''

단우는 마루에 노트북을 펴고
USB를 꽂더니, 6월 4일 파일의
정보 카드부터 띄웠다.

``영상은 안 봐?''

``시각부터.''

단우가 숫자 몇 개를 노트에
옮겨 적고 끄덕였다.

``진짜 숙제 검사네.'' 채아가
0612를 가리켰다. ``하나 더는
이거야. 우리 아빠 배도 나와.
너무 가깝게.''

초록색 그물 조각이 나오자,
단우의 눈이 파일 이름에 가서
박혔다.

``6월 12일\ldots.''

``응. 그게 왜?''

``할머니가 새끼를 마지막으로
본 게 13일이야.'' 단우가
화면을 짚었다. ``이거, 그 바로
전날이야.''

채아의 팔에 오소소 소름이
돋았다.

``이거 박사님한테 지금 보내도
돼? 구조팀에도.''

``그러라고 가져온 거야.''

안쪽 방에서 나온 단우
할머니가, 말없이 화면 앞에
무릎을 꿇었다. 반달 옆에서,
초록색을 매단 작은 지느러미가
오르내렸다.

``말젯이다.''

할머니의 손가락이 화면 뒤쪽
검은 바위들로 옮겨 갔다.

``\ldots 저 뒤 바위,
그물먹는여주.''

이름처럼 그물이 잘 걸려서,
배들도 피해 다니는 바위였다.

빗소리가 현관 지붕을 두드렸다.

``근데 이거\ldots{} 한 달
전이잖아. 살아 있을까?''

아무도 대답하지 않았다.

\chapterartpage{images/pagefill/ch10-page.png}{images/pagefill/ch10.png}

\section{모르겠다고 말하는
법}\label{uxbaa8uxb974uxaca0uxb2e4uxace0-uxb9d0uxd558uxb294-uxbc95}

\dateline{7월 11일 월요일 · 방송 특집까지 D-5 · 다원}

다원은 마을 회의 때마다 같은
말을 제일 많이 들었다.

``그래서, 엄마는 뭐래?''

다원은 엄마 미숙의 말을 정리해
주는 아이였다. 월요일 저녁
마을회관, 엄마 옆자리도 다원
차지였다.

``오늘도 네가 옆에 있어 줘.''

엄마가 낮은 목소리로 속삭이며
다원 무릎을 쳤다.

어촌계와 주민이 함께 여는
회의였다. 진행하는 강덕수 계장
아저씨는 항구 공사 얘기를
한마디로 넘겼다.

``그건 나중에 공지하고. 자,
해녀들 생각부터 들어 보자.''

맨 안쪽엔 해녀회장님이 있었다.
계장 아저씨는 엄마를 보며
제주말을 쏟아 냈다.

``해녀들은 관광선 막자는 거냐?
그럼 마을은 뭘 먹고 살고?''

``뭐라고 하셔?''

불턱에선 다 알아듣던 엄마도,
회관에선 낯선 제주말이 빠르게
들리자 다원부터 보았다. 다원이
뜻을 풀어 주자 엄마는 천천히
대답했다.

``글쎄, 잘 모르겠어. 정말
모르겠어. 그건 회장님이 정할
일이야.''

눈이 다원에게 모였다. 예전
같으면 엄마 말을 어른들
입맛대로 다듬었을 거다.
`해녀들은 배가 걱정이래요.'
그러면 회의는 빨리 끝났다.
대신 며칠 뒤엔 `미숙이 배
싫다더라'가 포구를 돌았다.

다원은 노트 모서리를 접었다
폈다.

``모르겠대요.''

``\ldots 응?''

``엄마는 모르겠대요. 회장님이
정할 일이래요.''

``다원이 왜 그렇게 오래
설명하냐?''

``길게 모르겠다고 한 거예요.''

여기저기 웃음이 새어 나왔다.
계장 아저씨가 손을 내저었다.

``해녀들 배 싫어하는 거, 다
아는 거 아니냐. 그렇게 옮기면
되주.''

``엄마가 한 말은 그게
다예요.''

``그게 다면 설명이 왜 그렇게
길어?''

또 웃음이 터졌다. 다원은
노트를 덮지 않았다.

``제가 보태면, 그건 제
말이죠.''

맨 안쪽에서 회장님 목소리가
났다.

``해녀들 일은 해녀들끼리
정한다. 다음.''

엄마 손이 다원 손등을 꾹
눌렀다.

\Needspace{4\baselineskip}\begin{center}◇\end{center}\nopagebreak

다음은 6학년 채아와 고단우
차례였다.

단우가 노트북을 텔레비전에
꽂았다.

``이게 6월 12일이에요.''
채아가 영상을 틀었다. 초록
그물 걸친 새끼 곁에 관광선
뱃머리가 붙어 있었다.

``우리 아빠 배예요.
20미터쯤이에요.'' 채아
목소리가 떨렸다. ``연구소랑
구조팀엔 어제 먼저 보냈어요.''

맨 뒷줄 채아 아빠가 고개를
끄덕였다. 어른들 목소리가
높아졌다.

``그럼 배 때문에 다친 거네!''

``그물은 어선 거주! 관광선에
그물이 어디 있어!''

``그래서 그 새끼, 지금은 살아
있냐?''

단우가 고개를 저었다.

``\ldots 몰라요. 6월 13일
뒤로는 본 사람이 없어요.''

계장 아저씨가 팔짱을 꼈다.

``다 좋은데, 결국 그 AI가 맞힌
거 아니냐? 새끼가
없어졌잖아.''

``시험도 찍어서 맞으민
동그라미주.''

``학생, 숙제는 숙제고. 맞힌 건
맞힌 거다.''

귀가 빨개진 단우가 자판을
두드렸다. 타닥, 탁.

화면이 반으로 갈라졌다.

\begin{Verbatim}[samepage=true,frame=single,framesep=2.5mm,rulecolor=\color{gray},formatcom=\logfont]
화면 왼쪽   | 원본 영상
촬영        | 6월 4일 10:12
화면 오른쪽 | 라벨 기록
녹음된 때   | 6월 4일 10:12부터 10분
분리 위협   | =========
\end{Verbatim}

왼쪽 영상에선 반달 옆에 바짝
붙은 말젯이가 통 튀어 올랐다.

``자르지 않은 원본이라, 촬영
시각이 그대로 남아 있어요.''

``저 배도 우리 아빠 배예요.
그날은 70미터쯤. 거리는
지켰어요.''

단우가 막대를 짚었다.

``이건 같은 시각에 녹음된
소리에 붙은 라벨이에요. 십 분
동안 `분리 위협'이 아홉 번.''

왼쪽에서 말젯이가 또
뛰어올랐다. 그 순간 오른쪽
막대 끝 `분리 위협' 한 칸이
깜박였다.

``푸흡. 떨어지긴 뭘 떨어져. 딱
붙어서 노는구만.''

앞줄 어른이 무릎을 치자 웃음이
뒤로 번졌다.

``이 규칙은 배가 가까이 오면
`분리 위협'을 찍어요. 그날도
배가 왔을 뿐이에요. 날마다
같은 번호를 찍으면 언젠가
하루는 맞아요. 근데 찍어서
맞힌 걸 `알았다'고 텔레비전에
내보내면 안 되잖아요.''

회관이 조용해졌다. 계장 아저씨
눈이 작은 몸을 따라갔다.

``\ldots 뛰네, 뛰어.''

계장 아저씨 입에서 나온
소리였다. 다원은 노트에 단우
발표는 괜찮았고 화면 맞춤법은
두 군데 틀렸다고 적었다.
마지막엔 `제 규칙' 대신 `이
규칙'이라고 한 이유도 써
두었다.

``엄마, 봤어?''

엄마는 이번엔 다원 말고 화면
속 작은 등만 보았다.

``작은 돌고래, 봤어요.''

다원이 대신 정리한 말이
아니었다. 미숙이 직접 한
한국어였다. 일어선 미숙이 손을
들고 있었다.

``금요일 새벽, 그물먹는여
안쪽. 혼자. 몸에 뭐 걸려
있었어요. 초록색. 그래도
헤엄쳤어요.''

채아가 반쯤 일어났다.

``살아 있다는 거예요?
금요일에요?''

``네. 숨, 쉬었어요.''

뒤에서 누가 고개를 저었다.

``잘못 봤주. 거긴 배도 못
들어가는 덴데.''

``봤어요. 물질 가는 길에.''

엄마가 탁자에 손가락으로 바위
몇 개를 그리고 그 사이로
물길을 그었다.

``물 빠져도 여기는 바다랑
이어져요. 작은 고기도
들어와요.''

손가락이 물길 안쪽을 톡톡
두드렸다.

``태풍 오면 그물이 여기
쌓여요. 더 많이.''

다원은 더 설명하지 않았다.
미숙이 이미 자기 목소리로
말하고 있었다. 노트를 가만히
덮었다.

``거긴 배가 못 들어가요. 우리
눈엔 안 띄지.''

어부 시절부터 이 바다를 훤히
안다는 오삼촌이, 뒷벽에서
팔짱을 풀었다. 자리에 앉은
엄마가 노트에 손을 얹었다.
아까보다 뜨거운 손이었다.

\Needspace{4\baselineskip}\begin{center}◇\end{center}\nopagebreak

회의 뒤, 다원은 방파제 입구
갯바위에 앉았다. 단우가 옆에
와 앉았다.

``\ldots 나는? 발표 어땠어?''

``괜찮았어. 맞춤법 두 개
빼고.''

``어디?''

``안 알려 줘. 그건 설명료 받고
해 주는 거야.''

``설명료 얼마인데?''

``컵라면. 큰 거.''

단우가 피식 웃었다.

``아까 그냥 해녀들은 배
싫어한다고 해 줘도 됐잖아.
다들 그렇게 아는데.''

다원은 운동화 끝으로 따개비를
건드렸다.

``아는 척하면, 엄마가 다쳐.''

단우는 더 묻지 않았다.

``너는 아까 왜 `이 규칙'이라고
했어?''

``\ldots 뭐?''

``라벨은 네 아이디로 붙었다며.
그럼 그 규칙은 누가 만든
건데?''

단우는 대답 대신 가방끈만
감았다 풀었다.

``\ldots`돌려줘'는 누가 한
말일까?''

말 돌리는 줄 알면서 다원은
받아 주었다.

``아직 몰라. 누가, 무엇을
돌려달라는 건지부터
알아야지.''

오삼촌 트럭에서 라디오가
울렸다.

``\ldots 제7호 태풍은 목요일
밤 제주를 지날 것으로
보입니다\ldots.''

단우가 벌떡 고개를 들었다.
다원은 탁자 위를 톡톡 두드리던
엄마 손가락을 떠올렸다.

그물먹는여. 그물. 더 많이.

\dobira{6장}{정정문에 담긴 고백}{images/pagefill/ch11.png}{images/pagefill/ch11-page.png}

\section{정정합니다}\label{uxc815uxc815uxd569uxb2c8uxb2e4}

\dateline{7월 14일 목요일 · 방송 특집까지 D-2}

목요일 밤, 곰새코지의 불이
한꺼번에 나갔다. 냉장고 소리가
뚝 끊기자, 할머니가 밥상
가운데 초를 세웠다.

제주시 병원에 있는 엄마는
태풍에 발이 묶였다. 문자는 한
줄, 문 잘 잠그라고.

촛불 아래서 단우는 화요일에 온
탁 박사 메시지를 다시 열었다.

\begin{chatthread}
\chatbubble{탁 박사}{채아 영상은 구조팀에도 전달했어.\\소장님은 방송 날 생중계로 찾아보자고 하셔.\\미숙 씨가 본 건 지난주라 지금 위치는 모른대.\\위치를 모르면 구조팀도 출동할 수 없대.\\헤엄칠 수 있으면 회의가 먼저.\\걸려서 못 움직이면 바로 간대.}
\end{chatthread}

단우는 마지막 줄을 두 번
읽었다. 미숙이 본 새끼는
헤엄치고 숨 쉬었으니, 지금은
회의가 먼저였다.

그런데 태풍이 왔다. 회관
탁자를 톡톡 두드리던 미숙의
손가락이 떠올랐다.

``태풍 오면 그물이 여기
쌓여요. 더 많이.''

``할머니, 태풍 지나가면
그물먹는여에 그물이 많이
쌓여요?''

``물때 봐야 알주.''

휴대폰이 부르르 떨었다.

\dateline{탁 박사: 정전. 연료 0잔. 위험.}

단우가 풉 웃자 촛불이 누웠다
일어섰다.

``휴대폰 불 꺼라. 그것도 느
연료여.''

``제 연료는 배터리예요,
할머니.''

``그러니 아끼라고.''

단우가 화면을 끄자, 방이 촛불
크기만큼 줄어들었다.

할머니는 올해가 마지막
물질이라던 주황색 테왁 끈을
매만졌다.

``할머니는 왜 올해 그만둬요?
숨은 아직 저보다 오래
참잖아요.''

``바당 밑에서 숨을 세다 보민,
올라올 숫자가 모자라.''

할머니가 엄마 문자를
들여다보며 말했다.

``느 어멍이 열다섯 살에 처음
들어갔다 나와서 무섭다 했주.''

``엄마가요? 바다에요?''

``나는 `무섭긴 뭐가 무섭나' 그
한마디만 했다.''

할머니는 더 말하지 않고,
촛불을 단우 쪽으로 슥 밀었다.

``\ldots 할머니. 나도 무서운
거 있어요.''

``말해라.''

``텔레비전에 나온 거, 제가 한
거 아니에요. 근데 그 말을
하면, 숨긴 것도 다 말해야
돼요. 말로는 잘 안 나와요.
써서 읽어 드려도 돼요?''

할머니가 고개를 끄덕였다.

\Needspace{4\baselineskip}\begin{center}◇\end{center}\nopagebreak

단우는 노트북을 열어, `옛날
것' 폴더 맨 아래
test\_old.py에 원래 이름을
다시 쳤다.

``labelbot.py.''

첫 줄 `정정합니다.'는 금방
썼다. 처음 쓴 둘째 줄은 `제
봇이 라벨 붙이는 걸 조금
도왔습니다.'였다.

``조금은 무슨.''

단우는 그 줄을 통째로 지웠다.

소장실에서 들은 말이
끼어들었다. `어떻게 맞혔든,
맞은 건 맞은 거예요.' 그
말대로라면 봇도, 금색 왕관도
괜찮았다.

바탕화면 구석의 대회 출품 파일
`내 라벨이 만든 첫 문장'은
열지 않았다.

자정 무렵 글이 끝났지만,
와이파이는 텅 비어 있었다.
공유기도 정전이었다.

``아, 맞다.''

단우는 글을 휴대폰으로
다원에게 보냈다.

\begin{chatthread}
\chatbubble{다원}{띄어쓰기 틀린 데 열두 군데.\\내용은 안 고쳐도 돼.\\근데 이제 알겠다. 왜 '이 규칙'이었는지.}
\chatbubble{단우}{그땐 내 거라고 못 했어.}
\chatbubble{다원}{이제 했네.}
\end{chatthread}

``맞춤법 검사야, 뭐야.''

열두 군데를 고치는 동안
배터리가 한 칸 줄었다.

\Needspace{4\baselineskip}\begin{center}◇\end{center}\nopagebreak

단우는 노트북을 들고 할머니
앞에 앉았다.

``어려운 말이 많아요. 맨 위
`정정'은, 틀린 걸 바로잡는다는
말이고요.''

``읽어 봐라. 다 못 알아들어도
들으마.''

\begin{Verbatim}[samepage=true,frame=single,framesep=2.5mm,rulecolor=\color{gray},formatcom=\logfont]
정정합니다

생중계 자막 '돌려줘.'는
잘못 나간 말입니다.
돌고래가 그렇게 말했다는
증거는 없습니다.
그 아래 찍힌 제 아이디로 붙은 라벨도
제가 듣고 붙인 것이 아닙니다.
\end{Verbatim}

할머니는 테왁 끈을 쥔 채
가만히 있었다.

\begin{Verbatim}[samepage=true,frame=single,framesep=2.5mm,rulecolor=\color{gray},formatcom=\logfont]
제가 만든 봇 프로그램이
제가 자는 동안 라벨을 붙였습니다.
배가 가깝고 휘슬이 많으면
무조건 '분리 위협'이었습니다.
배가 온 것을
새끼가 떨어진 것으로 잘못 붙인 겁니다.
연구소 AI의 마지막 칸이
그 '분리 위협'을 '돌려줘'로 바꿨습니다.
\end{Verbatim}

``봇이 그거냐. 그 프로그램.''

``네. 제 대신 라벨 붙이는
거요. 제가 만들었어요.''

``느 대신.''

``\ldots 네.''

할머니가 턱으로 화면을
가리켰다. 계속하라는 뜻이었다.

\begin{Verbatim}[samepage=true,frame=single,framesep=2.5mm,rulecolor=\color{gray},formatcom=\logfont]
새끼가 없어졌으니
결국 맞은 게 아니냐는 말도 들었습니다.
하지만 6월 4일 오전에는
새끼가 어미 바로 옆에서 뛰고 있었습니다.
그 시각에 녹음된 십 분짜리 소리에도
제 봇은 '분리 위협'을 아홉 번 붙였습니다.
어떻게 맞혔든 맞으면 된다면,
제 봇이 한 일도 괜찮은 게 됩니다.
\end{Verbatim}

단우는 불턱의 해녀들처럼 숨을
길게 들이마셨다.

\begin{Verbatim}[samepage=true,frame=single,framesep=2.5mm,rulecolor=\color{gray},formatcom=\logfont]
그런데 제 봇은 바다를 한 번도 보지 않았습니다.
가입할 때 저는 이 칸에 체크했습니다.
'모든 라벨은 제가 듣고 붙이겠습니다.'
하지만 봇이 붙인 소리를 저도 듣지 않았습니다.
그래서 고칩니다.
제가 자는 동안 붙은 라벨 9,308개를
모두 빼 주세요.
봇은 원래 이름 labelbot.py로 함께 올립니다.
\end{Verbatim}

할머니의 손이 테왁 끈을 꽉
쥐었다.

\begin{Verbatim}[samepage=true,frame=single,framesep=2.5mm,rulecolor=\color{gray},formatcom=\logfont]
아직 모르는 것
'돌려줘'가 된 그 휘슬은
이름 같은 소리일 수 있습니다.
하지만 누구 이름인지 모릅니다.
이름이 맞는지도 모릅니다.
새끼가 지금 어디 있는지 모릅니다.

고단우 (라벨러 danu_0101)
\end{Verbatim}

``이거 올리는 게 무서워요.
다들 알게 되는 거.''

할머니는 한참 촛불만 보다가
입을 열었다.

``무섭주.''

그러고는 화면 맨 아래 줄을
손가락으로 짚었다.

``느 이름 쓴 거냐.''

``네.''

``게민 됐다.''

할머니가 거친 손바닥으로
단우의 머리를 한 번, 세게
쓸었다. 화면 속 글자가 번져
보여서, 단우는 소매로 눈가를
문질렀다.

할머니는 다시 테왁 끈을
감았다. 한 바퀴, 또 한 바퀴.

``할머니, 연료 좀 쓸게요.''

할머니는 끈에서 눈을 떼지
않았다.

``써라.''

엄마는 신호가 한참 가서야
받았다. 뒤에서 기계 소리가
삑삑 울렸다.

``단우야? 무슨 일 있어?''

``나 글 하나 썼어. 저번에
엄마가 진짜 내가 한 거냐고
물었잖아. 그때 `반쯤'이라고 한
거, 나머지 반.''

전화 너머가 조용해졌다.

``보내 봐.''

엄마가 읽는 동안 누가 엄마를
부르는 소리가 났다. 엄마는
대답하지 않았다.

``다 읽었어. 이번엔 반쯤
아니네. 올리기 전에 탁
박사님한테 먼저 보내. 그리고
그 출연 동의서. 방송에 네 얼굴
내는 거엔, 엄마 서명 안 했어.
앞으로도 안 해.''

``\ldots 그럼 방송은?''

``그건 탁 박사님이랑 엄마가
얘기할게. 글은 네 거니까 네가
정해. 얼굴은\ldots{} 엄마가
지켜.''

누가 또 엄마를 불렀다. 끊기
전에 엄마가 한 번 더 말했다.

``문 잠갔지?''

``응.''

단우는 탁 박사에게 정정문을
보냈다. 보내기 버튼 위에서
엄지는 망설이지 않았다.

\Needspace{4\baselineskip}\begin{center}◇\end{center}\nopagebreak

새벽 2시, 휴대폰이 울렸다.
초는 손가락 한 마디만큼 남아
있었다.

``안 자고 있을 줄 알았어.''

``박사님도요.''

``글 읽었어. 연구소 공지로
올리자. 이번엔 내가 올릴게.''

``제 이름도요?''

``공개할 땐 가려. `시민 라벨러
A'로. 그 대신 `제가'는 그대로
둬. 네 글이니까.''

``아이디는 생중계에 벌써
나왔잖아요. 이름도요.''

``알아. 나온 건 못 주워 담아.
그래도 우리가 먼저 퍼뜨리진
않아. 이름이 퍼지면 사람들은
글보다 너를 봐.''

``소장님은요?''

``보냈어. 봇 얘기는 소장님도
이 글로 처음 아셨어. 내가 먼저
말씀드렸어야 했는데. 아직 답이
없으셔.''

공지 맨 아래엔 새로운 원칙이
붙었다. ``공개할 때는 미성년자
이름을 가린다.''

처음 연구소에 간 날엔 세 번째
줄에서 눈을 피했던 단우가,
이번에는 끝까지 읽었다.

``이제 자. 올라가는 건 아침에
봐.''

단우는 초를 불어 끄고 밥상
옆에 엎드렸다. 공지 화면을
새로 고치다 눈이 감겼다.

웅, 냉장고가 다시 돌아가는
소리에 눈을 떴다. 창밖은
희뿌옜고 바람도 잦아들었다.

공지는 금요일 새벽 5시 40분에
올라와 있었다. 맨 먼저 댓글을
단 사람은 탁 박사였다.

\begin{Verbatim}[samepage=true,frame=single,framesep=2.5mm,rulecolor=\color{gray},formatcom=\logfont]
탁은재: '돌려줘'를 만든 마지막 칸,
탁은재: 그 책임은 저에게도 있습니다.
탁은재: 반대만 하고 막지 않았습니다.
\end{Verbatim}

탁 박사의 이름은 댓글 앞에
그대로 있었다. 가려진 건 글
끝의 이름 하나, 단우
것뿐이었다.

버릇처럼 `귀 기울임'
리더보드를 열자, 맨 위에
\texttt{1위\ 새벽세시감귤}이
보였다. 단우는 목록을 끝까지
내렸다.

``진짜 없네.''

어깨 힘이 쭉 빠졌다.

연구소 실시간 화면에서 방파제
마이크 칸은 회색이었고, 큰여
마이크엔 태풍이 남긴 잡음만
지글거렸다.

휘슬 표시가 깜박이는 건 등대
마이크 칸뿐이었다. 노트
귀퉁이의 메모가 떠올랐다.

\dateline{등대 쪽만? 나중에}

그 `나중'이 지금인지도 몰랐다.
깜박임이 점점 빨라졌다. 단우가
본 중에 가장 빨랐다.

등대 마이크에서 W-07이 1분에
열한 번.

그러다 뚝 끊겼다.

\chapterartpage{images/pagefill/ch12-page.png}{images/pagefill/ch12.png}

\section{등대 쪽
휘슬}\label{uxb4f1uxb300-uxcabd-uxd718uxc2ac}

\dateline{7월 15일 금요일 · 방송 특집까지 D-1}

태풍으로 학교가 쉰 금요일
오전, 깜빡 잠든 단우를 오삼촌
전화가 깨웠다.

``등대 마이크에 그 휘슬 다시
잡혔다. W-07. 모양도
그대로다.''

``진짜요? 어디요?''

``그건 모르지.'' 오삼촌이
하품을 삼켰다. ``오후에 방파제
마이크 건지러 간다. 엄마
허락은 받았다. 조끼 입고,
방파제 밑까지만.''

오후, 물마루 1호가 방파제
밑에서 엔진을 끄자 물소리가
커졌다.

``엔진 끄면 다 들려.''

단우는 아재 개그를 기다렸지만,
오삼촌은 줄만 당겼다.

``삼촌, 개그는요?''

``오늘은 쉰다.''

줄 끝에 올라온 방파제 마이크
전선이 초록색 덩어리에 칭칭
감겨 있었다.

``\ldots 그물이에요?''

``다 해진 그물이다. 봐라,
칼자국 하나 없지. 파도에
끌려다니다 전선을 당긴 거다.
자른 게 아니라 감긴 거야.''

`선장들이 잘랐다는 말도 있고.'
지난주 마을방의 그 `선장들'
안엔 채아 아빠도 있었다.

단우는 6월 12일 영상을 띄운
휴대폰을 그물 옆에 대 보았다.

``삼촌, 영상 속 그물이랑
똑같아요.''

오삼촌이 등대 아래 바위 쪽으로
턱짓했다.

``이 물이 다 그물먹는여로 감아
들어. 그물도, 뭐든. 너 마을방
있지? 사진 올려라. 내 휴대폰은
화면이 어두워서 안 보인다.''

오삼촌이 작은 상자를 발로
밀었다.

``그리고 이거. 휴대용
마이크다. 박사님이 네
프로그램에 꽂아 보래.''

전선을 찍어 올리자 답이
달렸다.

\begin{chatthread}
\chatroomtitle{곰새코지 마을방}
\chatbubble{바당지기}{그물이었네. 선장들한테 미안.}
\chatbubble{구멍가게}{선장들한테 미안.\\채아 아빠한테도 미안.}
\end{chatthread}

``뭐라고들 하냐?''

``다들 미안하대요. 다 똑같은
말로요.''

``삼촌, 등대 마이크가
그물먹는여랑 제일 가깝죠?''

``그래. 조용한 날엔 그쪽
소리도 넘어온다.''

W-07은 어떤 날 등대 마이크
하나에만 잡혔다. 그런 날
휘슬은 세 마이크 삼각형 바깥,
그물먹는여 쪽에 있었을지도
몰랐다. W-07이 크게 울릴 때면
사진엔 거의 늘 반달이 있었다.

``반달이 그물먹는여 안쪽을
드나들었던 거예요. 다원이
엄마가 작은 돌고래를 봤다는
데도 거기고요.''

``드나들었을 수도 있지.''
오삼촌이 `수도'에 힘을 줬다.

``\ldots 새끼가 거기 있다는
증거는 아니죠.''

``아니지. 그래도 그 휘슬, 오늘
아침에도 그쪽에서 울렸다.''

단우는 노트의
\texttt{등대\ 쪽만?\ 나중에}
아래에 밑줄을 쳤다. 물음표는
지우지 않았다. 오삼촌이 쉰다던
개그를 꺼냈다.

``단우야, 물고기 반대말이 뭔지
아냐?''

``불고기요.''

``\ldots 아는구나.''

멀리서 불쑥이가 물속으로
사라졌다.

``관객이 하나 줄었네.''
오삼촌이 열쇠를 돌렸다.

\Needspace{4\baselineskip}\begin{center}◇\end{center}\nopagebreak

저녁, 단우는 휴대용 마이크를
꽂고 listen.py를 열었다.

``이제 이걸 귀로 만들어야지.''

W-07만 모으던 필터를, 지금
들리는 휘슬 이름을 바로 띄우는
감지기로 바꿀 참이었다.
라디오에서 노래 제목이 뜨듯,
소리 이름을 화면에 띄우게 하는
일이었다.

``맞는 상자가 있으면 이름.
없으면\ldots.''

labelbot.py에 쳤던 마지막 두
줄이 생각났다.

\begin{Verbatim}[samepage=true,frame=single,framesep=2.5mm,rulecolor=\color{gray},formatcom=\logfont]
아니면
      라벨 = "이동"
\end{Verbatim}

봇은 모르는 소리에도 아무
이름이나 붙였다. 휴대폰은
모르는 번호엔 이름 대신 숫자만
띄운다.

``없으면, 모른다고.''

\begin{Verbatim}[samepage=true,frame=single,framesep=2.5mm,rulecolor=\color{gray},formatcom=\logfont]
listen.py | 시험
W-07 녹음     | W-07
딱총새우 녹음 | 미분류-모름
              | 잡음 가능성 80%
\end{Verbatim}

`딱총새우 잡음이겠지' 하고
접어 두었던 따닥따닥 소리였다.
그 줄을 지우려다 말았다.

``이번엔 안 접어.''

손가락이 먼저
\texttt{listen\_live\_최종}까지
쳤다.

``\ldots 버그.''

원래 이름으로 저장하자마자
초인종이 울렸다. 머그잔 든 탁
박사였다.

``연료 여덟 잔째. 어머님이
조건 세 개 말씀하셨어.''

그때 엄마 전화가 왔다.

``배는 타도 돼. 박사님이랑
오삼촌이 옆에 있는대. 구명조끼
벗으면 바로 데려오시라고 했어.
네 얼굴이랑 이름은 화면에 안
나오게 해 달라고 했고.''

엄마는 금방 끊었다.

``원래는 배 여러 척이 돌고래
가까이 가면 안 돼. 그런데
구조팀이 물마루 1호랑 채아네
배를 `도와주는 배'로 허락받아
줬어.''

탁 박사가 단우 눈을 똑바로
보았다.

``대신 시키는 것 말고는
아무것도 건드리지 않기.''

``네.''

``너는 물마루 1호에서 헤드폰을
쓰고, 채아는 아빠 배에서
연구소 마이크로 또 하나의 귀가
돼. 다원이는 바닷가에서
해녀분들이랑 있고.''

``카메라는요?''

``너는 안 비춰. 그렇게
약속받았어. 귀만 빌려줘.''

내일 단우 자리는 카메라가
비추지 않는 갑판 구석이었다.
그게 별로 아깝지 않은 게
단우는 조금 이상했다.

``어르신, 내일 배에 같이 타
주실 수 있을까요? 바위는
어르신이 제일 잘
아신다면서요.''

망사리를 깁던 할머니가
바늘에서 눈을 떼지 않았다.

``배는 안 탄다.''

스크루에 다쳤다던 할머니 물벗
이야기가 떠올라 단우는 말을
보태지 않았다.

``무전기나 하나 얻어 와라.
갯바위에서 듣주.''

\Needspace{4\baselineskip}\begin{center}◇\end{center}\nopagebreak

밤 10시, 마이크 상자를 안은
단우와 할머니가 가로등 아래
섰다.

``여기서 기다리마.''

가로등 옆 승합차에서 방송팀이
상자를 배로 나르고 있었다.
오삼촌이 하품하며 손을
내밀었다.

``마이크 이리 줘라. 오늘 실어
둬야지.''

``삼촌, 저건 다 뭐예요?''

``방송 쪽 짐이다. 큰 스피커도
있다더라.''

한 상자 옆면에 흰 스티커가
보였다.

\dateline{재생 장비}

``삼촌, 이거 소리 트는
거예요?''

오삼촌이 스티커를
들여다보았다. 또 하품을 하려던
입이 그대로 닫혔다.

``나야 모르지. 실으라니 싣는
거다.''

오삼촌은 상자를 안고 내려갔다.
연구소 액자의 첫 줄이
떠올랐다.

\mbox{\texttt{1. 녹음한 소리를 바다에 틀지 않는다}}

단우는 탁 박사에게 메시지를
썼다.

\begin{chatthread}
\chatbubble{danu\_0101}{재생 장비가 실렸어요.\\원칙 1번은요?}
\end{chatthread}

두 시간 전, 현관에서 한 약속이
떠올랐다.

`시키는 것 말고는 아무것도
건드리지 않기.'

단우는 둘째 줄을 지웠다가 다시
썼다. 묻는 것도 건드리는 걸까.

보내기 버튼 위에서 엄지가
멈췄다. 어젯밤 정정문을 보낼
때는 망설이지 않던 엄지였다.
승합차 뒷문이 다시 열리고,
상자 하나가 더 물마루 1호로
넘어갔다.

\dobira{7장}{엔진을 끄자, 소리가 들렸다}{images/pagefill/ch13.png}{images/pagefill/ch13-page.png}

\section{세 개의 귀가 찾은
곳}\label{uxc138-uxac1cuxc758-uxadc0uxac00-uxcc3euxc740-uxacf3}

\dateline{7월 16일 토요일 · D-day · 06:30}

태풍이 남긴 구름에 포구는 아직
잿빛이었다. 어젯밤 가로등
아래서 멈췄던 엄지는, 결국
이불 속에서 보내기를 눌렀다.

\begin{chatthread}
\chatbubble{danu\_0101 (7월 16일 새벽 00:47)}{재생 장비가 실렸어요.\\원칙 1번은요?}
\chatbubble{탁 박사 (7월 16일 새벽 01:12)}{잘 보냈어. 판단은 내가 해.}
\end{chatthread}

틀겠다는 말도, 안 틀겠다는
말도 없었다.

``삼촌, 그거 보이긴 해요?''

오삼촌이 하품을 하며 휴대폰을
보고 있었다. 화면이 어두워
옆에서는 글자 하나 안 보였다.

``나는 보여.''

갑판 뒤쪽 상자들은 방수포를
덮어쓰고 있었다. 탁 박사가
머그잔을 들고 배에 올랐다.

``박사님, 그 답장이요.''

``말 그대로야. 넌 귀만
빌려줘.'' 탁 박사가 휴대폰을
내밀었다. ``그리고 이거,
소장님이 새벽에.''

\begin{chatthread}
\chatbubble{봉 소장}{정정문 봤어요. 깜짝 놀랐어요.\\방송은 예정대로 진행할게요.}
\end{chatthread}

소장 말끝마다 붙던
`미안해요'는 없었다.

7시, 물마루 1호가 포구를
나섰다. 마이크가 빠진 방파제
앞에는 채아네 배가 대신 나와
있었다.

``물마루 1호, 여기 곰새코지
3호. 자리 잡았습니다.''

곧바로 채아가 끼어들었다.

``아빠가 마이크 물에 넣었어!
방금 휘슬, 우리도 잡혔어!''

7시 반, 태블릿의 등대, 큰여,
채아네 배 마이크 칸에 W-07이
차례로 떴다.

``같은 휘슬인데 등대에 제일
먼저 닿았어. 큰여는 조금 늦게,
채아네 배는 더 늦게. 왜일까?''

``멀수록 늦게 오니까요? 번개
보고 천둥 칠 때까지 세는 거랑
같네요.''

``그거야. 셋이 얼마나 늦었는지
견줘 보면, 소리 난 자리가
나와.''

\begin{quote}
\textbf{단우의 메모 \#4} 도착
시간차 = 같은 소리가
마이크마다 늦게 닿는 차이.
늦게 닿을수록 멀다.
\end{quote}

탁 박사가 지도를 띄웠다.

\begin{Verbatim}[samepage=true,frame=single,framesep=2.5mm,rulecolor=\color{gray},formatcom=\logfont]
바다 지도 | 마이크 자리
*방파제 앞(채아네 배) . . *큰여
    .                   .
        .           .
            *등대
               (그물먹는여)
\end{Verbatim}

삼각형 바깥, 그물먹는여 둘레가
통째로 노랗게 칠해졌다.

``이게 다예요? 엄청
넓은데요.''

``운동장 몇 개는 되겠다.''
오삼촌이 말했다.

``마이크 셋 바깥이라 퍼진
거야. 이 넓이로는 구조팀이 못
들어와.''

\Needspace{4\baselineskip}\begin{center}◇\end{center}\nopagebreak

9시, 물마루 1호는 그물먹는여
바깥에 멈춰 섰다. 오삼촌이
닻을 내리고, 단우의 휴대용
마이크를 뱃전 너머로 내렸다.

``줄 끝까지 내렸다. 이제 네
귀다.''

단우는 헤드폰을 쓰고
listen.py를 켰다.

``물마루 1호, 여기 갯바위요.
엄마 바꿔 줄게요. 저 오늘은
심부름 휴가예요.''

다원이었다. 말끝에 쩝쩝, 귤
씹는 소리가 붙어 있었다.

``여기 물 색깔 달라요. 등대
보고 왼쪽, 바위 많은 데요.''

이번엔 다원 엄마, 미숙이었다.
곧 첨벙, 미숙이 테왁을 안고
물에 드는 소리가 났다.

탁 박사가 점을 하나 찍었다.
``미숙 씨 말대로면 이쯤이야.
그래도 넓네.''

그때 단우 화면에 W-07이 떴다.
등대 마이크보다 한 박자
빨랐다.

``박사님, 제 마이크가 등대보다
먼저예요.''

``그럼 우리 배 쪽이야. 먼 쪽
반은 지워도 돼.''

노란 칠의 먼 쪽 절반이
지워졌다. 남은 쪽은 미숙이
말한 점 쪽이었다. 그때 다원이
외쳤다.

``단우 바꿔 주세요! 급한
거예요.''

탁 박사가 으쓱하며 무전기를
넘겼다.

``단우야, 내기할래? 지금
올라오는 사람 누군지 맞히면 귤
하나.''

``그게 급한 거야?''

``응. 너 우리 엄마 숨비소리
그림으로 그렸다며.''

파도 소리가 커졌다. 다원이
무전기를 바다 쪽으로 내민
모양이었다.

호오이---

머릿속에 그림이 먼저 그려졌다.
평평하게 가다가 올라가고, 두
번 내려가는 선,
\mbox{F-U-D-D.}

``정\ldots.''

단우는 입을 다물었다. 불턱에서
정임에게 한 말이 떠올랐다. `저
혼자 보는 거예요.'

``다원아, 정임 삼춘한테 먼저
물어봐 줘. 소리로 이름 불러도
되냐고.''

``\ldots 방금 정답 말했잖아.''

멀리서 정임의 대답이 돌아왔다.

``내 소리? 해라, 해. 테레비
나오는 거냐?''

``무전이라고 전해 줘.''

``나 오늘 심부름 휴가라니까.''
다원이 투덜대더니 소리쳤다.
``삼춘! 테레비 아니고
무전이래요!''

``에이, 아깝다.''

미숙도 끼어들었다. ``나도
돼요.''

``들었지? 귤 하나 줄게.
다음.''

호오이---

이번엔 조금 가까웠다. 단우는
눈을 감았다. 평평, 올라감,
내려감, 다시 평평,
\mbox{F-U-D-F.}

``\ldots 다원이 엄마예요.''

무전기 너머가 한 박자
조용했다. 할머니도 아무 말이
없었다.

``\ldots 맞았어. 귤 두 개.''

불턱의 그 아침에도 그랬다.
미숙이 올라왔을 때만, 할머니
입에서 이름이 나오지 않았다.

``너 귀신이냐?''

오삼촌이 단우를 돌아보았다.

``허락받고 녹음한 사람만요.''

멀리서 미숙이 외쳤다. ``단우,
내 소리 아직 안 들려줬어!''

``오늘 끝나면 꼭요!''

``할머니, 노란 칠이 그물먹는여
안쪽까지 줄었어요. 등대 보고
왼쪽이요.''

잠깐 지직거리더니, 오늘
처음으로 할머니 목소리가 났다.

``태풍 뒤엔 물이 소등굴로 감아
들어. 거기다.''

``소등굴이요?'' 단우가
되물었다.

``그물먹는여 안쪽 바위틈 굴.
물 색 다르다는 데서 조금 더
들어간 자리.''

탁 박사가 미숙의 점 옆에 작은
동그라미를 그렸다. ``노란 칠
안이야. 딱 맞네.'' 운동장 몇
개만 하던 자리가 동그라미
하나로 줄었다.

``이 정도면 교실 한 칸이네.''
오삼촌이 말했다.

``만조가 한 시 반이여.''
할머니가 덧붙였다.

``만조요?''

``물이 제일 많이 차는 때. 물이
차면 바위틈이 다 잠겨. 새끼가
혹시 어디 걸려 있으면, 그때가
제일 위험해.''

``물때는 바다 버스 시간표다.''
오삼촌이 말했다.

탁 박사가 구조팀 수의사에게
전화를 걸어 스피커를 켰다.

``12시 40분에 도착해요. 자리가
계속 바뀌면 헤엄치는 거니까,
가서 지켜볼게요.''

12시 40분이면 만조까지 한
시간도 안 남는다. 태풍 밤 탁
박사의 답이 떠올랐다. 헤엄칠
수 있으면 먼저 지켜보고, 걸려
움직이지 못하면 바로
구조하기로 했다.

자리가 안 바뀌면, 그건 둘째
줄이었다.

\Needspace{4\baselineskip}\begin{center}◇\end{center}\nopagebreak

11시 반, 물 위로 회색 등
하나가 떠올랐다. 생중계는
정오였다.

``반달이다.''

오삼촌이 목소리를 낮췄다.
반달은 바위 가장자리만
맴돌았다.

``새끼는요?''

``안 보여.''

채아네 배는 그물먹는여
건너편에 있었다. 헤드폰 속
W-07 위로 웅웅 소리가 얇게
깔렸다.

``삼촌, 엔진 아직 켜져
있어요?''

``프로펠러는 세웠다. 규정대로.
엔진까지 끄라는 말은 없어.''

무전기에 낯선 목소리가
끼어들었다.

``여기 방송 쾌속정. 지금
출발합니다. 발전기 때문에
엔진은 켜 두고 갑니다.''

``켜 두고 오겠다고?'' 오삼촌이
무전기를 노려보았다.

그때 봉 소장의 전화가 왔다.
스피커를 켜자 태블릿에
스튜디오 화면이 열렸다. 금색
돌고래 넥타이, 그 옆에는
`심사위원장' 이름표를 단 사람.

``탁 박사님, 미안한데요. 거기
그림이 너무 조용해요.''

작은 칸에는 뱃머리 카메라가
잡은 잔잔한 바다뿐이었다.

``반달은 찾았어요. 새끼는
아직이고요.''

``재생 장비 실었죠? 엄마 소리
들으면 새끼가 나오지 않겠어요?
반달 소리 틀어서 불러내요.
그림이 필요해요.''

바람이 방수포 귀퉁이를
들추었다. 그 밑으로 흰
스티커가 보였다.

\dateline{재생 장비}

연구소 액자의 첫 줄이 다시
떠올랐다.

\mbox{\texttt{1. 녹음한 소리를 바다에 틀지 않는다}}

탁 박사는 휴대폰을 든 채
대답하지 않았다. 태블릿 구석의
시계가 11:50으로 바뀌었다.

\chapterartpage{images/pagefill/ch14-page.png}{images/pagefill/ch14.png}

\section{엔진을 끈
순간}\label{uxc5d4uxc9c4uxc744-uxb048-uxc21cuxac04}

\dateline{7월 16일 토요일 · D-day · 11:50}

오삼촌도 단우도 탁 박사만
보았다. 스피커 너머에서 소장이
한 번 더 말했다.

``탁 박사님, 엄마 소리 들으면
새끼가 나오지 않겠어요? 지금은
그게 제일 빠른 길 같아서요.''

탁 박사가 잠깐 눈을 감았다.
휴대폰을 쥔 손가락 마디가
하얬다.

``소장님, 안 틀어요.''

``\ldots 네?''

``연구소 원칙 1번이에요.
녹음한 소리를 바다에 틀지
않는다.'' 탁 박사가 눈을 떴다.
``엄마는 바로 저기 있어요.
진짜 소리가 벌써 바다에
있는데, 녹음을 틀면 새끼를
오히려 쫓아낼 수도 있어요. 저
소리가 반달한테 무슨 뜻인지도
우린 몰라요.''

``그래도 생방송인데\ldots.''

``알아요. 그래서 제가
책임질게요.''

스피커 너머가 잠깐 조용했다.

``\ldots 그럼 그 학생이라도
한마디만. 사람들이 같이 기다려
주게요.''

``그건 더 안 돼요. 소장님,
이따 다시 걸게요.''

전화가 끊기자 단우는 참았던
숨을 내쉬었다. 새벽의 답장이
떠올랐다.

`잘 보냈어. 판단은 내가 해.'

\Needspace{4\baselineskip}\begin{center}◇\end{center}\nopagebreak

12시 정각, 태블릿 화면이
반으로 갈렸다. 왼쪽은 밝은
스튜디오, 오른쪽은 바다였다.

``토요일 정오, 곰새코지 돌고래
특집 생방송입니다!''

인사가 끝나기도 전에, 배 한
척이 하얀 물살을 일으키며
달려왔다. 오삼촌이 무전기를
낚아챘다.

``방송 배, 여기 300미터
안이에요. 허락받은 배도
아니고요. 프로펠러 세우세요!''

``\ldots 세웁니다.''

방송 쾌속정이 그물먹는여
건너편에 섰다. 카메라는 반달이
맴도는 바다를 향했다. 단우가
앉은 조타실 뒤는 반대편이었다.

``이곳은 곰새코지
앞바다입니다!''

현장 진행자의 목소리가
쩌렁쩌렁 울렸고, 단우는
헤드폰을 두 손으로 눌렀다.
프로펠러는 다 섰는데도 헤드폰
속은 웅웅거렸다. 헛도는 엔진에
발전기 소리까지 얹혔다. 그
밑에서 W-07이 들렸다 말았다
했다.

``켜 두고 오겠다더니, 진짜
그대로네.'' 오삼촌이 혀를
찼다.

단우의 눈이 엔진 열쇠로 갔다.
손을 뻗으면 닿았다.

`시키는 것 말고는 아무것도
건드리지 않기.'

단우는 열쇠 대신 탁 박사를
돌아보았다.

``박사님.'' 단우가 헤드폰을
한쪽만 벗었다. ``프로펠러
말고, 엔진까지요.''

``엔진?''

``프로펠러를 세워도 엔진은
켜져 있잖아요. 신호에 선 차도
부릉거리는 것처럼요.
발전기도요. 그게 W-07을
덮어요.''

\begin{quote}
\textbf{단우의 메모 \#5} 배
소리 = 시끄러운 급식실. 친구가
내 이름을 불러도 묻힌다.
\end{quote}

탁 박사가 대답하기도 전에
오삼촌이 엔진 열쇠에 손을
얹었다.

``엔진 끄면 다 들려.''

연구소에서 예의상 ``하하''
웃고 넘긴 말이었다. 이번엔
하나도 웃기지 않았다. 처음부터
맞는 말이었으니까.

12시 3분, 오삼촌이 열쇠를
돌렸다. 발밑의 떨림이
가라앉았다. 단우는 벗었던
헤드폰 한쪽을 다시 귀에 댔다.

탁 박사가 단우를 한 번 보고
무전기를 들었다.

``현장 책임자 탁은재예요. 모든
배, 엔진까지 꺼 주세요.''

``곰새코지 3호, 닻 내리고
끕니다.'' 채아 아빠의 목소리
뒤로 채아가 외쳤다. ``아빠,
빨리!''

``여기 방송 배. 생방송이라
발전기가\ldots.''

``지금 꺼야 들려요. 그것도
제가 책임져요.''

대답 대신, 쾌속정 쪽 웅웅
소리가 뚝 멎었다.

바다에 파도 소리만 남았다.
헤드폰 속에서는 무언가
따닥따닥 튀었다.

그리고 그 밑에서, W-07이
끊기지 않고 끝까지 들렸다.
올라갔다 내려가고, 다시
올라갔다가 길게. 단우는
헤드폰을 누른 손을 저도 모르게
풀었다.

``들려요.'' 단우가 속삭였다.

태블릿 속에서도 현장 진행자가
카메라에 바짝 다가가 속삭였다.

``시청자 여러분, 지금\ldots{}
아무 소리도 안 들립니다.''

``현장? 현장, 들리십니까?''
스튜디오 진행자가 귀를 짚었다.
그 옆에서 소장이 진행자의 팔을
잡고, 입술에 손가락을 댔다.

꼬르륵. 오삼촌의 배에서 난
소리였다.

``\ldots 조용하니까
배고프네.''

``쉿!''

단우와 탁 박사가 한꺼번에
말했다. 오삼촌이 두 손으로
배를 꾹 눌렀다. 노트북 화면에
줄이 한 칸씩 내려왔다.

\begin{Verbatim}[samepage=true,frame=single,framesep=2.5mm,rulecolor=\color{gray},formatcom=\logfont]
listen.py | 물마루 1호
시각  배 소리     휘슬
12:03 ==========  W-07 ?
12:05 =====       W-07
12:06 .           W-07 U-D-U-F
12:06 .           미분류-모름 +1
                  잡음 가능성 62%
\end{Verbatim}

W-07 줄 아래, 처음 보는 줄이
하나 붙어 있었다.

``박사님, 이거요.''

오삼촌이 화면을 들여다보았다.

``딱총새우 잡음이겠지. 바위
밑에 천지다.''

구멍가게 평상에서 단우가
속으로 한 말 그대로였다.
그래도 단우는 그 줄을 접지
않았다.

``모르는 거예요.''

단우는 눈을 감았다. 파도 소리
사이로 아주 가는 소리가 하나
있었다. 조금 올라갔다가,
평평하게. 딱총새우의
따닥따닥과는 달랐다. W-07이
울리면, 그 소리가 꼭 뒤따랐다.

봇이었으면 아무 이름이나
붙이고 넘어갔을 소리였다.
감지기는 모른다고 적어 두었다.

무전기가 지직거렸다.

``물마루 1호! 여기 곰새코지
3호! 우리 헤드폰에도 가는 소리
들려요! 그 휘슬 바로
다음에요!''

채아였다. 탁 박사가 짧게
답했다.

``곰새코지 3호, 잘 잡았어.
마이크 그대로 둬.''

``같은 `모름' 소리인데, 배 두
척이랑 등대에 닿은 시간이 다
달라.''

``번개랑 천둥이요.''

탁 박사가 그 차이로 지도에 선
두 개를 그었다.

``이번엔 배 두 척이 그물먹는여
양쪽에 있어서 좁게 나와.
바위에 튕기니까 몇십 미터는
틀릴 수 있지만.''

두 선이 할머니의 동그라미 안쪽
끝에서 겹쳤다. 단우가 먼저
읽었다.

``소등굴 입구 쪽이에요.''

오삼촌이 머리를 긁적였다.
``\ldots 새우는 아니겠다.''

12시 8분, 소리가 또 울렸다. 두
선은 아까와 똑같은 자리에서
만났다.

``안 움직여요.'' 단우의
목소리가 떨렸다. ``아까랑 같은
자리예요.''

탁 박사가 스피커를 켠 채 다시
수의사에게 전화를 걸었다.
``소등굴 입구 쪽이에요. 두 번
다 같은 자리고요.''

``몇 분째 자리가 안 바뀌면
바위에 걸렸을 수 있어요.
걸려서 못 움직이면 바로 가요.
지켜보는 거 말고, 지금 가요.
속도 올릴게요.''

수의사가 조금 낮게 덧붙였다.

``아니면 다행이고요.''

갯바위에서 할머니 무전이
끼어들었다.

``만조까지 한 시간 이십 분. 물
벌써 들어온다.''

단우는 `미분류-모름' 줄만 모아
보았다.

\begin{Verbatim}[samepage=true,frame=single,framesep=2.5mm,rulecolor=\color{gray},formatcom=\logfont]
12:06:40 미분류-모름 U-F
12:07:30 미분류-모름 U-F
12:08:40 미분류-모름 U-F
12:10:10 미분류-모름 U-F
\end{Verbatim}

소리와 소리 사이가 조금씩
벌어지고 있었다. 점점 드물게
울린다는 뜻이었다. 단우는 다음
줄을 기다리며 속으로 수를
셌다. 하나, 둘, 셋\ldots.
헤드폰을 누른 손에 힘이
들어갔다.

\dobira{8장}{소등굴, 마지막 구조 신호}{images/pagefill/ch15.png}{images/pagefill/ch15-page.png}

\section{소등굴 구조
신호}\label{uxc18cuxb4f1uxad74-uxad6cuxc870-uxc2e0uxd638}

\dateline{7월 16일 토요일 · D-day · 12:10}

12시 11분이 되어도 다음 줄은
내려오지 않았다. 그 1분이 거의
다 가서야 한 줄이 붙었다.

\begin{Verbatim}[samepage=true,frame=single,framesep=2.5mm,rulecolor=\color{gray},formatcom=\logfont]
12:11:50 미분류-모름 U-F
\end{Verbatim}

``점점 늦어져요.''

``알아. 구조팀이 속도
올린댔어. 조금만 더.''

12시 25분, 하얀 물살 하나가
곧장 달려왔다.

``구조팀이다. 15분이나 일찍
왔네.''

구조팀 배가 옆에 섰다. 주황색
옷을 입은 사람이 무전기를
들었다. 아까 전화 너머에서
``지금 가요.'' 하던
목소리였다.

``구조팀 수의사 문지혜예요.
지금부터 현장은 제가
맡을게요.''

``넘겨 드릴게요.'' 탁 박사가
답했다.

``잠수사 둘이 입구를 찾을
거예요. 바위 안으로는 배가 못
들어가요. 바위를 아는 분
계세요? 배 위에서 짚어 주실
분이요.''

무전기가 조용해졌다. 어제저녁
할머니 대답이 떠올랐다.

`배는 안 탄다.'

그때 갯바위 쪽 무전이
지직거렸다.

``나다. 그물먹는여 바위는 내가
안다.''

``할머니, 배\ldots.''

``안다. 곰새코지 3호, 등대 밑
방파제 끝으로 와라.''

``곰새코지 3호, 갑니다.'' 채아
아빠 목소리였다.

오삼촌이 단우에게 쌍안경을
건넸다. 방파제 끝에 할머니가
주황색 테왁을 옆구리에 끼고 서
있었다.

``방파제 끝에 댔습니다. 천천히
오르세요.''

할머니가 난간을 쥐었는데,
쌍안경 속에서도 그 손이 떨리는
게 보였다. 배 스크루에 다친
물벗 이야기가 떠올랐다.

할머니는 난간을 쥔 채 한동안
서 있었다. 채아가 갑판을
건너와 할머니 팔을 붙잡았다.
할머니는 채아를 한 번
내려다보고, 한 발씩 배에
올랐다.

``어르신이\ldots{} 진짜
타셨네.'' 오삼촌이 낮게
말했다.

곰새코지 3호가 구조팀 배
옆으로 와 섰다. 물마루 1호와는
공을 던지면 닿을 거리라,
할머니의 하얀 입술이
맨눈으로도 보였다.

할머니는 테왁 끈에 긴 줄을
이어 맸다. 난간을 쥘 때와
달리, 매듭을 짓는 손은 빨랐다.
그러고는 테왁을 물에
내려놓았다.

``물이 소등굴로 감아 든다
했지. 봐라.''

테왁이 물살을 타고 바위 사이로
끌려 들어갔다. 바위 몇 개를
지나자 할머니가 줄을 꽉
쥐었다. 테왁이 그 자리에서
동동 떴다.

``거기다.''

``저 테왁 밑이 입구예요?''
문지혜가 물었다.

``그래.''

``잠수사 들어갑니다. 곰새코지
3호, 엔진 끄고 대기해
주세요.''

``끕니다.''

채아 아빠가 열쇠를 돌리자,
웅웅거리던 소리가 멎었다.
할머니가 그제야 난간에서 한
손을 뗐다.

잠수사 둘이 물로 뛰어들었다.
곧 하나가 고개를 내밀었다.

``입구는 어떻게 생겼어요?''

``입구 왼쪽에 소 등짝 같은
바위, 그 밑이 좁아.''

테왁 왼쪽에 둥그스름한 바위
등이 반쯤 나와 있었다.

``소 등짝\ldots{} 소등굴.''
단우가 입속으로 중얼거렸다.

``안쪽은요?'' 잠수사가 다시
물었다.

할머니가 줄을 쥔 채 잠깐 말이
없었다.

``거기까진 나도 모른다.''

무엇을 물어도 ``물때 봐야
알주.''라던 할머니였다. 그
입에서 `모른다'가 나온 건
처음이었다.

``충분해요. 거기서부터는
저희가 볼게요.''

잠수사 둘이 테왁 밑으로
사라졌다. 채아는 할머니 팔을
놓지 않았다.

\Needspace{4\baselineskip}\begin{center}◇\end{center}\nopagebreak

오후 1시 5분, 문지혜의 무전이
다시 울렸다.

``찾았어요. 가슴지느러미에
걸친 그물 조각이 틈에 쌓인 헌
그물이랑 엉켰어요. 아침 썰물
때 걸린 것 같아요. 물이
차오르면서 숨 쉬러 올라오기가
힘들어졌어요.''

소 등짝 바위는 파도가 칠
때마다 다 잠겼다 드러났다.
문지혜가 잠수사들에게
지시했다.

``바위에 엉킨 줄부터 끊어요.
그다음 물 위로 받쳐서 숨부터
쉬게 해요.''

잠시 뒤, 문지혜의 목소리가
높아졌다.

``줄 끊었어요. 올려요!''

테왁 옆에서 물거품이 한꺼번에
끓어올랐다. 작은 회색 등이 물
위로 솟았고, 잠수사 둘이
양옆에서 받치고 있었다.

푸우.

숨구멍에서 물보라가
솟아올랐다.

``숨 쉰다!'' 오삼촌이 저도
모르게 주먹을 쥐었다. 탁
박사는 머그잔을 든 채 입술을
깨물었다.

``\ldots 됐어요. 그물 다
걷었어요.'' 문지혜의
무전이었다.

잠수사가 초록색 그물 조각을
들어 올렸다. 6월 12일 영상 속
그 초록색이었다.

받치던 손이 떨어지자, 작은
등이 바위틈을 빠져나갔다. 바깥
물에서 지느러미 뒤쪽이
반달처럼 떨어져 나간 등이
나타났고, 작은 등이 그 옆에
바짝 붙었다.

``옳지, 옳지.'' 오삼촌이 제
무릎을 쳤다.

무전기에서 채아 목소리가 났다.

``말젯이다. 말젯이, 엄마 옆에
붙었어.''

비 오던 현관에서는
`이거'라고만 하던 채아였다.

W-07이 뜸해지더니, 1분에 두
번으로 줄었다. 태풍 지난
새벽에는 1분에 열한 번이던
소리였다.

``박사님, 울리는 게\ldots{}
6월 초랑 비슷해요.''

``비슷한 거고.'' 탁 박사가
말했다.

W-07 뒤를 따라, `모름' 소리가
한 줄 내려왔다.

\begin{Verbatim}[samepage=true,frame=single,framesep=2.5mm,rulecolor=\color{gray},formatcom=\logfont]
13:12:40 미분류-모름 U-F
\end{Verbatim}

``저거, 새끼 소리 아니냐?''
오삼촌이 물었다.

``그럴 수도요.''

단우는 그 줄을 고치지 않았다.

``쓸린 상처가 많고 말랐어요.''
문지혜였다. ``오늘은
여기까지예요. 붙잡고 있으면 더
위험해요. 앞으로는 지켜봐야
알아요.''

``만조 전에 됐다.''
할머니였다.

그때 방송 쾌속정 옆에서 커다란
회색 등이 불쑥 솟았다.
지느러미 끝이 뜯긴 불쑥이였다.
불쑥이가 숨을 뿜자 물보라가
카메라를 덮쳤고, 태블릿 속
현장 화면이 물방울투성이가
됐다.

``아, 카메라가\ldots!'' 현장
진행자의 목소리가 뒤집혔다.

``다 젖었네.'' 오삼촌이
키득거렸다.

곰새코지 3호 갑판에는 할머니가
털썩 주저앉아 어깨를 떨고
있었다. 단우가 소리쳤다.

``할머니, 괜찮아요?''

``배 냄새, 버그다.''

``할머니, 그거 쟤
말버릇이잖아요!'' 채아가
웃음을 터뜨렸다.

단우는 바람막이를 벗어 뭉쳤다.

``받아!''

채아가 바람막이를 받아 할머니
어깨에 덮었다.

닦아 낸 태블릿 화면 아래로
자막이 떴다.

\textbf{돌고래가 돌려달라던
새끼, 마침내 돌아오다} 바당이,
바다에서 말하다 \textbar{}
특집 생중계

``\ldots 돌려달라던.'' 단우가
중얼거렸다.

탁 박사가 빈 머그잔을
내려다보았다. ``정정할 게 하나
더 생겼네.''

\Needspace{4\baselineskip}\begin{center}◇\end{center}\nopagebreak

밤 9시, 마을방도 반 단톡방도
말젯이 이야기로 시끄러웠다. 그
사이로 첫 생중계의 `돌려줘.'
자막만 오려 낸 영상이 돌았다.

\begin{Verbatim}[samepage=true,frame=single,framesep=2.5mm,rulecolor=\color{gray},formatcom=\logfont]
돌려줘.
근거 녹음 조각 #A-3318
라벨러 danu_0101
\end{Verbatim}

\begin{Verbatim}[samepage=true,frame=single,framesep=2.5mm,rulecolor=\color{gray},formatcom=\logfont]
바다덕후77: 거짓말한 초딩 라벨러
ㅇㅇ: 어느 학교 애냐
익명: danu_0101 누군지 아는 사람?
\end{Verbatim}

연구소 공지 끝에는 여전히
`시민 라벨러 A'뿐이었다.

``\ldots 이름은 가려 줬는데.''

하지만 오려 낸 영상 속
아이디는 아무도 가려 줄 수
없었다. 귀가 뜨거운데도 엄지는
화면을 계속 밀어 올렸다.

``단우야.''

간호사 가방을 멘 엄마가 방문
앞에 서 있었다.

``엄마, 이거\ldots.''

``봤어. 병원에서.'' 엄마가
손을 내밀었다. ``이리 줘
볼래?''

엄마는 휴대폰을 뺏지 않았다.
받아서 책상 위에 엎어 놓았을
뿐이었다. 그러고는 자기
휴대폰을 들고 베란다로 나갔다.

``아이 아이디 나온 영상, 내려
주세요.''

방송국에 거는 전화였다. 끊고
나서는 영상마다 신고 창을
열었다.

돌아온 엄마가 엎어 둔 휴대폰을
손바닥으로 한 번 눌렀다.

``오늘은 그만 봐.'' 엄마가
잠깐 말을 골랐다.
``\ldots 밥은?''

``\ldots 아직.''

``엄마도 아직.''

라면 물 끓는 소리를 따라
단우도 부엌으로 나갔다. 식탁
위에, 이번에는 그릇이 두
개였다.

\Needspace{4\baselineskip}\begin{center}◇\end{center}\nopagebreak

밤 11시 40분, 단우는 노트북을
열었다. 한쪽 창에서는
listen.py가 등대 마이크 소리를
받고 있었다.

메신저 창에도 답이 와 있었다.

\begin{chatthread}
\chatbubble{탁 박사}{바당이에서 라벨 9,308개 다 뺐어.\\오늘 커피는 몇 잔째인지도 모르겠네.}
\chatbubble{채아}{6월 4일 영상 넣어. 아빠도 괜찮대.\\근데 너, 돌고래 실제로 봤냐?}
\end{chatthread}

단우는 피식 웃다가 맨 아래
줄에서 멈췄다. 답은 짧게 쳤다.

\dateline{단우: 응. 봤어.}

이번에는 `데이터로는'을 붙이지
않았다. 마감은 밤 11시 59분.
출품 창 제목 칸에는 그 제목이
그대로 있었다.

\dateline{내 라벨이 만든 첫 문장}

단우는 제목을 지우고, 입속으로
먼저 말해 보았다.

``찍어서 맞은 거랑, 들어서 안
거.''

단우는 그 말을 줄여 제목 칸에
치고, 파일을 끌어다 놓았다.

\begin{Verbatim}[samepage=true,frame=single,framesep=2.5mm,rulecolor=\color{gray},formatcom=\logfont]
푸른코드 청소년 AI 대회 | 출품
제목  맞은 것과 안 것
1  정정합니다
2  labelbot.py
3  listen.py
4  6월 4일 기록과 영상
5  내 라벨(봇 9,308개 뺌)
\end{Verbatim}

두 파일을 나란히 열자, 마지막
줄이 보였다.

\begin{Verbatim}[samepage=true,frame=single,framesep=2.5mm,rulecolor=\color{gray},formatcom=\logfont]
labelbot.py | 아니면 라벨 = "이동"
listen.py   | 없으면 미분류-모름
\end{Verbatim}

둘 다 단우가 친 줄이었다.

밤 11시 58분, 단우는 `제출'을
눌렀다. 막대가 끝까지
차올랐다.

\mbox{\texttt{23:58:41 업로드 완료}}

책상 위 휴대폰이 부르르
떨었지만, 단우는 뻗던 손을
거두었다. 휴대폰은 엎어진
그대로였다.

listen.py 창에 줄이 하나
내려왔다.

\begin{Verbatim}[samepage=true,frame=single,framesep=2.5mm,rulecolor=\color{gray},formatcom=\logfont]
listen.py | 등대 마이크
23:58:30 W-07 U-D-U-F
23:59:02 미분류-모름 +1
\end{Verbatim}

\chapterartpage{images/pagefill/ch16-page.png}{images/pagefill/ch16.png}

\section{잘라 낸
장면들}\label{uxc798uxb77c-uxb0b8-uxc7a5uxba74uxb4e4}

\dateline{7월 23일 토요일 · 생중계 후 7일 · 채아}

채아는 지난 엿새 동안 온 뉴스
알림을 하나도 지우지 않았다.
언덕길에서 알림들을 다시
보았다.

\begin{Verbatim}[samepage=true,frame=single,framesep=2.5mm,rulecolor=\color{gray},formatcom=\logfont]
[7/19] 마을 약속: 물질 땐 관광선 멀리
[7/20] 윤나래, 포획설 글 정정
[7/20] 오션톡 "충격 정정!" 200만 회
[7/21] 대회 결과 발표: 9월 1일
\end{Verbatim}

윤나래 정정 글은 ``제가
틀렸습니다.''로 시작했다. 처음
글 첫 줄 `예전에도 이 바다에서
돌고래를 잡아\ldots'엔 줄이
그어져 있었다. 그 말을 영상
끝에 남겨 둔 건 채아였다.

오션톡 줄에선 입이 삐죽
나왔다.

``정정 영상이 내 백만 조회를
넘었네.''

휴대폰이 또 떨었다.

\dateline{고단우: 엄마가 오늘은 집에 있으래.}

``알았어. 넌 집에 있어.''
채아가 화면에 대고 말했다.

\Needspace{4\baselineskip}\begin{center}◇\end{center}\nopagebreak

오전 10시, 연구소 마당은
기자회견 카메라로 빼곡했다.

채아는 맨 뒷줄에서 휴대폰을
단상으로만 들었다. 아빠 말대로
얼굴은 빼고, 마이크와 단상만.
생방송을 켜자 입꼬리가
올라갔다.

``안녕, 돌핀---''

채아는 입을 다물었다 다시
열었다.

``오늘은 안 잘라요. 다시보기도
그대로 둘게요.''

\dateline{물멍: 오프닝 어디 감?}

탁 박사가 종이를 쥐고 단상에
올랐다. 봉 소장을 당겨 보다가
채아는 눈을 깜박였다. 넥타이
자리에 단추 푼 셔츠
깃뿐이었다.

``넥타이 없으니까 딴 사람
같다.''

\dateline{물멍: ㅋㅋ 다 들림}

채아는 입술을 다물었다. 탁
박사가 종이를 펼쳤다.

```돌려줘'는 번역이
아니었습니다.''

마당이 조용해졌다.

``그 소리는 이름일 가능성이
있습니다. 누구의 이름인지는
모릅니다. `돌려달라던' 자막도
틀렸습니다. 구조는
`돌려줘'라는 문장 덕분이
아니었습니다.''

탁 박사가 눈을 들었다.

``문장을 지우고 들었기
때문입니다.''

기자 하나가 손을 들었다.

``문장을 지웠다는 게 무슨
뜻입니까?''

``모르는 소리는 `모른다'고
적어 두고 들었습니다. 그날은
배들도 엔진까지 다 껐습니다.''

``시민 데이터 쪽 문제
아닙니까?''

``아닙니다. 라벨 하나를
문장으로 바꿔 발표한 건
연구소입니다.''

``그럼 이번엔 누가
들었습니까?''

탁 박사가 마이크를 당겼다.
``그건 말씀드리지
않겠습니다.''

채아는 휴대폰 쥔 손에 힘을
줬다.

``\ldots 잘했다.''

``처음 자막의 확신도
31\%는요?''

``31\%는 맞을 확률이
아니었습니다. 닮은꼴
점수였습니다.''

``연예인 닮은꼴 앱. 31\%
닮았대도 그 사람은 아니잖아.''
채아가 중얼거렸다.

소장이 마이크 앞에 섰다. 손이
버릇처럼 목께로 갔다가 잡을 게
없어 내려왔다.

``제가 서둘렀습니다.''

그 한 줄뿐이었다. 탁 박사가
종이를 넘겼다.

``오늘부터 바당이의 문장
만들기는 끕니다. 프로젝트
이름도 `이름 지도'로
바꿉니다.''

``그럼 누구 이름입니까?''

``모릅니다.''

``추측이라도요.''

``모릅니다.''

기자들이 일제히 손을 들었다.

``팔 아파도 안 잘라요.''

언덕을 내려오는데 고단우가
전화했다.

``봤어. 끝까지.''

``넥타이 얘기도 들었냐.''

``다 들림.''

``\ldots 시끄러.''

\Needspace{4\baselineskip}\begin{center}◇\end{center}\nopagebreak

그날 저녁, 아빠가 식탁에서
봉미숙을 뜯었다. `과태료'가
벌금 같은 거란 건 채아도
알았다. 그 아래 날짜, 6월
12일. 휴지통에서 되살린 영상
속 그날이었다.

아빠는 종이를 접지 않았다.
휴대폰 밑 빨간 도장 종이를
꺼내 새 종이 옆에 나란히
놓았다. 그러고는 숟가락을
들었다.

등 뒤 텔레비전에서 뉴스가
나왔다.

``물마루 AI 연구소는 실수를
바로잡고, 연구를 이어 가기로
했습니다.''

채아는 식탁을 한 장 찍었다.
업로드는 않고 숟가락을 들었다.

\Needspace{4\baselineskip}\begin{center}◇\end{center}\nopagebreak

7월 29일 금요일, 채아는 탈의장
앞에서 회장님을 기다렸다.

``회장님, 시위하던 날 찍힌
거요. 영상에 넣어도 돼요?''

회장님은 망사리를 내려놓지도
않았다.

``난 싫다.''

``네. 뺄게요.''

그날 밤, 채아는 제목 칸부터
채웠다.

``제목은 뭐냐.'' 소파에서
아빠가 물었다.

``잘라 낸 장면들. 버렸던 거
다시 붙이는 거.''

\begin{Verbatim}[samepage=true,frame=single,framesep=2.5mm,rulecolor=\color{gray},formatcom=\logfont]
잘라 낸 장면들 | 편집
6월 12일 아빠 배 20m  | ?
포구 시위 고함        | 넣음(얼굴 흐림)
"단체 생각 아니에요"  | 넣음(얼굴 흐림)
회장님이 본 장면      | 뺌
\end{Verbatim}

채아는 6월 12일 컷을 틀었다.
초록 그물을 매단 새끼 아래로
아빠 배 뱃머리가 불쑥
들어왔다.

``아빠, 이거 진짜 넣어? 그
종이 날짜인데.''

아빠는 뱃머리를 한참 보았다.

``\ldots 그 종이, 내 영상
때문에 온 거지?''

``배는 내가 몰았다.''

아빠가 턱으로 화면을 가리켰다.

``올려라. 내 배 맞다.''

``사람들이 또 아빠 욕할
텐데.''

``종이 받았잖아. 이제 숨길
것도 없다.''

채아는 물음표를 지우고
`넣음'을 눌렀다.

``얼굴은.'' 아빠가 물었다.

``안 나와. 날짜는 적고, 찍은
자리는 지웠어.''

업로드하자 첫 댓글이 달렸다.

\dateline{새벽세시감귤: 원본 고마워요}

``새벽 세 시는 아직
멀었는데.'' 처음 보는
닉네임이었다.

댓글이 줄줄이 달렸다.

\begin{Verbatim}[samepage=true,frame=single,framesep=2.5mm,rulecolor=\color{gray},formatcom=\logfont]
ㅇㅇ: 선장 아빠 결국 걸렸네
익명: 그 라벨러 초딩도 이 동네 애임?
익명: 이름 아는 사람 댓글 ㄱㄱ
\end{Verbatim}

채아는 아래 두 줄을 지웠다.

``이건 안 돼.''

첫 줄 앞에선 손가락이 멈췄다.
아빠가 의자를 끌어와 앉았다.

``그건 둬라. 맞는 말이다.''

채아는 그 줄을 두고 맨 위에
댓글을 고정했다.

\begin{Verbatim}[samepage=true,frame=single,framesep=2.5mm,rulecolor=\color{gray},formatcom=\logfont]
채아의 돌핀로그 | 고정 댓글
아이 이름, 아이디 언급 댓글은 지웁니다.
\end{Verbatim}

아빠가 그 줄을 읽고 채아를
봤다. ``너도 들어가냐, 거기.''

``\ldots 응. 나도 아이잖아.''

아빠가 고개를 끄덕였다. 그사이
구독자 수가 몇 번 바뀌었다.

\begin{Verbatim}[samepage=true,frame=single,framesep=2.5mm,rulecolor=\color{gray},formatcom=\logfont]
채아의 돌핀로그 | 구독자
업로드 전   3만
밤 11시     2만 8천
밤 11시 반  2만 9천5백
\end{Verbatim}

``롤러코스터 타냐.'' 채아가
중얼거렸다.

``나는 아직 구독 중이다.''
아빠가 픽 웃었다.

그때 셋만 있는 단톡방이
울렸다.

\begin{chatthread}
\chatbubble{고단우}{31일이 할머니 마지막 물질이래.\\새벽 6시에 간대. 와.}
\chatbubble{다원}{나 그날 심부름 안 해. 놀 거야.}
\end{chatthread}

\dobira{9장}{마지막 물질, 처음 듣는 휘슬}{images/pagefill/ch17.png}{images/pagefill/ch17-page.png}

\section{마지막 물질의
아침}\label{uxb9c8uxc9c0uxb9c9-uxbb3cuxc9c8uxc758-uxc544uxce68}

\dateline{7월 31일 일요일 · 생중계 후 15일}

새벽 5시 50분, 탈의장 앞
가로등은 아직 켜져 있었다.
채아와 다원이 걸어왔다. 다원의
주머니가 납작한 돌로 불룩했다.

``나 오늘 심부름 안 해. 엄청
열심히 놀 거야.''

할머니는 탈의장 안쪽에서 테왁
끈을 고치고 있었다.

그때 채아가 단우 옆구리를 쿡
찔렀다. 탈의장 문 앞에 간호사
가방을 멘 사람이 서 있었다.
전날 저녁 근무를 마치고 온
길인 듯했다.

``너네 엄마도 와?''

``\ldots 몰랐어.''

할머니는 문 쪽을 한 번 보더니,
말없이 고무옷을 입었다. 엄마는
문지방을 넘지 않고 가방끈만
고쳐 멨다.

``다들 부지런하네.''

오삼촌이 하품을 하며 평상에
앉았다. 늘 그렇듯 어두운
휴대폰 화면이었는데,
어스름이라 오늘은 옆에서도
글자가 보였다.

\begin{Verbatim}[samepage=true,frame=single,framesep=2.5mm,rulecolor=\color{gray},formatcom=\logfont]
귀 기울임 | 내 정보
새벽세시감귤
이번 주 리더보드 1위
\end{Verbatim}

``\ldots 삼촌.''

``응?''

``오삼촌이었어요?!''

채아가 고개를 쭉 뺐다. ``그럼
제 영상에 `원본 고마워요' 단
것도요?''

오삼촌은 대답 대신 귤을 깠다.

``새벽 세 시에 깨는 버릇은
그물 놓고도 안 고쳐져.''

``그걸 다 손으로요? 전 삼촌도
봇인 줄 알았어요.''

``손으로 해야 귀에 남지.''

``그때 `축하드려요 1위님
\textsuperscript{}'도요?''

``축하한 거다. 진짜로.''
오삼촌이 귤 반쪽을 단우에게
쥐여 주었다. ``웃는 얼굴이야,
그거.''

탁 박사가 단우에게 종이컵을
내밀었다.

``넌 코코아. 연료는 어른 거야.
이름 지도 라벨러를 새로 받아.
조건은 하나. 라벨마다 붙인
사람 이름이랑 이유.''

``이유도요?''

``왜 그렇게 붙였는지. 모르면
모른다고 쓰고.''

``나는 벌써 냈다.''
오삼촌이었다.

``그거 미성년자도 돼요?''

``이름은 가리고.''

``서울 가도 할 수 있어요?''

``녹음 조각은 어디서든 열려.''
탁 박사가 단우 귀를 톡
가리켰다. ``귀만 들고 가.''

단우는 탈의장 문 쪽을 보았다.

``엄마한테 먼저 물어볼게요.''

``그래야지.''

문 앞의 엄마가 입을 열었다.

``들었어. 저녁에 얘기해.''

\Needspace{4\baselineskip}\begin{center}◇\end{center}\nopagebreak

6시 반, 해녀들이 바다로
내려갔다. 정임이 할머니를
돌아보았다.

``덕순 삼춘, 배도
탔다면서요!''

``한 번이다. 다신 안 탄다.''

할머니도 테왁을 밀며 큰여
쪽으로 나갔다. 단우는 불턱 앞
바위에 앉아 눈을 감았다.

호오이---

평평하게 가다가 올라가고, 두
번 내려가는 선.

``정임 삼춘!''

테왁 옆에서 볼이 동그란 얼굴이
손을 흔들었다.

``맞혔냐? 테레비 나오는
거냐?''

``아니요! 그냥 들은 거예요!''

``에이, 아깝다.''

방파제 끝에서 퐁, 퐁, 퐁, 돌이
튀었다. 다원이 팔을 번쩍
들었다.

``다섯 번!''

``네 번이야. 내가 찍었어.''
채아가 휴대폰을 흔들었다.

``그럼 영상 보고 정정해.''

얼마 뒤, 할머니가 망사리를
끌고 와 단우 옆에 앉아 끈부터
감았다.

``큰여 바깥에 반달이
지나가더라. 말젯이 뵌다.''

단우는 눈을 번쩍 떴다. 큰여
바깥에 회색 등 두 개가 나란히
떠올랐다가 가라앉았다.

``느, 아까부터 눈 감고 뭐
하냐.''

``듣는 거요.''

``녹음은?''

``오늘은 안 해요. 그냥
들을래요.''

단우는 다시 눈을 감았다.

호오이---

이번 소리는 가까웠다. 머릿속에
선이 그려졌다. 평평, 올라감,
내려감, 다시---

``미숙 올라왔다.''

할머니였다. 끈을 감는 손은
멈추지 않았고, 고개도 바다
쪽으로 돌리지 않았다.

단우는 눈을 떴다. 바위 앞에서
미숙이 이쪽을 보고 있었다.
물안경 너머 눈이 동그래지더니,
입이 활짝 벌어졌다.

``네, 삼춘!''

방파제의 다원이 돌을 던지다
말고 돌아보았다.

``삼춘! 방금 우리 엄마
불렀어요?'' 다원이 외쳤다.

``미숙 올라왔다 했다. 왜.''

``\ldots 아니요!''

다원이 돌을 휙 던졌다. 퐁, 퐁,
퐁, 퐁, 퐁, 퐁.

``여섯 번!''

``단우!'' 미숙이 바위로
올라오며 단우 팔을 툭 쳤다.
``그때 `오늘 끝나면 꼭요!'
했지? 그 오늘, 한참 지났어.''

``내 소리, 나는 한 번도 못
들었어.''

``그날 이 바위 앞에서 받은
거예요. 잘 받아 두라고 하셨던
거.''

단우가 헤드폰을 미숙의 귀에
씌우고 재생을 눌렀다. 미숙은
두 손으로 헤드폰을 누른 채
눈을 감았다가, 떴다가, 다시
감았다.

``이게\ldots{} 나야?''

``네. 평평하다가 올라가고,
내려갔다가 다시 평평.
\mbox{F-U-D-F요.}''

미숙이 손가락으로 공중에
옆으로, 위로, 아래로, 다시
옆으로 선을 그었다.

``끝이 평평하네. 삼춘 거는
끝에 내려가는데.''

할머니가 끈을 다 감고
일어섰다.

``한 번 더 든다.''

미숙이 헤드폰을 꼭 쥔 채
단우를 보았다.

``나도. 한 번 더 들을래.''

\Needspace{4\baselineskip}\begin{center}◇\end{center}\nopagebreak

할머니의 테왁이 큰여
바깥쪽으로 점점 작아졌다.
미숙이 헤드폰을 벗어 단우
무릎에 올려놓았다.

``이게 삼춘 마지막이야.''

단우는 대답 대신 눈을 감았다.
파도 소리, 갈매기 소리.

숨비소리는 한동안 들리지
않았다. 단우는 헤드폰 줄을 꽉
쥐었다.

아주 먼 데서, 바람을 타고 소리
하나가 왔다.

호오이---

올라가서 오래 평평하게 있다가
내려오는 소리,
\mbox{U-F-F-D였다.}

``할머니.''

``삼춘.''

단우와 미숙은 한꺼번에 말하고,
서로 쳐다보다가 같이 웃었다.
웃는데, 단우는 목 안쪽이
뜨거워졌다.

테왁이 바위 쪽으로 다가왔다.
할머니가 바위를 짚고
올라오는데, 그 앞에 마른
수건을 든 엄마가 먼저 와
있었다.

엄마는 아무 말 없이 수건을
내밀었다. 할머니는 수건을 한
번, 엄마 얼굴을 한 번 보더니
받아 들었다.

엄마는 비켜서지 않고 할머니
옆에 섰다. 할머니가 바다 쪽을
본 채 수건으로 얼굴을 닦았다.

``\ldots 밥은?''

엄마가 입을 열었다가 닫았다.
다시 열었다.

``\ldots 아직.''

할머니가 고개를 끄덕이고
탈의장 쪽으로 걸었다. 엄마가
반걸음 뒤에서 따라갔다. 단우는
벌떡 일어났다.

``할머니! 불턱에서 녹음한
숨비소리요. 계속 갖고 있어도
돼요?''

``뭐 하젠?''

단우는 잠깐 머뭇거렸다.
``\ldots 아직 몰라요.''

``해라.''

할머니 말이 이어졌다.

``대신 이건 내 거다. 쓸 때마다
물어봐라.''

``\ldots 네. 쓸 때마다요.''

김복자 회장이 소리쳤다.

``덕순 삼춘! 은퇴 잔치는
고기여, 회여?''

``물때 봐야 알주.''

해녀들이 와르르 웃었다. 다원과
채아가 동시에 외쳤다.

``회요!''

``고기요!''

\Needspace{4\baselineskip}\begin{center}◇\end{center}\nopagebreak

그날 밤 9시 반, 단우는 책상
앞에 앉아 있었다. 9시만 넘으면
감기던 눈꺼풀이 오늘은
멀쩡했다. 방문 너머에서 엄마
목소리가 들렸다.

``연구소에 낼 땐 이름은
가리고. 그건 꼭이야.''

``알았어. 밥 먹을 때도
말했잖아.''

노트북에는 이름 사전이 열려
있었다. \#A-3318 옆에 휘슬
번호 W-07과 바당이 문장
`돌려줘'가 적혀 있었고, 메모
칸은 물음표였다. 단우의
손가락이 그 줄에서 멈췄다.

``\ldots 이건 이제 알지.''

단우는 물음표를 지우고 고쳐
썼다.

\begin{Verbatim}[samepage=true,frame=single,framesep=2.5mm,rulecolor=\color{gray},formatcom=\logfont]
메모        오역. 맞아 보인 건 우연.
\end{Verbatim}

옆 창에서는 listen.py가 돌고
있었다. 저녁에 오삼촌이 방파제
끝에 매달아 준 휴대용
마이크였다. 작은 송신기가 그
소리를 집까지 보내 주었다.
휴대폰에는 계장의 마을방
공지가 떠 있었다.

\begin{chatthread}
\chatroomtitle{곰새코지 마을방}
\chatbubble{강덕수}{오늘 밤 항구 측량선이 시험 운항합니다.\\밤 9시 30분부터 10시 30분까지 방파제 바깥을 지납니다.}
\end{chatthread}

창밖, 방파제 너머로 작은 불빛
하나가 천천히 지나갔다. 9시
41분, listen.py 창에 줄이
내려오기 시작했다.

\begin{Verbatim}[samepage=true,frame=single,framesep=2.5mm,rulecolor=\color{gray},formatcom=\logfont]
listen.py | 방파제 휴대 마이크
21:41:07 미분류-모름 U-U-D
21:41:09 미분류-모름 D-F-U
21:41:12 미분류-모름 U-U-D
21:41:14 미분류-모름 D-F-U
21:41:15 미분류-모름 F-D-D-U
\end{Verbatim}

헤드폰 속에서 가는 소리 둘이
번갈아 울리다가, 끝에 다른
소리 하나가 짧게 끼어들었다.

``\ldots 처음 보는 모양인데.''

불빛은 아직 지나가는 중이었다.

``불빛이 지나가는
동안에만\ldots?''

W-07은 없었다. 오삼촌이 새끼
소리 아니냐던 \mbox{U-F도}
아니었다. 그런데 끝의 한 줄을
보자, 생중계 날 밤의 `+1'이
떠올랐다.

단우는 기록을 그날 밤까지
거슬러 올라갔다. `+1'을 눌러
펼치자 모양이 나왔다.

\begin{Verbatim}[samepage=true,frame=single,framesep=2.5mm,rulecolor=\color{gray},formatcom=\logfont]
listen.py | 등대 마이크
23:58:30 W-07 U-D-U-F
23:59:02 미분류-모름 F-D-D-U
\end{Verbatim}

오늘 밤 마지막 줄과 같은
모양이었다.

단우는 이름 사전 맨 아래에 새
줄 세 개를 만들었다.

번호는 이름이 아니었다.
아직은, 서로 다른 모양을 담아
둔 빈 상자 이름일 뿐이었다.

\begin{Verbatim}[samepage=true,frame=single,framesep=2.5mm,rulecolor=\color{gray},formatcom=\logfont]
이름 사전 | 새 줄
휘슬 번호   W-31, W-32, W-33
연구소 번호 없음
할머니 별명 없음
붙인 사람   고단우
이유        모름
메모        F-D-D-U는 생중계 날 밤에도
오늘 밤     측량선 시험 운항
\end{Verbatim}

저장 창이 떴다. 손가락이 먼저
움직였다.

\dateline{이름사전\_최종\_진짜최종}

``\ldots 버그.''

그 글자를 지우고,
\texttt{이름사전} 그대로
저장했다.

단우는 헤드폰을 다시 썼다.

\addtocontents{toc}{\protect\tocgroup{책 뒤에}}

\bookmarksetup{startatroot}

\chapter*{진짜 과학
노트}\label{uxc9c4uxc9dc-uxacfcuxd559-uxb178uxd2b8}
\addcontentsline{toc}{chapter}{진짜
과학 노트}

\markboth{진짜 과학 노트}{진짜
과학 노트}

\backmatterstyle

어디까지가 사실일까?

이 책은 지어낸 이야기다.
돌고래 소리에 관한 내용은 실제
연구에서 왔다. 그래서 사실과
상상을 `진짜인 것', `아직
모르는 것', `지어낸 것' 세
묶음으로 나눴다. 괄호의 숫자는
그 일이 나오는 장이다.

연구한 동물도 표시했다.

\begin{itemize}
\tightlist
\item
  \textbf{{[}남방큰돌고래{]}}:
  제주에 사는 종을 직접
  연구했다.
\item
  \textbf{{[}가까운 종{]}}:
  사촌뻘인 큰돌고래를
  연구했다.
\item
  \textbf{{[}다른 종{]}}: 그
  밖의 고래를 연구했다.
\end{itemize}

큰돌고래에서 밝혀진 일이
남방큰돌고래도 같다는 보장은
없다. 그래서 어느 돌고래의
일인지 꼭 구분해서 읽자.

\section*{진짜인
것}\label{uxc9c4uxc9dcuxc778-uxac83}
\addcontentsline{toc}{section}{진짜인
것}

\markright{진짜인 것}

\begin{enumerate}
\def\labelenumi{\arabic{enumi}.}
\tightlist
\item
  {[}남방큰돌고래{]} 휘슬과
  클릭은 다른 소리다. 클릭은
  평균 약 96킬로헤르츠(kHz)라
  사람 귀에 거의 안 들린다.
  우리가 듣는 `딸깍'은 일부다.
  (2장)
\item
  {[}남방큰돌고래{]} 이름 같은
  `시그니처 휘슬'이
  남아공·일본·호주에서
  확인됐다. 2026년에는 제주
  무리에서도 보고됐다. (4장)
\item
  {[}남방큰돌고래{]} 호주에서
  새끼와 떨어진 어미는
  시그니처 휘슬을 많이 냈다.
  마트에서 엄마를 놓치고
  ``엄마!''만 부르는 것과
  비슷하다. (4장)
\item
  {[}가까운 종{]} 큰돌고래
  어미는 새끼가 곁에 있으면
  휘슬 꼭대기가 높아진다.
  (4장)
\item
  {[}가까운 종{]} 20년이
  지나도 헤어진 동료의 휘슬을
  알아본 큰돌고래가 있다.
  (3장)
\item
  {[}가까운 종{]} 천천히 가는
  작은 배 하나만 있어도 서로
  들리는 범위가 약 4분의 1
  줄었다. (7장)
\item
  {[}남방큰돌고래{]} 제주
  연구팀이 휘슬을 골라내는
  AI를 만들었다. 고른 100개 중
  96개쯤이 진짜 휘슬이었다.
  (2장)
\item
  {[}남방큰돌고래{]} 제주
  바다에 약
  110\textasciitilde120마리가
  산다. 2012년에
  해양보호생물이 됐다. (3장)
\item
  {[}다른 종: 범고래{]} 2024년
  캐나다에서 갇힌 새끼에게
  범고래 소리를 틀었더니
  오히려 쫓아냈다. 한 번의
  사례다. (7장)
\end{enumerate}

\section*{아직 모르는
것}\label{uxc544uxc9c1-uxbaa8uxb974uxb294-uxac83}
\addcontentsline{toc}{section}{아직
모르는 것}

\markright{아직 모르는 것}

책 속 인물도, 실제 과학자도
모르는 것이다.

\begin{enumerate}
\def\labelenumi{\arabic{enumi}.}
\setcounter{enumi}{9}
\tightlist
\item
  휘슬이 사람 이름과 똑같은
  일을 하는지는 모른다. 탁
  박사도 ``그건 몰라''라고
  했다.
\item
  {[}남방큰돌고래{]}가 친구의
  휘슬을 흉내 내 부르는지,
  새끼 곁에서 휘슬이
  달라지는지는 확인되지
  않았다. 기억이 얼마나
  가는지도 모른다.
\item
  돌고래 소리를 사람 문장으로
  옮기는 기술은 없다. 향유고래
  소리도 뜻은 거의 모른다.
  (1장)
\item
  W-07이 누구의 휘슬인지는
  이야기 속에서도 끝까지
  모른다. 단우가 아는 것은
  숫자가 바뀐 날짜뿐이다.
\end{enumerate}

\section*{지어낸
것}\label{uxc9c0uxc5b4uxb0b8-uxac83}
\addcontentsline{toc}{section}{지어낸
것}

\markright{지어낸 것}

이야기를 위해 작가가 만든
것이다. 진짜 돌고래나 사람과
관계없다.

\begin{enumerate}
\def\labelenumi{\arabic{enumi}.}
\setcounter{enumi}{13}
\tightlist
\item
  반달, W-07, 불쑥이, 말젯이,
  곰새코지, 그물먹는여,
  소등굴은 모두 지어낸 이름과
  곳이다.
\item
  물마루 AI 연구소, 바당이,
  `귀 기울임' 앱, 4장의 W-07
  그래프, 푸른코드 대회는
  작가가 만들었다.
\item
  제주 해양동물 구조팀과
  `도와주는 배' 허락은 쉽게
  바꾼 상상이다. 숨비소리로
  사람을 알아보는 할머니의
  귀도 지어낸 능력이다.
\end{enumerate}

\section*{제주에서 돌고래를 볼
때 (2026년
기준)}\label{uxc81cuxc8fcuxc5d0uxc11c-uxb3ccuxace0uxb798uxb97c-uxbcfc-uxb54c-2026uxb144-uxae30uxc900}
\addcontentsline{toc}{section}{제주에서
돌고래를 볼 때 (2026년 기준)}

\markright{제주에서 돌고래를
볼 때 (2026년 기준)}

50미터 안으로 가지 않는다.
300미터 안에서는 프로펠러를
멈추고, 배 3척 이상이 함께
들어가지 않는다. 만지거나
먹이를 주지 않는다. 어기면
과태료를 200만 원 이하로 낸다.
규정은 바뀔 수 있으니 가기
전에 다시 확인하자. 엔진까지
끄라는 규칙은 2026년에는 없다.
이야기 속 2033년에는 같다고
가정했다.

근거와 참고 문헌 전체는 교사용
자료에 있어요.

책에 나온 활동 12개와 토론
질문은 교사용 자료(무료)에
있어요. 받는 곳: {[}{[}교사용
자료 주소{]}{]}

\bookmarksetup{startatroot}

\chapter*{단우의
메모}\label{uxb2e8uxc6b0uxc758-uxba54uxbaa8}
\addcontentsline{toc}{chapter}{단우의
메모}

\markboth{단우의 메모}{단우의
메모}

\backmatterstyle

낱말 모음

책에 나온 낯선 낱말 가운데 꼭
필요한 것만 가나다순으로
모았다. 괄호 안은 그 낱말이
처음 나온 장이다. 전체 낱말은
교사용 자료에 있다.

\textbf{곰새기} (3장)\\
제주 해녀들이 돌고래를 부르던
이름 가운데 하나다.

\textbf{라벨} (1장)\\
소리 조각에 붙이는 이름표다.
사진에 다는 해시태그와
비슷하다.

\textbf{말젯이} (1장)\\
제주어로 `셋째'라는 뜻이다.
`막내'가 아니다.

\textbf{물때} (1장)\\
바닷물이 들고 나는 때다. 욕실
타일 틈의 거뭇한 때와는 다른
말이다.

\textbf{미분류-모름} (6장)\\
아는 모양과 맞지 않는 소리에
붙이는, `모른다'는 표시다.
지우지 않고 모른다고 남겨 두는
칸이다.

\textbf{바당이의 네 칸}
(2장)\\
바당이가 소리에서 문장을
지어내는 네 단계다. 듣기,
나누기, 짝 맞추기, 문장 만들기
순서다.

\textbf{불턱} (3장)\\
옛날 해녀들이 옷을 갈아입고
불을 쬐던 돌담 자리다. 지금은
대부분 탈의장으로 바뀌었다고
한다.

\textbf{삼춘} (1장)\\
제주에서 동네 어른을 부르는
말이다. `덕순 삼춘'처럼 이름
뒤에 붙인다.

\textbf{숨비소리} (3장)\\
해녀가 물 위로 솟으며 참았던
숨을 내쉬는 소리다. 이 책의
할머니는 숨비소리만 듣고도
누구인지 알아본다.

\textbf{시그니처 휘슬} (4장)\\
돌고래가 자기만 내는 휘슬이다.
이름 같은 것이다.

\textbf{정정문} (6장)\\
틀리게 알린 것을 바로잡는
글이다. 단우는 첫 줄에
`정정합니다.'라고 썼다.

\textbf{테왁} (1장)\\
해녀가 물에 띄워 놓고 몸을
기대는 둥근 부표다. 물 위에서
해녀가 있는 자리를 알려 주기도
한다.

\textbf{확증편향} (1장)\\
보고 싶은 것만 골라 보는
마음의 버릇이다. 응원하는 팀이
지면 심판이 상대편 같아 보이는
것과 같다.

\textbf{휘슬} (2장)\\
돌고래끼리 주고받는 높은
휘파람 소리다. 대부분 사람
귀에도 들린다.

\bookmarksetup{startatroot}

\chapter*{작가의
말}\label{uxc791uxac00uxc758-uxb9d0}
\addcontentsline{toc}{chapter}{작가의
말}

\markboth{작가의 말}{작가의
말}

\backmatterstyle

이 이야기는 질문 하나에서
시작했어요.

``AI가 돌고래 말을 번역했대!''

이런 뉴스가 나오면 우리는
믿을까요? 믿는다면, 그 문장은
누가 만든 걸까요? 단우는 그
질문 한가운데에 선 아이예요.

이야기를 다 읽은 여러분도
궁금할 거예요. 어디까지가
진짜이고, 어디부터 지어낸
걸까요? 그래서 여기에 나눠
적어 둡니다.

\section*{진짜인
것}\label{uxc9c4uxc9dcuxc778-uxac83-1}
\addcontentsline{toc}{section}{진짜인
것}

\markright{진짜인 것}

돌고래 소리를 듣는 연구는
실제로 있어요. 많은 돌고래가
저마다 자기만의 휘파람 소리를
가지고 있어요. 과학자들은
이것을 `시그니처 휘슬'이라고
불러요.

2026년에는 제주
남방큰돌고래에게서도 이 휘슬이
보고됐어요. 그 연구는 휘슬
모양을 `내려감, 올라감,
평평'으로 쪼개 비교했어요.
단우의 \mbox{U-D-F} 표기와
닮았지요.

돌고래를 생중계로 보여 주는
일도 실제로 있어요. 우리나라
고래 연구 기관도 그런 방송을
한 적이 있어요.

제주 남방큰돌고래는
110\textasciitilde120마리쯤으로
추정돼요. 2012년에
해양보호생물이 되었어요.

배가 돌고래를 볼 때 지킬
규칙도 있어요. 50미터 안으로
다가가면 안 돼요. 300미터
안에서는 프로펠러를 멈춰야
해요. 배 세 척 이상이 한꺼번에
300미터 안에 들어가도 안
되고요. 돌고래를 만지거나
먹이를 주는 것도 안 돼요. 이
책에는 이 규칙을 2026년 기준
그대로 썼어요. 이야기 속
2033년의 규칙이 어떨지는
아무도 몰라요.

제주 바다에서는 그물이나
낚싯줄에 걸린 돌고래가 여러 번
발견됐어요. 스스로 풀려난
돌고래도 있고, 끝내 살지 못한
돌고래도 있어요. 말젯이는 그
누구도 아니에요. 말젯이가 걸린
모습도, 구조된 날도, 구조
방법도 모두 지어냈어요.

\section*{지어낸
것}\label{uxc9c0uxc5b4uxb0b8-uxac83-1}
\addcontentsline{toc}{section}{지어낸
것}

\markright{지어낸 것}

곰새코지 마을, 소등굴,
그물먹는여는 이 이야기를 위해
지어낸 곳이에요. 물마루 AI
연구소와 AI 바당이, 앱 `귀
기울임'도 지어냈어요. 푸른코드
청소년 AI 대회도 없는
대회예요.

등장인물도 모두 지어낸
사람이에요. 해녀 미숙도 실제
어떤 분을 그린 게 아니에요.

말젯이를 구한 구조팀도
상상이에요. 제주에는 실제로
돌고래를 도우려 애쓰는
사람들이 있어요.
제주돌고래긴급구조단, 제주도
전담팀, 해양동물 구조치료기관
같은 곳이에요. 작중 구조팀은
그중 어느 곳도 아니에요.

그리고 가장 큰 상상이 하나
있어요. 돌고래 휘슬을 사람
문장으로 옮기는 기술이에요.
휘슬 하나가 무슨 뜻인지 아는
사람은 아직 없어요.

\section*{해녀
이야기}\label{uxd574uxb140-uxc774uxc57cuxae30}
\addcontentsline{toc}{section}{해녀
이야기}

\markright{해녀 이야기}

제주 해녀들은 돌고래를
`곰새기'라고도 불렀대요. 이
이름은 진관훈 선생님이
제이누리에 쓴 글(2024)에서
배웠어요.

숨비소리는 해녀가 물 위로
올라와 참았던 숨을 내쉬는
소리예요. `호오이' 하고 길게
울리지요. 하지만 숨비소리만
듣고 누군지 알아보는 건 이 책
속 할머니(그리고 할머니 소리를
알아듣는 미숙)의 능력이에요.
모든 해녀가 그렇다는 기록은
찾지 못했어요.

곰새코지 해녀들은 아직 돌담
불턱에 모여요. 하지만 요즘
해녀들은 대부분 탈의장을
쓴대요.

해녀 이야기는 국가기록원의
해녀 자료와 제이누리의 해녀
글을 읽으며 썼어요. 제주 말은
국립국어원 우리말샘에서
찾아보았어요.

\section*{모름
칸}\label{uxbaa8uxb984-uxce78}
\addcontentsline{toc}{section}{모름
칸}

\markright{모름 칸}

이 책을 쓰는 동안 저도 `모름'
칸을 여러 번 만들었어요.
찾아도 나오지 않는 것이
많았거든요. 그 칸들은 비워
두지 않고 `모름'이라고
적었어요. 그래야 다음에 무엇을
찾을지 보였어요.

제주 말과 해녀 이야기를 살펴
준 {[}{[}감수자 이름:
제주어·해녀{]}{]} 님, 돌고래
과학을 살펴 준 {[}{[}감수자
이름: 돌고래 과학{]}{]} 님께
고맙습니다.

{[}{[}날짜: 예) 2026년
가을{]}{]}\\
{[}{[}작가 이름{]}{]}

\colophonstart

\textbf{곰새코지 돌고래 목소리
사건}

{[}{[}초판 1쇄 발행일{]}{]}

지은이 이희근\\
그린이 {[}{[}삽화가
이름{]}{]}\\
펴낸곳 에르고스피어\\
{[}{[}출판사 신고번호 · 주소 ·
전자우편{]}{]}

ISBN {[}{[}종이책 ISBN ---
발급 전{]}{]}\\
값 \mbox{14,000원}

{[}{[}AI 활용 고지 문구 ---
결정 전{]}{]}

이 책에 나오는 사람, 기관,
마을, 바다 이름은 지어낸
것입니다.

\colophonend


\backmatter


\end{document}
